% ============================================================
% SHA-256 as a mathematical object: a field guide and a
% research program.  (Baseline variant.)
%
% Build:  make            (all variants)
%         latexmk -pdf main.tex   (this variant only)
%
% Structure: preamble.tex holds all packages/macros; each section
% lives in sections/NN-name.tex; backmatter.tex holds the appendix
% and the clustered references. Sibling variants main-algebra.tex
% and main-projects.tex share everything but the front matter.
% ============================================================
\providecommand{\DocPdfTitle}{SHA-256 as a Mathematical Object}
% ============================================================
% Shared preamble for all variants of the document.
% Each top-level file (main.tex, main-algebra.tex,
% main-projects.tex) does:
%     \providecommand{\DocPdfTitle}{...}   % optional override
%     % ============================================================
% Shared preamble for all variants of the document.
% Each top-level file (main.tex, main-algebra.tex,
% main-projects.tex) does:
%     \providecommand{\DocPdfTitle}{...}   % optional override
%     % ============================================================
% Shared preamble for all variants of the document.
% Each top-level file (main.tex, main-algebra.tex,
% main-projects.tex) does:
%     \providecommand{\DocPdfTitle}{...}   % optional override
%     \input{preamble}
% then supplies its own title/abstract and \input's the sections.
% ============================================================
\documentclass[11pt]{article}

\usepackage[T1]{fontenc}
\usepackage[utf8]{inputenc}
\usepackage[margin=1.15in]{geometry}

% ---- Math (load before newpx; newpxmath supplies AMS symbols) ----
\usepackage{amsmath,amsthm}

% ---- Typography ----
\usepackage{newpxtext,newpxmath}   % Palatino-family text + matching math
\usepackage{microtype}

% ---- Tables ----
\usepackage{booktabs,array}
\usepackage{float}   % [H] to pin the batteries table at its reference

% ---- Color, graphics, boxes ----
\usepackage[dvipsnames]{xcolor}
\usepackage{graphicx}
\usepackage{tikz}
\usetikzlibrary{arrows.meta,positioning,calc,shapes.geometric,decorations.pathreplacing}
\usepackage[most]{tcolorbox}

\definecolor{linknavy}{RGB}{0,60,120}
\definecolor{citewine}{RGB}{120,20,50}
\definecolor{projfill}{RGB}{245,247,240}
\definecolor{projframe}{RGB}{90,120,70}
\definecolor{msgfill}{RGB}{213,229,245}   % message bits (blue)
\definecolor{lenfill}{RGB}{250,224,188}   % length field (orange)
\definecolor{padfill}{RGB}{198,224,180}   % single pad bit (green)
\definecolor{archgreen}{RGB}{20,110,60}   % local-archive links in bibliography
\definecolor{codegray}{RGB}{120,120,120}

% ---- Code listings (appendix) ----
\usepackage{listings}
\lstset{
  language=Python,
  basicstyle=\ttfamily\footnotesize,
  keywordstyle=\bfseries\color{linknavy},
  commentstyle=\itshape\color{codegray},
  stringstyle=\color{archgreen},
  showstringspaces=false,
  numbers=left,
  numberstyle=\tiny\color{codegray},
  numbersep=8pt,
  breaklines=true,
  keepspaces=true,
  columns=flexible,
  frame=leftline,
  framerule=0.8pt,
  rulecolor=\color{projframe},
  xleftmargin=1.5em,
  aboveskip=1em,
}

% ---- Citations: natbib (author-year) + BibTeX ----
\usepackage[round]{natbib}

% ---- Hyperlinks with SECTION-level back-references ----
% backref=section makes each bibliography entry list the *sections*
% in which it is cited, hyperlinked back into the text.
\providecommand{\DocPdfTitle}{SHA-256 as a Mathematical Object}
\usepackage[colorlinks=true,
            linkcolor=linknavy,
            citecolor=citewine,
            urlcolor=linknavy,
            backref=section,
            pdftitle={\DocPdfTitle},
            pdfauthor={}]{hyperref}
\usepackage{cleveref}

% Last-updated stamp: `make` regenerates lastupdated.tex with the
% current Pacific date/time; fall back to \today if absent.
\InputIfFileExists{lastupdated}{}{}
\providecommand{\lastupdated}{\today}

% The bibliography is thematically clustered: tools/clusterbbl.py
% (run automatically after bibtex via .latexmkrc) inserts cluster
% headers into the .bbl. Suppress natbib's automatic heading; we
% typeset our own.
\renewcommand{\bibsection}{}

% Format back-references as "Cited in §2, §4.1, and §6."
\renewcommand*{\backrefsep}{, \S}
\renewcommand*{\backreftwosep}{ and \S}
\renewcommand*{\backreflastsep}{, and \S}
\renewcommand*{\backref}[1]{}% disable the plain list; use \backrefalt
\renewcommand*{\backrefalt}[4]{%
  \ifcase #1\relax
  \or {\footnotesize\itshape Cited in \S#2.}%
  \else {\footnotesize\itshape Cited in \S#2.}%
  \fi}

% ---- Theorem environments ----
\newtheorem{theorem}{Theorem}[section]
\newtheorem{fact}[theorem]{Fact}
\newtheorem{conjecture}[theorem]{Conjecture}
\theoremstyle{definition}
\newtheorem{definition}[theorem]{Definition}
\theoremstyle{remark}
\newtheorem{remark}[theorem]{Remark}

% ---- Undergraduate project boxes ----
% \begin{project}{Short title}{prerequisites}{time window} ... \end{project}
% Plain (non-\refstepcounter) counter: keeps citations inside boxes
% anchored to the enclosing section, so backref dedup works.
\newcounter{projectctr}
\newtcolorbox{project}[3]{%
  enhanced, breakable,
  code={\stepcounter{projectctr}},
  colback=projfill, colframe=projframe,
  fonttitle=\bfseries\small,
  title={Project \theprojectctr: #1},
  attach boxed title to top left={yshift=-2mm, xshift=4mm},
  boxed title style={colback=projframe},
  after upper={\par\smallskip\footnotesize\textit{Prerequisites:} #2\quad
               \textit{Window:} #3},
  left=2mm, right=2mm, top=3mm, bottom=2mm
}

% ---- Reading-spine box for variant front matter ----
% \begin{readingspine}{title} ... \end{readingspine}
\newtcolorbox{readingspine}[1]{%
  enhanced, breakable,
  colback=projfill, colframe=projframe,
  fonttitle=\bfseries\small,
  title={#1},
  attach boxed title to top left={yshift=-2mm, xshift=4mm},
  boxed title style={colback=projframe},
  left=2mm, right=2mm, top=3mm, bottom=2mm
}

% Local archived copy of a cited work: visually distinct (green) and
% clickable — opens refs/<file> relative to this PDF's location.
\newcommand{\localpdf}[1]{\href{./refs/#1}{\textcolor{archgreen}{$\blacktriangleright$\,\texttt{refs/#1}}}}

% ---- Notation macros ----
\newcommand{\Ztwo}{\mathbb{Z}/2^{32}\mathbb{Z}}
\newcommand{\Ftwo}{\mathbb{F}_2}
\newcommand{\State}{\mathcal{S}}
\newcommand{\shaxor}{\oplus}
\newcommand{\shaplus}{\boxplus}
\newcommand{\SHA}{\textsf{SHA-256}}
\newcommand{\bitstrings}{\{0,1\}}

% then supplies its own title/abstract and \input's the sections.
% ============================================================
\documentclass[11pt]{article}

\usepackage[T1]{fontenc}
\usepackage[utf8]{inputenc}
\usepackage[margin=1.15in]{geometry}

% ---- Math (load before newpx; newpxmath supplies AMS symbols) ----
\usepackage{amsmath,amsthm}

% ---- Typography ----
\usepackage{newpxtext,newpxmath}   % Palatino-family text + matching math
\usepackage{microtype}

% ---- Tables ----
\usepackage{booktabs,array}
\usepackage{float}   % [H] to pin the batteries table at its reference

% ---- Color, graphics, boxes ----
\usepackage[dvipsnames]{xcolor}
\usepackage{graphicx}
\usepackage{tikz}
\usetikzlibrary{arrows.meta,positioning,calc,shapes.geometric,decorations.pathreplacing}
\usepackage[most]{tcolorbox}

\definecolor{linknavy}{RGB}{0,60,120}
\definecolor{citewine}{RGB}{120,20,50}
\definecolor{projfill}{RGB}{245,247,240}
\definecolor{projframe}{RGB}{90,120,70}
\definecolor{msgfill}{RGB}{213,229,245}   % message bits (blue)
\definecolor{lenfill}{RGB}{250,224,188}   % length field (orange)
\definecolor{padfill}{RGB}{198,224,180}   % single pad bit (green)
\definecolor{archgreen}{RGB}{20,110,60}   % local-archive links in bibliography
\definecolor{codegray}{RGB}{120,120,120}

% ---- Code listings (appendix) ----
\usepackage{listings}
\lstset{
  language=Python,
  basicstyle=\ttfamily\footnotesize,
  keywordstyle=\bfseries\color{linknavy},
  commentstyle=\itshape\color{codegray},
  stringstyle=\color{archgreen},
  showstringspaces=false,
  numbers=left,
  numberstyle=\tiny\color{codegray},
  numbersep=8pt,
  breaklines=true,
  keepspaces=true,
  columns=flexible,
  frame=leftline,
  framerule=0.8pt,
  rulecolor=\color{projframe},
  xleftmargin=1.5em,
  aboveskip=1em,
}

% ---- Citations: natbib (author-year) + BibTeX ----
\usepackage[round]{natbib}

% ---- Hyperlinks with SECTION-level back-references ----
% backref=section makes each bibliography entry list the *sections*
% in which it is cited, hyperlinked back into the text.
\providecommand{\DocPdfTitle}{SHA-256 as a Mathematical Object}
\usepackage[colorlinks=true,
            linkcolor=linknavy,
            citecolor=citewine,
            urlcolor=linknavy,
            backref=section,
            pdftitle={\DocPdfTitle},
            pdfauthor={}]{hyperref}
\usepackage{cleveref}

% Last-updated stamp: `make` regenerates lastupdated.tex with the
% current Pacific date/time; fall back to \today if absent.
\InputIfFileExists{lastupdated}{}{}
\providecommand{\lastupdated}{\today}

% The bibliography is thematically clustered: tools/clusterbbl.py
% (run automatically after bibtex via .latexmkrc) inserts cluster
% headers into the .bbl. Suppress natbib's automatic heading; we
% typeset our own.
\renewcommand{\bibsection}{}

% Format back-references as "Cited in §2, §4.1, and §6."
\renewcommand*{\backrefsep}{, \S}
\renewcommand*{\backreftwosep}{ and \S}
\renewcommand*{\backreflastsep}{, and \S}
\renewcommand*{\backref}[1]{}% disable the plain list; use \backrefalt
\renewcommand*{\backrefalt}[4]{%
  \ifcase #1\relax
  \or {\footnotesize\itshape Cited in \S#2.}%
  \else {\footnotesize\itshape Cited in \S#2.}%
  \fi}

% ---- Theorem environments ----
\newtheorem{theorem}{Theorem}[section]
\newtheorem{fact}[theorem]{Fact}
\newtheorem{conjecture}[theorem]{Conjecture}
\theoremstyle{definition}
\newtheorem{definition}[theorem]{Definition}
\theoremstyle{remark}
\newtheorem{remark}[theorem]{Remark}

% ---- Undergraduate project boxes ----
% \begin{project}{Short title}{prerequisites}{time window} ... \end{project}
% Plain (non-\refstepcounter) counter: keeps citations inside boxes
% anchored to the enclosing section, so backref dedup works.
\newcounter{projectctr}
\newtcolorbox{project}[3]{%
  enhanced, breakable,
  code={\stepcounter{projectctr}},
  colback=projfill, colframe=projframe,
  fonttitle=\bfseries\small,
  title={Project \theprojectctr: #1},
  attach boxed title to top left={yshift=-2mm, xshift=4mm},
  boxed title style={colback=projframe},
  after upper={\par\smallskip\footnotesize\textit{Prerequisites:} #2\quad
               \textit{Window:} #3},
  left=2mm, right=2mm, top=3mm, bottom=2mm
}

% ---- Reading-spine box for variant front matter ----
% \begin{readingspine}{title} ... \end{readingspine}
\newtcolorbox{readingspine}[1]{%
  enhanced, breakable,
  colback=projfill, colframe=projframe,
  fonttitle=\bfseries\small,
  title={#1},
  attach boxed title to top left={yshift=-2mm, xshift=4mm},
  boxed title style={colback=projframe},
  left=2mm, right=2mm, top=3mm, bottom=2mm
}

% Local archived copy of a cited work: visually distinct (green) and
% clickable — opens refs/<file> relative to this PDF's location.
\newcommand{\localpdf}[1]{\href{./refs/#1}{\textcolor{archgreen}{$\blacktriangleright$\,\texttt{refs/#1}}}}

% ---- Notation macros ----
\newcommand{\Ztwo}{\mathbb{Z}/2^{32}\mathbb{Z}}
\newcommand{\Ftwo}{\mathbb{F}_2}
\newcommand{\State}{\mathcal{S}}
\newcommand{\shaxor}{\oplus}
\newcommand{\shaplus}{\boxplus}
\newcommand{\SHA}{\textsf{SHA-256}}
\newcommand{\bitstrings}{\{0,1\}}

% then supplies its own title/abstract and \input's the sections.
% ============================================================
\documentclass[11pt]{article}

\usepackage[T1]{fontenc}
\usepackage[utf8]{inputenc}
\usepackage[margin=1.15in]{geometry}

% ---- Math (load before newpx; newpxmath supplies AMS symbols) ----
\usepackage{amsmath,amsthm}

% ---- Typography ----
\usepackage{newpxtext,newpxmath}   % Palatino-family text + matching math
\usepackage{microtype}

% ---- Tables ----
\usepackage{booktabs,array}
\usepackage{float}   % [H] to pin the batteries table at its reference

% ---- Color, graphics, boxes ----
\usepackage[dvipsnames]{xcolor}
\usepackage{graphicx}
\usepackage{tikz}
\usetikzlibrary{arrows.meta,positioning,calc,shapes.geometric,decorations.pathreplacing}
\usepackage[most]{tcolorbox}

\definecolor{linknavy}{RGB}{0,60,120}
\definecolor{citewine}{RGB}{120,20,50}
\definecolor{projfill}{RGB}{245,247,240}
\definecolor{projframe}{RGB}{90,120,70}
\definecolor{msgfill}{RGB}{213,229,245}   % message bits (blue)
\definecolor{lenfill}{RGB}{250,224,188}   % length field (orange)
\definecolor{padfill}{RGB}{198,224,180}   % single pad bit (green)
\definecolor{archgreen}{RGB}{20,110,60}   % local-archive links in bibliography
\definecolor{codegray}{RGB}{120,120,120}

% ---- Code listings (appendix) ----
\usepackage{listings}
\lstset{
  language=Python,
  basicstyle=\ttfamily\footnotesize,
  keywordstyle=\bfseries\color{linknavy},
  commentstyle=\itshape\color{codegray},
  stringstyle=\color{archgreen},
  showstringspaces=false,
  numbers=left,
  numberstyle=\tiny\color{codegray},
  numbersep=8pt,
  breaklines=true,
  keepspaces=true,
  columns=flexible,
  frame=leftline,
  framerule=0.8pt,
  rulecolor=\color{projframe},
  xleftmargin=1.5em,
  aboveskip=1em,
}

% ---- Citations: natbib (author-year) + BibTeX ----
\usepackage[round]{natbib}

% ---- Hyperlinks with SECTION-level back-references ----
% backref=section makes each bibliography entry list the *sections*
% in which it is cited, hyperlinked back into the text.
\providecommand{\DocPdfTitle}{SHA-256 as a Mathematical Object}
\usepackage[colorlinks=true,
            linkcolor=linknavy,
            citecolor=citewine,
            urlcolor=linknavy,
            backref=section,
            pdftitle={\DocPdfTitle},
            pdfauthor={}]{hyperref}
\usepackage{cleveref}

% Last-updated stamp: `make` regenerates lastupdated.tex with the
% current Pacific date/time; fall back to \today if absent.
\InputIfFileExists{lastupdated}{}{}
\providecommand{\lastupdated}{\today}

% The bibliography is thematically clustered: tools/clusterbbl.py
% (run automatically after bibtex via .latexmkrc) inserts cluster
% headers into the .bbl. Suppress natbib's automatic heading; we
% typeset our own.
\renewcommand{\bibsection}{}

% Format back-references as "Cited in §2, §4.1, and §6."
\renewcommand*{\backrefsep}{, \S}
\renewcommand*{\backreftwosep}{ and \S}
\renewcommand*{\backreflastsep}{, and \S}
\renewcommand*{\backref}[1]{}% disable the plain list; use \backrefalt
\renewcommand*{\backrefalt}[4]{%
  \ifcase #1\relax
  \or {\footnotesize\itshape Cited in \S#2.}%
  \else {\footnotesize\itshape Cited in \S#2.}%
  \fi}

% ---- Theorem environments ----
\newtheorem{theorem}{Theorem}[section]
\newtheorem{fact}[theorem]{Fact}
\newtheorem{conjecture}[theorem]{Conjecture}
\theoremstyle{definition}
\newtheorem{definition}[theorem]{Definition}
\theoremstyle{remark}
\newtheorem{remark}[theorem]{Remark}

% ---- Undergraduate project boxes ----
% \begin{project}{Short title}{prerequisites}{time window} ... \end{project}
% Plain (non-\refstepcounter) counter: keeps citations inside boxes
% anchored to the enclosing section, so backref dedup works.
\newcounter{projectctr}
\newtcolorbox{project}[3]{%
  enhanced, breakable,
  code={\stepcounter{projectctr}},
  colback=projfill, colframe=projframe,
  fonttitle=\bfseries\small,
  title={Project \theprojectctr: #1},
  attach boxed title to top left={yshift=-2mm, xshift=4mm},
  boxed title style={colback=projframe},
  after upper={\par\smallskip\footnotesize\textit{Prerequisites:} #2\quad
               \textit{Window:} #3},
  left=2mm, right=2mm, top=3mm, bottom=2mm
}

% ---- Reading-spine box for variant front matter ----
% \begin{readingspine}{title} ... \end{readingspine}
\newtcolorbox{readingspine}[1]{%
  enhanced, breakable,
  colback=projfill, colframe=projframe,
  fonttitle=\bfseries\small,
  title={#1},
  attach boxed title to top left={yshift=-2mm, xshift=4mm},
  boxed title style={colback=projframe},
  left=2mm, right=2mm, top=3mm, bottom=2mm
}

% Local archived copy of a cited work: visually distinct (green) and
% clickable — opens refs/<file> relative to this PDF's location.
\newcommand{\localpdf}[1]{\href{./refs/#1}{\textcolor{archgreen}{$\blacktriangleright$\,\texttt{refs/#1}}}}

% ---- Notation macros ----
\newcommand{\Ztwo}{\mathbb{Z}/2^{32}\mathbb{Z}}
\newcommand{\Ftwo}{\mathbb{F}_2}
\newcommand{\State}{\mathcal{S}}
\newcommand{\shaxor}{\oplus}
\newcommand{\shaplus}{\boxplus}
\newcommand{\SHA}{\textsf{SHA-256}}
\newcommand{\bitstrings}{\{0,1\}}


\title{\SHA{} as a Mathematical Object:\\
       \large A Field Guide and a Program of Small Questions}
\author{[Author name here]\thanks{Working draft for mathematician
  colleagues. Comments very welcome.}}
\date{Last updated \lastupdated}

\begin{document}
\maketitle

\begin{abstract}
\noindent
The function \SHA{} maps arbitrary bit-strings of bounded length to
$256$-bit strings. It is one of the most heavily used mathematical objects
on Earth, and one of the least understood: a quarter-century of study has
produced neither a proof of any of its conjectured properties nor any
observable deviation from the behavior of a random function. This note
introduces \SHA{} historically and empirically---what it is for, why it has
lasted---and then frames it geometrically, as a self-map of the finite space
$(\Ztwo)^8$ built by deliberately alternating two incompatible group
structures on the same $2^{256}$-point set---a clash the $2$-adic metric
brings into sharp focus, sorting the machine operations by whether they
respect it and opening an unexpected line to the nonarchimedean dynamics
of $p$-adic disks. We describe the frameworks that
exist for analyzing such objects, why a large gap separates ``passes every
test we have'' from ``provably pseudorandom,'' and how the
complex-analytic tradition of Flajolet and collaborators, brought to bear
on the statistics of random mappings, offers concrete, measurable
coordinates for locating \SHA{} in the space of pseudorandom generators. We then open the adversary's own toolbox---%
differential cryptanalysis, the method that felled this function's
ancestors, with its remarkable secret history---and close with the view
from computational complexity: \textsc{sat} solvers, Impagliazzo's five
worlds, and what theorem provers did and did not prove. Throughout, we flag self-contained subprojects
suitable for undergraduates in short research windows.
\end{abstract}

\tableofcontents
\bigskip

% ------------------------------------------------------------
% Front-matter reader's note, shared by all variants. Input
% after \tableofcontents, BEFORE the first \section. Deliberately
% free of \cite commands: citations outside numbered sections
% would corrupt the section-level backrefs (both works mentioned
% here are cited properly in \cref{sec:roundshape}).
% ------------------------------------------------------------
\noindent\textit{A word on how to read what follows.} This is a
deliberately leisurely introduction: we approach \SHA{} the long way
around, with excursions into history, epistemology, finite geometry, and
analytic combinatorics, because the object of interest is not the
algorithm so much as where it sits---its context, its significance, and
the peculiar shape of what is and is not known about it. Readers
impatient for the algorithm itself should know that it is short and
entirely public: the Wikipedia page on SHA-2 contains a complete and
perfectly accessible, if dense, definition in pseudocode, and the
official standard (NIST's FIPS 180-4) is only a few pages longer; both
are cited where the round function is anatomized (\cref{sec:roundshape}).
Nothing about the definition is hidden. Everything about its behavior is.
\bigskip


% ------------------------------------------------------------
% §1  Opener: the Josephus problem as epistemological archetype
% ------------------------------------------------------------
\section{A pattern where none was promised}\label{sec:josephus}
% Save this section's hyperref anchor so citations appearing after the
% floated figure still record a section-level backref (not figure.1),
% keeping the bibliography's "Cited in §…" list free of duplicates.
\makeatletter\edef\josephussectionhref{\@currentHref}\makeatother

\emph{Concrete Mathematics} opens with a small massacre~\citep[\S 1.1]{GKP1994}.
Number $n$ people $1,2,\dots,n$ standing in a circle, and going around the
circle eliminate every second person until one survives. Which position
$J(n)$ survives? For $n=10$ the eliminations run
$2,4,6,8,10,3,7,1,9$, and the survivor is $J(10)=5$
(\cref{fig:josephus}). Computing small cases and staring at them reveals
that if $n = 2^m + \ell$ with $0 \le \ell < 2^m$, then
\begin{equation}\label{eq:josephus}
  J(2^m + \ell) \;=\; 2\ell + 1 ,
\end{equation}
which says something better in binary: \emph{$J(n)$ is the one-bit
cyclic left shift of $n$}. For $n = 10 = (1010)_2$ we get $J(n) =
(0101)_2 = 5$. An iterative algorithmic process ends up having a
solution involving a simple bit manipulation of the binary
representation of the input. (Hold on to one question as you read: what
happens if every \emph{third} person is eliminated instead? We will not
answer it until \cref{sec:josephus3}, where it becomes a parable about
the constants in \SHA{}'s own definition.)

\begin{figure}[!ht]
  \centering
  % One circle of 10, with the given people crossed out (cumulative):
  % #1 = list of position/elimination-order pairs to cross, #2 = sublabel
  \newcommand{\joscircle}[2]{%
    \begin{tikzpicture}[baseline=(current bounding box.center)]
      \foreach \i in {1,...,10} {
        \pgfmathsetmacro{\ang}{90 - (\i-1)*36}
        \node[circle, draw, minimum size=4.5mm, inner sep=0pt,
              font=\scriptsize] (p\i) at (\ang:1.35) {\i};
      }
      \foreach \i/\elim in {#1} {
        \pgfmathsetmacro{\ang}{90 - (\i-1)*36}
        \draw[gray] ($(p\i)+(-0.19,-0.19)$) -- ($(p\i)+(0.19,0.19)$);
        \node[gray, font=\tiny] at (\ang:1.8) {\elim};
      }
      \node[font=\footnotesize, align=center] at (0,-2.35) {#2};
    \end{tikzpicture}}%
  \joscircle{2/1, 4/2, 6/3, 8/4, 10/5}{first lap:\\$2,4,6,8,10$ go}%
  \,$\to$\,%
  \joscircle{2/1, 4/2, 6/3, 8/4, 10/5, 3/6, 7/7}{second lap:\\$3,7$ go}%
  \,$\to\cdots\to$\,%
  \begin{tikzpicture}[baseline=(current bounding box.center)]
    \foreach \i in {1,...,10} {
      \pgfmathsetmacro{\ang}{90 - (\i-1)*36}
      \node[circle, draw, minimum size=4.5mm, inner sep=0pt,
            font=\scriptsize] (p\i) at (\ang:1.35) {\i};
    }
    \foreach \i/\elim in {2/1, 4/2, 6/3, 8/4, 10/5, 3/6, 7/7, 1/8, 9/9} {
      \pgfmathsetmacro{\ang}{90 - (\i-1)*36}
      \draw[gray] ($(p\i)+(-0.19,-0.19)$) -- ($(p\i)+(0.19,0.19)$);
      \node[gray, font=\tiny] at (\ang:1.8) {\elim};
    }
    \node[circle, draw, very thick, fill=projfill,
          minimum size=4.5mm, inner sep=0pt, font=\scriptsize]
          at (90-4*36:1.35) {5};
    \node[font=\footnotesize, align=center] at (0,-2.35)
      {end: $5$ survives};
  \end{tikzpicture}
  \caption{Circular decimation for $n=10$, read left to right; the
  ellipsis elides the intermediate eliminations of the final lap. Small
  gray indices give the global elimination order; position $5$ survives,
  and $J(10) = 5 = (0101)_2$ is the one-bit cyclic shift of
  $n = 10 = (1010)_2$.}
  \label{fig:josephus}
\end{figure}
% Restore the section anchor so citations after this float record
% section-level backrefs (not figure.1), keeping "Cited in §…" deduped.
\makeatletter\let\@currentHref\josephussectionhref\makeatother

Dwell for a moment on what \emph{kind} of fact this is. The problem is
defined \emph{iteratively}: the survivor is specified as the outcome of
$n-1$ eliminations performed in order, and the obvious way to compute
$J(n)$ is to simulate them, one by one, around the circle. The closed
form abolishes the simulation. The answer is obtained by merely
\emph{rearranging the bits of the input}---one cyclic shift, no iteration
at all: a computation that appeared to require $n-1$ sequential steps
collapses to a single stroke. Nothing in the problem statement promised
that such a collapse existed; it was simply there, waiting to be noticed.

It will pay immediately to name this operation the way the \SHA{}
standard does~\citep{FIPS180-4}. For a $w$-bit word $x$ and $0 \le r < w$,
write
\[
  \mathrm{ROTL}^{r}(x), \qquad \mathrm{ROTR}^{r}(x), \qquad \mathrm{SHR}^{r}(x)
\]
for the cyclic left rotation, the cyclic right rotation, and the ordinary
(bit-discarding) right shift of $x$ by $r$ positions. In this notation the
Josephus law \cref{eq:josephus} is simply
\[
  J(n) \;=\; \mathrm{ROTL}^{1}(n),
\]
the rotation taken in the $(\lfloor \lg n \rfloor + 1)$-bit window
that $n$ occupies. The rotation operators are \emph{literal components
of \SHA{}}. Its round function stirs the state with fixed maps
$\Sigma_0, \Sigma_1, \sigma_0, \sigma_1$, each an exclusive-or of two
or three rotations of a $32$-bit word (\cref{sec:roundshape})---the
survivor map of a circle of $32$ bits, composed with itself in
threes. The same innocent operation that \emph{creates} the visible
structure in the Josephus problem is deployed in \SHA{} to
\emph{destroy} visible structure; and one of the very few theorems
known about \SHA{}'s internals is a statement about exactly these
xor-of-rotation maps (\cref{sec:rivest}).

\emph{Innocent discrete processes secretly respect structure}, and the
structure is usually found empirically---by computing small cases,
arranging them suggestively, and noticing. The recurrence came first; the
closed form \cref{eq:josephus} was extracted from data.

The Josephus problem is not a lone curiosity. A whole gallery of famous
iterative processes collapse the same way, each to a one-line function
of the binary digits of its input: the Tower of Hanoi, Gray codes, the
game of Nim, the parity of Pascal's triangle, even---in a sense that
must be stated carefully---the notorious $3x+1$ map.
\Cref{app:gallery} collects these specimens, in increasing order of
discomfort: three where the collapse is total, one (Kummer's carry
theorem) where the collapse names the central character of this whole
paper, and one where a collapsing change of coordinates provably
\emph{exists} and provably \emph{helps not at all}. The moral compounds
across the gallery: bit operations are not just where these collapses
land, they are a recurring \emph{language} in which discrete iteration,
when it is secretly simple, tends to become simple---and
(\cref{sec:statespace}) they are precisely the alphabet from which
\SHA{} is built.

This paper is about the mirror image of that phenomenon. The function
\SHA{}, which we will introduce slowly, is a discrete process on binary
strings that was \emph{engineered so that no analyst would ever have a
Josephus moment with it}: no change of representation, no clever
coordinate system, no generating function is supposed to reveal any
respected structure whatsoever---even though the function is built from a
few dozen lines of elementary operations, all of them individually as
transparent as the decimation rule. \SHA{} too is defined iteratively---a
fixed round transformation applied sixty-four times, itself chained block
after block---and the working conjecture is that, unlike the decimation,
its iteration admits no collapse. Since the entire document is an
extended commentary on that claim, it is worth displaying once, in the
form of a conjecture---informal, but falsifiable in every clause:

\begin{conjecture}[the no-collapse conjecture; informal]
\label{conj:nocollapse}
No change of representation tames \SHA{}: no bit-shift pattern, no
closed form, no algorithm of any kind computes it materially faster
than performing its sixty-four rounds; and no nontrivial property of
its outputs can be predicted materially more cheaply than by computing
them.
\end{conjecture}

\noindent
Nothing remotely like this is proved, and \cref{sec:worlds} explains
why a proof is beyond the current reach of mathematics. What is known
traces the conjecture's boundary. Inward, the few honest theorems about
the internals say which components, taken alone, are transparent
(\cref{thm:rivest} is the loveliest); the round-by-round records of
\cref{sec:siege} measure how far into the iteration analysis currently
penetrates. Outward, the theory of time-bounded Kolmogorov complexity
(\cref{sec:liupass}) is the formal home of ``no shortcut exists,'' and
turns the conjecture's vague ``materially faster'' into exact
statements. In between lies everything this paper does. Whether the
engineers succeeded is, remarkably, not a settled question but an
empirical one, now running as an uncontrolled worldwide experiment for
a quarter of a century. Making parts of that experiment controlled,
small, and measurable is the research program of this note.

% ------------------------------------------------------------
% §2  History, uses, and the epistemological situation
% ------------------------------------------------------------
\section{The object, observed from outside}\label{sec:history}

Before any internals, here is the object as its users see it. \SHA{} is a
single, fixed, publicly specified function
\[
  \SHA : \bigcup_{0 \,\le\, \ell \,<\, 2^{64}} \bitstrings^{\ell}
         \;\longrightarrow\; \bitstrings^{256},
\]
published by the U.S.\ National Institute of Standards and Technology in
2002 and unchanged since~\citep{FIPS180-4}. It takes a bit string of any
length below $2^{64}$---the empty string, a single bit, the concatenated
works of a library---and returns $256$ bits. (We stress \emph{bit} string:
the standard defines the function on arbitrary bit lengths, not merely on
whole bytes; software interfaces usually round to bytes, but the
mathematical object does not. This matters in \cref{sec:geometry}.)

Two set-theoretic observations frame everything that follows. The domain
is astronomically larger than the codomain, so the function is nowhere
near injective: by pigeonhole, colliding pairs exist in unimaginable
abundance---the average fiber over a $256$-bit output is itself a set of
size around $2^{2^{64}-256}$. Whether the function is \emph{surjective},
on the other hand, is not known: nobody has proved that every $256$-bit
string is attained, and more generally nothing nontrivial is proved about
the image size or about how evenly the domain is spread across the
range---the property Bellare and Kohno isolate as hash-function
\emph{balance}, showing that deviations from it would measurably
accelerate collision search~\citep{BellareKohno2004}; for \SHA{} the
balance is simply unknown. Heuristically, of course, a random function
with this domain misses a given output with probability too small to
name, and \cref{sec:randommaps} is about taking such heuristics seriously
as measurements.

The function is deterministic and involves no secrets: everyone can evaluate
it, and everyone computes the same values.\footnote{``No secrets'' means
no secret \emph{inputs}; it emphatically does not mean no design
intelligence. Every constant and structural parameter was chosen with
visible care, by a designer---the NSA---demonstrably in possession of
nonpublic cryptanalytic knowledge. The clearest episode: NIST published
the family's ancestor, now called SHA-0, in 1993; within two years the
NSA had it withdrawn, citing an undisclosed weakness, and reissued with
exactly one change, a single one-bit rotation added to the message
schedule~\citep{FIPS180-1}. In 1998 academic cryptanalysis independently
discovered a differential collision attack on SHA-0 that this tweak
happens to block~\citep{ChabaudJoux1998}, and by 2005 explicit SHA-0
collisions had been computed~\citep{BihamChen2005}, while the amended
SHA-1 stood for another decade (\cref{sec:falling}). Whatever process
selects these functions, it is not simple random choice; \cref{sec:nums}
discusses how the SHA-2 generation made its constants auditable in
response to exactly this kind of asymmetry.} It is also fast, in a way
worth pausing on: the definition is in effect a small fixed circuit of
bitwise operations and $32$-bit additions, so it is directly expressible
in hardware. A laptop core computes it at gigabytes per second; dedicated
silicon does far better---purpose-built \SHA{} ASIC designs were an
academic subject within two years of
standardization~\citep{Dadda2004}, and the economics of proof-of-work
mining (\cref{sec:uses}) later industrialized them, the arc from CPU
software to modern mining ASICs gaining roughly six orders of magnitude
in throughput per device and five in energy per hash~\citep{Taylor2017}.

What is it \emph{for}? The honest answer is: for behaving, in public, like
a random function---a function sampled uniformly at random from the set of
all functions into $\bitstrings^{256}$---while actually being a short,
fixed program. All of its uses are corollaries of facets of that one
pretension.

\subsection{A catalogue of uses}\label{sec:uses}

\begin{description}
  \item[Fingerprinting.] If $H$ behaves like a random function, two
  different documents have the same $256$-bit hash with probability
  $2^{-256}$; a hash is a practically unforgeable fingerprint. The
  version-control system Git addresses every object it stores by its hash:
  file contents, directory trees, commits. Equality of fingerprints is
  treated as equality of content---mathematically unjustified and, at
  planetary scale, empirically flawless so far.

  \item[Commitment.] Publish $H(x)$ today, reveal $x$ next month; anyone
  can verify you did not change your mind, yet the hash revealed nothing
  about $x$. A random-looking function gives sealed envelopes without
  paper. (This is how one archives a prediction, a bid, or a claim of
  priority.)

  \item[Merkle trees.] Hash the leaves; hash the concatenations of sibling
  hashes; repeat. The root is a $256$-bit commitment to an arbitrarily
  large collection, and membership of any single leaf is certified by a
  logarithmic-length path~\citep{Merkle1988,Merkle1989}. Certificate Transparency
  (the public audit log of all TLS certificates on the web) and every
  blockchain are Merkle trees at heart. Notice the epistemic shape of
  what the tree provides: the inclusion of one transaction among
  billions is \emph{proved} by a chain of a few dozen hashes---a
  certificate anyone can verify in microseconds, revealing essentially
  nothing about the other leaves. Pushed further, this seed grows into
  the theory of zero-knowledge and succinct proof
  systems~\citep{GMR1989}: modern hash-based protocols such as STARKs
  commit to the execution trace of an entire computation with a Merkle
  tree, then prove its correctness by opening a few hash
  chains~\citep{BenSasson2018}. The pattern recurs everywhere in
  mathematical cryptography and deserves explicit statement: a task---%
  here, inverting or colliding \SHA{}---is \emph{presumed} hard, and the
  presumed hardness then serves as a primitive from which richer
  structures (commitments, proofs, signatures, digital money) are built
  by explicit reduction. Hardness conjectures function the way axioms do
  elsewhere in mathematics: unproved, load-bearing, and generative.

  \item[Message authentication.] With a shared secret key $k$, the
  construction $\mathrm{HMAC}(k, m)$, built from two nested \SHA{} calls,
  authenticates messages; its analysis~\citep{BCK1996} is a rare case of a
  \emph{theorem} in this landscape---assuming the compression function
  behaves as a pseudorandom family, a notion we take up in
  \cref{sec:frameworks}.

  \item[Proof of work.] Bitcoin asks miners to find a nonce making
  $\SHA(\SHA(\text{block header}))$ land below a target
  threshold~\citep{Nakamoto2008}. If the function behaves randomly, there
  is no better strategy than exhaustive guessing, so a solution \emph{is}
  verifiable evidence of expended computation. As of this writing the
  Bitcoin network evaluates the function roughly $10^{21}$ times per
  second---surely the most intensively exercised nontrivial function in
  the history of mathematics---and this entire expenditure is, among other
  things, a brute-force search for regularities that has found none.

  \item[Key derivation and pseudorandom generation.] Operating systems
  stretch and mix entropy pools through \SHA{}; standardized deterministic
  random-bit generators iterate it. Here the function is used precisely
  \emph{as} a pseudorandom number generator, which frames the question of
  \cref{sec:frameworks}: where does it sit in the space of PRNGs?
\end{description}

Note how different the demands of these uses are: fingerprinting needs
\emph{collision resistance} (no collision $x \ne y$ with $H(x) = H(y)$
has ever been \emph{published}---though note the epistemic fine print:
were one discovered, its discoverer might have excellent reasons to keep
it private),
commitments need \emph{hiding} and \emph{binding}, proof-of-work needs
\emph{progress-freeness} (no shortcut to partial preimages), key
derivation needs \emph{pseudorandomness of outputs}. A taxonomy of these
properties and their logical relations exists~\citep{RogawayShrimpton2004},
and the standard reference surveying the whole subject at textbook depth
is the Handbook chapter of \citet[Ch.~9]{MenezesOorschotVanstone1996};
the striking fact is that one fixed function is trusted, simultaneously,
with all of them.

\subsection{The original hash functions}\label{sec:elderhash}

The word ``hash'' is decades older than its cryptographic sense, and the
two meanings are worth separating before the cryptographic one takes over
the paper. Hashing as a storage technique is nearly as old as
stored-program computing: scatter-storage schemes circulated inside IBM
from 1953 (in an internal memorandum of H.~P. Luhn), reached the open
literature with Peterson's study of random-access
addressing~\citep{Peterson1957}, and were codified, history and all, in
Knuth's chapter on searching~\citep[\S 6.4]{Knuth1998}. In that tradition
a hash function need only \emph{scatter}; collisions are a fact of life
to be chained through, probed around, and amortized away, and nobody is
trying to \emph{manufacture} one.

Here is a complete, honest member of that tradition, as a working
one-liner,
\begin{center}
  \texttt{def h(k): return k \% 997}
\end{center}
---division hashing, filing the integer key $k$ into one of $997$
buckets~\citep[\S 6.4]{Knuth1998}. Mathematically it is the reduction
map $\mathbb{Z} \to \mathbb{Z}/997\mathbb{Z}$, and because that map is a
ring homomorphism, every anatomical question one could ask of it has a
one-line proof. Its collisions are not merely abundant but
\emph{listable}: the keys colliding with $k$ are exactly the arithmetic
progression $k + 997\mathbb{Z}$. It is provably injective on a natural
piece of its domain: any $997$ consecutive keys land in $997$ distinct
buckets. It is provably surjective, with exact equidistribution: among
any $997m$ consecutive keys, every bucket is hit exactly $m$ times. And
it respects algebraic structure outright:
$h(j+k) \equiv h(j) + h(k)$ and $h(jk) \equiv h(j)\,h(k) \pmod{997}$.

How far does this transparency stretch toward \SHA{}'s actual machinery?
Further than one might guess. Trade the prime modulus for a power of two
and make the map affine,
\[
  h(k) = (a\,k + b) \bmod 2^{n}, \qquad a \text{ odd},
\]
and \emph{carries} enter the story: the multiplication and addition now
ripple information from low-order bits into high ones---the $2$-adically
continuous operations we meet again in \cref{sec:statespace}. Yet nothing
is lost to analysis. With $a$ odd the map is a \emph{bijection} of the
$n$-bit words (no longer a compressing bucket hash but a reversible
mixing map), invertible in closed form as $k = a^{-1}(h - b) \bmod 2^{n}$,
where $a^{-1}$ is the inverse of $a$ modulo $2^{n}$; iterated, it is
precisely a linear congruential generator, whose period is a classical
theorem (\cref{sec:prngs}). Now throw an exclusive-or on top,
\[
  h(k) = \bigl((a\,k + b) \bmod 2^{n}\bigr) \shaxor c,
\]
and---the pleasant surprise---it stays analyzable: an $\shaxor$ with a
constant is itself a bijection, living in the \emph{rival}
$\Ftwo$-geometry of \cref{sec:statespace}, and a composition of bijections
is a bijection, so the whole map is a reversible permutation one can still
invert by hand. What one now holds is a one-round baby \SHA{}: two of its
three ARX primitives (carry-bearing arithmetic and $\shaxor$) already in
place, and analyzable only because nothing yet forces the two geometries
to \emph{interfere}. \SHA{} adds the third primitive---bitwise rotation,
the operation that couples the geometries---and iterates the mixture
$64$ times; it is that composition, not any single ingredient, that
carries the design out of reach of every one-line argument above.

\subsection{Cryptographic hash functions}\label{sec:falling}

\SHA{} is the survivor of an explicit evolutionary lineage, and the
extinctions in that lineage are the best available evidence about what
kind of object it is. The cryptographic branch split from the storage
tradition of \cref{sec:elderhash} the moment someone set out to
\emph{manufacture} collisions rather than merely tolerate them---the
moment, that is, an adversary entered the room. One-way functions
appeared first as a systems trick for protecting stored login passwords
(in use by the late 1960s; first carefully engineered constructions
in~\citealp{Purdy1974}; the classic Unix case history
in~\citealp{MorrisThompson1979}); one-wayness was then named as a
foundational primitive of the new public-key
cryptography~\citep{DiffieHellman1976}; and the modern object---a
compressing, collision-resistant, \emph{one-way hash function}---was
isolated in Merkle's 1979 dissertation~\citep{Merkle1979}. Same word as
before, escalating demand: from ``collisions are a performance nuisance''
to ``collisions exist in abundance, but no adversary on Earth can exhibit
one.'' The escalation is a change of adversary model, not of
machinery---and it took thirty years to complete.

Put the four-question anatomy exam of \cref{sec:elderhash} in front of
\SHA{} now, and every clean answer becomes an open problem.
Collisions: they exist in astronomical abundance (pigeonhole), yet none
has ever been exhibited, and no known method materially beats generic
birthday search at $2^{128}$ evaluations (\cref{sec:telescope}).
Injectivity on a restricted domain: is \SHA{} injective on the
$2^{128}$ inputs of exactly $128$ bits? For a random function this
would be very nearly a fair coin flip---the expected number of
colliding pairs is about $\tfrac12$---and for \SHA{} nobody can improve
on that shrug in either direction. Surjectivity: unknown; a random
function of these dimensions would miss at most a fantastically small
fraction of its range, but nothing of the kind is proved. Respected
structure: the design \emph{intent} is the wholesale negation of the
homomorphism property, and ``respects nothing'' is precisely the
conjunction of unproved statements this paper keeps circling. Same
species as the one-liner, same questions; but every answer has moved
from a line of algebra to an open problem. The rest of the paper lives
in that gap.

The modern era begins with Merkle's and Damg{\aa}rd's 1989 work putting
hash functions on a design principle. Rather than design a function on
arbitrary-length inputs directly, one fixes a single \emph{compression
function} $f$ that maps a fixed-size state and one message block to a new
state, and extends it to messages of any length by iterating: split the
(padded) message into blocks $m_1, \dots, m_k$, start from a fixed
initial value, and fold the blocks through $f$ one at a time,
\[
  h_0 = \mathrm{IV}, \qquad
  h_i = f(h_{i-1}, m_i), \qquad
  H(m) = h_k,
\]
reading off the final chaining value $h_k$ as the digest. The virtue of
this construction is a \emph{theorem}: any collision in the full hash $H$
can be unwound into a collision in the single map $f$, so collision
resistance of the whole reduces to collision resistance of $f$ alone (the
Merkle--Damg{\aa}rd theorem, \cref{thm:md}). This is exactly the skeleton
\SHA{} is built on, compression function and all---the finite map $f$ and
the two theorems that wrap it are anatomized in
\cref{sec:merkledamgard,sec:constructions}. Rivest's concrete proposals
MD4 (1990) and MD5 (1992)~\citep{RFC1321} were the first mass-market
instances---$128$-bit hashes built as fast iterated compression---and MD5
was, for a decade, everywhere. NIST's SHA-1 (1995) enlarged the state to
$160$ bits and adjusted the internals; it inherited MD5's ubiquity. The
SHA-2 family---of which \SHA{} is the $256$-bit member---was standardized
in 2002 with longer digests and a reworked, more conservative round
function.

Then the towers fell, in order of age. In 2005 Wang and Yu exhibited
explicit collisions in MD5 by hand-crafted differential
analysis~\citep{WangYu2005}; within a few years refinements produced
colliding documents in seconds on a laptop, and forged certificates in the
wild. The same summer, Wang, Yin, and Yu broke SHA-1's collision
resistance on paper, at cost $2^{69}$ evaluations versus the $2^{80}$
birthday bound~\citep{WangYinYu2005}; twelve years later the attack was
industrialized: the SHAttered project published two distinct PDF files
with equal SHA-1 hashes, at a cost of roughly $2^{63}$ evaluations---about
6{,}500 CPU-years plus 100 GPU-years, real but affordable for a
well-resourced organization~\citep{Stevens2017}. Git, which had bet its
object model on SHA-1, began a still-ongoing migration to \SHA{}.

Two lessons, both load-bearing for this paper. First, \emph{the choice of
algorithm matters enormously}: MD5, SHA-1, and \SHA{} look like
siblings---the same iterated-compression skeleton, the same diet of
$32$-bit additions, rotations, and Boolean operations---yet differential
cryptanalysis (anatomized in \cref{sec:differential}) shattered the first
two while \SHA{}, attacked by the same schools with the same tools for
twenty years, has not visibly moved
(current state of the art in \cref{sec:frameworks,sec:siege}). Whatever separates
broken from unbroken here is not the general species of design; it lives
in the fine structure of round counts, rotation offsets, and message
scheduling---quantitative differences producing a qualitative divide, and
nobody can currently compute, from the specification alone, which side of
the divide a design falls on. Second, breaks are not black swans; they are
the \emph{normal} fate of these objects, which makes \SHA{}'s quarter
century of stability a phenomenon in want of explanation rather than a
default to be assumed.

\subsection{The epistemological situation}\label{sec:epistemology}

How do we know what we know about \SHA{}? Not the way we know things about
the Josephus function. There is no theorem asserting any of the properties
in \cref{sec:uses}; as we will see in \cref{sec:frameworks}, for a fixed
unkeyed function it is delicate even to \emph{state} them. What exists
instead is a corpus of observations: a published specification, a fossil
record of broken relatives (\cref{sec:falling}), a large literature of
failed and partially successful attacks on weakened versions, and an
industrial base exercising the function at cosmological rates without
surfacing an anomaly.

This is the epistemic posture of a \emph{naturalist}, not an axiomatist:
we know \SHA{} the way nineteenth-century biology knew a species---by
morphology, habitat, behavior under stress, and comparison with extinct
relatives. The cryptographic literature has even formalized the ignorance:
Rogaway's ``human ignorance'' framework~\citep{Rogaway2006} proves
theorems of the form ``any explicit method for breaking protocol $P$
yields an explicit collision in \SHA{}''---transferring trust not from an
axiom but from the observed, decades-long failure of a very motivated
species to produce such a collision. It is Popperian science embedded
inside pure mathematics: the conjecture ``\SHA{} has no exploitable
structure'' is falsifiable, massively tested, and unproven.

The rest of this paper takes the naturalist's program seriously and asks
what a \emph{quantitative} natural history of this object would look like:
first the anatomy (\cref{sec:geometry}), then the field measurements one
can actually take, where work of Flajolet supplies exactly the right
instrument (\cref{sec:randommaps}), then the existing theoretical
frameworks and the genuine gap between them (\cref{sec:frameworks}).

% ------------------------------------------------------------
% §3  The construction and its finite geometry
% ------------------------------------------------------------
\section{Anatomy: a self-map of a finite space}\label{sec:geometry}

We now open the organism. The dissection has two stages, and only the
second is specific to \SHA{}: first a completely general---and genuinely
theoretical---reduction of ``hash anything'' to ``iterate one map on a
finite space''; then the finite space itself and the map that \SHA{}
iterates on it.

\subsection{From arbitrary bit strings to one finite map}
\label{sec:merkledamgard}

A function on $\bigcup_{\ell < 2^{64}} \bitstrings^\ell$ looks
unwieldy---its domain is (for practical purposes) infinite and
ragged. The first structural fact about \SHA{} is that all of this
unwieldiness is handled by a single, provably safe reduction, the
\emph{Merkle--Damg{\aa}rd construction}, and that everything interesting
then happens on a finite set.

\paragraph{Padding.} A message $m$ of bit length $\ell$ (any $\ell$,
including $0$; bits, not bytes) is padded by appending a single
$\mathsf{1}$ bit, then the minimum number of $\mathsf{0}$ bits, then the
number $\ell$ written as a $64$-bit integer, so that the total length is a
multiple of $512$. The padded message splits into $512$-bit blocks
$m_1, \dots, m_k$ (\cref{fig:padding}). Padding is where ``arbitrary bit
length'' lives: the map
\[
  \mathrm{pad} : \bigcup_{\ell < 2^{64}} \bitstrings^{\ell}
  \;\longrightarrow\; \bigl(\bitstrings^{512}\bigr)^{+}
\]
is injective (the appended $\mathsf{1}$ and the length field between them
guarantee it), and encoding the \emph{length} in the final block---so-called
Merkle--Damg{\aa}rd strengthening---is not bookkeeping but the hypothesis
that makes \cref{thm:md} below true.

\begin{figure}[ht]
  \centering
  \begin{tikzpicture}[font=\small, >={Stealth[length=1.8mm]}]
    % ---- the linear message ----
    \draw[fill=msgfill] (0,2.2) rectangle (9.03,2.8);
    \draw (3.7,2.2) -- (3.7,2.8);
    \draw (7.4,2.2) -- (7.4,2.8);
    \node at (1.85,2.5)  {\textsf{A}};
    \node at (5.55,2.5)  {\textsf{B}};
    \node at (8.21,2.5)  {\textsf{C}};
    \node[anchor=west] at (9.25,2.5) {message $m$, $\ell$ bits};
    % ---- guide lines: where each portion lands ----
    \draw[dashed, gray] (0,2.2)    -- (0,0.62);
    \draw[dashed, gray] (3.7,2.2)  -- (3.95,0.62);
    \draw[dashed, gray] (7.4,2.2)  -- (7.9,0.62);
    \draw[dashed, gray] (9.03,2.2) -- (9.53,0.62);
    % ---- the 512-bit blocks ----
    \draw[fill=msgfill] (0,0)    rectangle (3.7,0.6);
    \draw[fill=msgfill] (3.95,0) rectangle (7.65,0.6);
    \node at (1.85,0.3) {\textsf{A}};
    \node at (5.8,0.3)  {\textsf{B}};
    % final block: tail of m | pad bit 1 | 0...0 | length
    \draw[fill=msgfill]  (7.9,0)   rectangle (9.25,0.6);
    \draw[fill=padfill]  (9.25,0)  rectangle (9.75,0.6);
    \draw[fill=black!8]  (9.75,0)  rectangle (10.85,0.6);
    \draw[fill=lenfill]  (10.85,0) rectangle (11.6,0.6);
    \draw (7.9,0) rectangle (11.6,0.6);
    \node at (8.575,0.3) {\textsf{C}};
    \node[font=\bfseries] at (9.5,0.3) {$\mathsf{1}$};
    \node at (10.3,0.3)  {$\mathsf{0}\cdots\mathsf{0}$};
    \node at (11.22,0.3) {$\langle \ell \rangle$};
    % pad bit and zero-fill annotations (staggered to avoid collision)
    \node[font=\scriptsize, align=center] (onebit) at (8.15,1.7)
      {single pad bit $\mathsf{1}$\\ \tiny(mandatory: exactly one, always
       present, marks end of message)};
    \draw[->, thin] (onebit.south) .. controls (8.9,1.15) .. (9.5,0.63);
    \node[font=\scriptsize, align=center] (zeros) at (10.9,1.15)
      {zero fill\\ \tiny($k$ zeros, $0\le k\le 511$; \emph{possibly none})};
    \draw[->, thin] (zeros.south west) -- (10.3,0.63);
    % ---- block names and sizes ----
    \node at (1.85,-0.33) {$m_1$};
    \node at (5.8,-0.33)  {$m_2$};
    \node at (9.4,-0.33) {$m_3$};
    \draw[decorate, decoration={brace, mirror, amplitude=4pt}]
      (0,-0.62) -- (3.7,-0.62)
      node[midway, below=5pt, font=\scriptsize] {512 bits};
    \draw[decorate, decoration={brace, mirror, amplitude=3pt}]
      (10.85,-0.08) -- (11.6,-0.08)
      node[midway, below=4pt, font=\scriptsize, align=center]
      {$64$ bits:\\bit length $\ell$};
  \end{tikzpicture}
  \caption{Padding and blocking, for a message needing $k=3$ blocks. The
  linear message (blue; portions \textsf{A}, \textsf{B}, \textsf{C})
  is consumed left to right by $512$-bit blocks. The final block is
  completed by a single $\mathsf{1}$ bit (green), a run of $\mathsf{0}$s,
  and the $64$-bit binary encoding of the message length $\ell$
  (orange)---the Merkle--Damg{\aa}rd strengthening. The rule is exact and
  admits no options: to a message of $\ell$ bits one \emph{always} appends
  one $\mathsf{1}$, then the fewest $\mathsf{0}$s (possibly none) so that
  the total reaches a multiple of $512$ once the $64$-bit length is
  affixed---i.e.\ $k$ is the least $\ge 0$ solution of
  $\ell + 1 + k + 64 \equiv 0 \pmod{512}$. The $\mathsf{1}$ bit and the
  length field are never omitted; only the zero-count $k$ depends on
  $\ell$, and if the tail \textsf{C} leaves no room for the $65$ closing
  bits, $k$ simply runs on into an extra all-padding block.}
  \label{fig:padding}
\end{figure}

\paragraph{Chaining.} All real work is delegated to one
\emph{compression function}
\[
  f : \State \times \bitstrings^{512} \longrightarrow \State,
  \qquad \State = (\Ztwo)^8 ,
\]
which eats a $512$-bit block and updates a point of the finite space
$\State$ (about which everything in \cref{sec:statespace}). Fix a public
starting point $h_0 = \mathrm{IV} \in \State$ and set
\[
  h_i = f(h_{i-1}, m_i), \qquad \SHA(m) = h_k ,
\]
reading the final state out as $256$ bits (\cref{fig:md}).

\begin{figure}[ht]
  \centering
  \begin{tikzpicture}[
      node distance=9mm,
      blk/.style={draw, minimum width=13mm, minimum height=7mm, font=\small},
      f/.style={draw, fill=projfill, minimum size=8mm, font=\small},
      >={Stealth[length=2mm]}]
    \node[f] (f1) {$f$};
    \node[f, right=16mm of f1] (f2) {$f$};
    \node[right=10mm of f2, font=\small] (dots) {$\cdots$};
    \node[f, right=10mm of dots] (f3) {$f$};
    \node[blk, above=7mm of f1] (m1) {$m_1$};
    \node[blk, above=7mm of f2] (m2) {$m_2$};
    \node[blk, above=7mm of f3, fill=black!8] (m3) {$m_k$};
    \node[left=9mm of f1, font=\small] (iv) {$\mathrm{IV}$};
    \node[right=9mm of f3, font=\small] (out) {$\SHA(m)$};
    \draw[->] (iv) -- (f1);
    \draw[->] (f1) -- node[below, font=\scriptsize]{$h_1$} (f2);
    \draw[->] (f2) -- (dots);
    \draw[->] (dots) -- node[below, font=\scriptsize]{$h_{k-1}$} (f3);
    \draw[->] (f3) -- (out);
    \draw[->] (m1) -- (f1);
    \draw[->] (m2) -- (f2);
    \draw[->] (m3) -- (f3);
    \node[above=1mm of m3, font=\scriptsize, align=center]
      {last block contains\\the bit length $\ell$};
  \end{tikzpicture}
  \caption{The Merkle--Damg{\aa}rd construction: an arbitrary-length
  message is folded through one compression function
  $f : \State \times \bitstrings^{512} \to \State$ on the finite state
  space $\State = (\Ztwo)^8$.}
  \label{fig:md}
\end{figure}

This is not an ad hoc recipe; it is a construction with a theorem, proved
independently in 1989 by Damg{\aa}rd---stated and proved as Theorem~3.1
of~\citet[pp.~419--420]{Damgard1989}---and by Merkle, under the name
``meta-method,'' in~\citet{MerkleOWHF1989}. Both original papers are
short, readable, and archived alongside this note; the theorem is exactly
why the design has been reused by MD5, SHA-1, and SHA-2 alike:

\begin{theorem}[Merkle, Damg{\aa}rd]\label{thm:md}
Any collision of the iterated hash yields, constructively, a collision of
the compression function: from $m \ne m'$ with $\SHA(m) = \SHA(m')$ one
can extract $(h, b) \ne (h', b')$ with $f(h,b) = f(h',b')$.
\end{theorem}

\begin{proof}[Proof idea]
Walk the two chains backward from the equal outputs. If the messages have
different lengths, the final blocks already differ (they encode the
lengths---this is where strengthening is used) yet map to the same output:
a collision in $f$. If the lengths agree, compare positions from the end;
at the last position where the inputs to $f$ differ, the outputs agree,
again a collision in $f$.
\end{proof}

The moral: \emph{the arbitrary-length-input problem reduces, with proof,
to the study of one map on one finite space}. (The domain is itself
finite---every allowed input is shorter than $2^{64}$ bits---but its
inputs run through unboundedly many \emph{lengths}, and it is that
length-variability, not any actual infinity, that the construction tames.)
The reduction is airtight only for
collision resistance---the construction leaks other structure, such as the
famous length-extension property (from the digest $H(m)$ alone, without
knowing $m$, one can compute $H(m \Vert \mathrm{pad} \Vert m')$; defined
and dissected in \cref{sec:models}), whose repair spawned its own theory
(indifferentiability) and whose avoidance shaped
the post-2000 designs, including the sponge construction underlying
SHA-3~\citep{BDPV2007}. But it means the rest of this paper may honestly
focus on the finite object $f$.

\subsection{The state space and its two rival geometries}
\label{sec:statespace}

The state space is
\[
  \State \;=\; (\Ztwo)^{8},
  \qquad |\State| = 2^{256} :
\]
an $8$-dimensional module-like array of $32$-bit words---the
``$8{\times}32$'' geometry in which the entire life of \SHA{} unfolds. As
a plain set, $\State$ is also $\Ftwo^{256}$, a $256$-dimensional vector
space over the two-element field. Everything hinges on the fact that
\emph{the same underlying set carries two incompatible abelian group
structures}, both used incessantly:
\begin{itemize}
  \item bitwise addition $\shaxor$ (exclusive or): the $\Ftwo$-vector
  space structure of $\Ftwo^{256}$;
  \item wordwise addition $\shaplus$ modulo $2^{32}$: the product group
  structure of $(\Ztwo)^8$, where carries propagate within each word.
\end{itemize}
These two structures share the zero element and almost nothing else. A map
that is linear for $\shaxor$ (a matrix over $\Ftwo$) is, generically,
maximally nonlinear from the point of view of $\shaplus$, and vice versa;
the carry chains of integer addition are precisely the terms that
$\Ftwo$-linear algebra cannot see.

It helps to picture the additive ($\shaplus$) structure literally, as a
clock (\cref{fig:odometer}). Read the state as an odometer counting
modulo $2^{256}$: every modular addition advances the needle by some
amount, and the crucial event is whether it \emph{wraps} past $0$---an
overflow, a carry out of the top, thrown away by the reduction. That one
bit of ``did it go around?'' is exactly what $\shaxor$-algebra is blind to
and what the $2$-adic metric, next, is built to record.

\begin{figure}[ht]
  \centering
  \begin{tikzpicture}[font=\small, >={Stealth[length=2mm]}]
    \def\R{1.75}
    \draw[thick] (0,0) circle (\R);
    \fill (90:\R) circle (1.7pt);
    \node[above=1pt] at (90:\R) {\scriptsize $0 \equiv 2^{256}$};
    \fill (-90:\R) circle (1.3pt);
    \node[below=1pt] at (-90:\R) {\scriptsize $2^{255}$};
    \node[align=center, font=\scriptsize] at (0,0)
      {state as a\\ counter\\ $\bmod\; 2^{256}$};
    % arc 1: no wrap (does not cross the top marker at 90)
    \fill[projframe] (25:\R) circle (1.5pt);
    \draw[->, thick, projframe] (25:{\R+0.3}) arc (25:-120:{\R+0.3});
    \node[projframe, font=\scriptsize, align=center] at (-52:{\R+1.25})
      {$x \shaplus c_1$:\\ stays short of $2^{256}$\\ \emph{no carry}};
    % arc 2: wraps past the top marker (crosses 90)
    \fill[citewine] (145:\R) circle (1.5pt);
    \draw[->, thick, citewine] (145:{\R+0.3}) arc (145:35:{\R+0.3});
    \node[citewine, font=\scriptsize, align=center] at (130:{\R+1.5})
      {$x \shaplus c_2$: sweeps\\ past $0$ --- \emph{carry out}\\ (the wrap)};
  \end{tikzpicture}
  \caption{Wordwise addition $\shaplus$ as motion on a clock. As an
  odometer counting mod $2^{256}$, the state advances by each addition;
  it \emph{wraps} past $0$ exactly when the sum overflows---a carry out of
  the top bit, discarded by the reduction. Whether a given addition wraps
  is invisible to $\Ftwo$-linear ($\shaxor$) algebra and is precisely the
  information the $2$-adic metric records. (The true state is not one
  circle of $2^{256}$ positions but eight \emph{coupled} clocks of
  $2^{32}$ each, $(\Ztwo)^8$: the single odometer is the intuition, the
  product of eight is the reality.) A round is a great many such
  additions, and iterating \SHA{} (\cref{sec:iterated}) sends the counter
  wandering---sometimes around the circle, sometimes not.}
  \label{fig:odometer}
\end{figure}

There is a third, more analytic way to say this, and it is the one we
expect resonates with a $p$-adic sensibility. Identify $\Ztwo$ with the
quotient $\mathbb{Z}_2 / 2^{32}\mathbb{Z}_2$ of the $2$-adic integers.
Then the basic machine operations sort themselves by how they treat the
$2$-adic metric (in which two integers are close when they agree in many
low-order bits):
\begin{center}
  \begin{tabular}{@{}lccc@{}}
    \toprule
    operation & $\Ftwo$-linear? & $2$-adically $1$-Lipschitz? & where it is ``tame'' \\
    \midrule
    $x \shaxor c$, bitwise $\wedge,\vee,\lnot$ & yes ($\shaxor$) / no & yes & both lenses see it \\
    $x \shaplus c$, \ $x \shaplus y$          & no  & yes & only the $2$-adic lens \\
    rotation $x \lll r$, shift $x \gg s$       & yes & no  & only the $\Ftwo$ lens \\
    \bottomrule
  \end{tabular}
\end{center}
Bitwise logic and modular addition are \emph{$1$-Lipschitz} ($2$-adically
continuous in the strongest sense: output bit $i$ depends only on input
bits $0,\dots,i$, the triangular or ``causal'' maps); compositions of
these are the \emph{T-functions} of Klimov and
Shamir~\citep{KlimovShamir2003}, and there is a genuine ergodic theory,
due to Anashin, characterizing when such maps act
measure-preservingly or ergodically on $\mathbb{Z}_2$~\citep{Anashin2006}.
Rotations, by contrast, are $\Ftwo$-linear but violently discontinuous
$2$-adically: they let high-order bits influence low-order ones.

The $2$-adic lens is not an ad hoc bookkeeping device; it opens onto a
genuine subject. Nonarchimedean dynamics---the iteration theory of
self-maps of $p$-adic disks, opened up by \citet{Lubin1994}---asks
exactly this paper's kind of question one level of structure up: when
must the iteration of a $p$-adic map secretly respect algebraic
structure? Lubin's paper records a suspicion of exactly the Josephus
flavor---that a noninvertible series commuting with an invertible
nontorsion one must have a formal group hidden behind it---and its main
result makes a sharp instance of that suspicion into a theorem
(\citealp{Lubin1994}, Main Theorem~6.3):

\begin{theorem}[Lubin]\label{thm:lubin}
Let $o$ be the ring of integers in a finite extension of $\mathbb{Q}_p$,
with maximal ideal $\mathfrak m$, and let $u, f \in o[[x]]$ be power
series fixing $0$, with $u$ \emph{invertible} (its linear coefficient
$u'(0)$ a unit) and $f$ \emph{noninvertible} ($f'(0)\in\mathfrak m$) of
finite Weierstrass degree $\delta$---the number of zeros $f$ has in the
open disk, counted with multiplicity. If $u$ and $f$ commute,
$u\circ f = f\circ u$, then \emph{either} some iterate of $u$ is the
identity series ($u$ is a root of unity in the group of invertible
series) \emph{or} $\delta = p^{d}$ for some $d\ge 1$.
\end{theorem}

The content is rigidity. The only Weierstrass degrees a noninvertible
map may carry while sharing its orbits with a genuine (nontorsion)
invertible symmetry are the powers of $p$---and those are precisely the
degrees of the endomorphisms of a one-dimensional formal group. So bare
commutation, a purely dynamical hypothesis with no algebra in it, already
forces the numerical signature of a formal-group skeleton; whether the
skeleton can ever be truly absent is Lubin's own still-open question.

The phenomenon is visible in miniature over our own ring $\mathbb{Z}_2$
($p=2$). The noninvertible $f(x) = 4x + x^{2}$---linear coefficient
$4\in 2\mathbb{Z}_2$, Weierstrass degree $2 = 2^{1}$, a $p$-power on the
nose---commutes with the invertible $u(x) = 9x + 6x^{2} + x^{3}$ (unit
linear coefficient $9$), as one checks by expanding
$u\circ f = f\circ u$ directly. The companion Corollary~3.2.1 of the same
paper is just as concrete: every zero of $f'$ inside the disk must be a
zero of some iterate of $f$. Here $f'(x) = 2x + 4$ vanishes at $x=-2$,
and indeed $f(f(-2)) = f(-4) = 0$. Two lines of arithmetic thus exhibit a
symmetry ($u$), a forced degree ($2^{1}$), and a forced coincidence among
the iterates---structure precipitating out of commutation alone, the
Josephus moral in $p$-adic dress.

How near does this reach \SHA{}? Adjacent, honestly, not overhead.
\Cref{thm:lubin} governs \emph{analytic} series on the $p$-adic disk,
whereas \SHA{}'s ARX pieces are merely $1$-Lipschitz T-functions---the
coarse, non-analytic relatives tabulated above---so the theorem does not
bind the function on the nose, and Anashin's ergodic theory is what
survives at that lower regularity~\citep{Anashin2006, KlimovShamir2003}.
What \emph{does} transfer is the question and its shape: does any
nontrivial map commute with a (reduced, word-narrowed) \SHA{} component,
and if one did, would it drag in hidden algebraic structure the way it
must for analytic series? At toy scale that is a finite, machine-checkable
hunt---search for a commuting symmetry; find none, and the absence is
itself a measurement of structurelessness. In this language \SHA{}'s
designers are in the business of manufacturing $1$-Lipschitz dynamics
that commute with nothing and hide no formal group, and \cref{thm:lubin}
is the evidence that, one regularity class up, such structurelessness is
genuinely hard to achieve by accident.

\SHA{}'s round function is an \emph{ARX} design---Addition, Rotation,
Xor---and the table explains the acronym's cunning: each primitive is
transparent in one geometry and wild in the other, and the design
alternates them so that the composition is wild in both. An analyst who
adopts $\Ftwo$-linear algebra kills the X's and R's but is left helpless
against the carries; an analyst who adopts the $2$-adic/T-function lens
domesticates additions and logic but is destroyed by the rotations. This
is not an accident of \SHA{}; it is the design principle of the whole ARX
family. What \emph{is} specific to \SHA{} is the exact schedule of
alternation, and \cref{sec:falling} showed how much such specifics matter.

\begin{project}{Single operations under the $2$-adic lens}{elementary
  number theory; curiosity about $p$-adic numbers}{3--4 weeks}
  Anashin's theory gives clean criteria for when a T-function like
  $x \mapsto x \shaplus c$ or $x \mapsto x \shaxor c \shaplus (x \wedge d)$
  is a bijection of $\Ztwo[k]$ for all $k$, and when its iteration visits
  every residue (ergodicity)~\citep{Anashin2006, KlimovShamir2003}.
  Verify the criteria experimentally on $8$- and $16$-bit words, visualize
  orbits, then break the hypotheses with a single rotation and measure
  what dies. A perfect first encounter with $p$-adic ideas producing
  visible, computable phenomena.
\end{project}

\subsection{The compression function, structurally}
\label{sec:roundshape}

We describe $f$ at the level of shape; the full specification fits on two
pages of \citet{FIPS180-4} (with an accessible pseudocode rendering
in~\citealp{WikipediaSHA2}) and is worth reading only after the shape is
clear. The state is eight words $(a,b,c,d,e,f,g,h)$. The $512$-bit message
block is first \emph{expanded}: its sixteen words $W_0, \dots, W_{15}$
are stretched to sixty-four by the recurrence
\[
  W_t \;=\; \sigma_1(W_{t-2}) \shaplus W_{t-7} \shaplus
            \sigma_0(W_{t-15}) \shaplus W_{t-16},
  \qquad 16 \le t < 64,
\]
where $\sigma_0, \sigma_1$ are fixed $\shaxor$-combinations of rotations
and shifts of a single word ($\Ftwo$-linear), while the additions are
modular (2-adic). Then sixty-four rounds are applied; round $t$ computes
\[
  T_1 = h \shaplus \Sigma_1(e) \shaplus \mathrm{Ch}(e,f,g)
          \shaplus K_t \shaplus W_t,
  \qquad
  T_2 = \Sigma_0(a) \shaplus \mathrm{Maj}(a,b,c),
\]
and shifts the state pipeline:
$(a, \dots, h) \leftarrow
 (T_1 \shaplus T_2,\; a,\; b,\; c,\; d \shaplus T_1,\; e,\; f,\; g)$.
Here $\Sigma_0, \Sigma_1$ are again xors of three rotations, and
\[
  \mathrm{Ch}(x,y,z) = (x \wedge y) \shaxor (\lnot x \wedge z),
  \qquad
  \mathrm{Maj}(x,y,z) = (x \wedge y) \shaxor (x \wedge z) \shaxor (y \wedge z)
\]
are bitwise Boolean functions of degree two (``choose'' and ``majority'').
Finally---a crucial asymmetry called the Davies--Meyer feedforward---the
input state is added back: the block's net effect is
$h \mapsto h \shaplus \tilde{f}(h, m)$, which prevents running the rounds
backward even though each individual round is invertible. This
feedforward is not a detail but the hinge on which one-wayness turns, and
it is one of the two places where \SHA{}'s architecture carries a genuine
theorem; we return to it in \cref{sec:constructions}.

\Cref{fig:schedule} traces how a single block's data reaches every one of
the sixty-four rounds. The block supplies only the first sixteen schedule
words; the recurrence then \emph{stretches} those seeds forward, each new
word reaching back to four earlier ones, until all sixty-four exist---so a
bit set in the incoming block propagates, through the schedule alone,
into arbitrarily late rounds. Each round $t$ then draws exactly two
external inputs, its schedule word $W_t$ and its constant $K_t$, and folds
them into the running eight-word state; everything else in the round is
internal feedback.

\begin{figure}[ht]
  \centering
  \begin{tikzpicture}[font=\small, >={Stealth[length=1.7mm]},
      seed/.style={draw, fill=msgfill, minimum width=3.1mm,
                   minimum height=5mm, inner sep=0pt},
      der/.style={draw, fill=projfill, minimum width=3.1mm,
                  minimum height=5mm, inner sep=0pt},
      rnd/.style={draw, fill=projfill, minimum width=7mm,
                  minimum height=6mm, inner sep=0pt, font=\scriptsize},
      lbl/.style={font=\scriptsize}]
    % ---- the 512-bit block ----
    \draw[fill=msgfill] (0,5.0) rectangle (5.28,5.5);
    \node[lbl] at (2.64,5.25) {$512$-bit message block $m$};
    % ---- 16 seed words W0..W15 ----
    \foreach \i in {0,...,15} {
      \node[seed] (w\i) at ({\i*0.33+0.165},4.2) {};
      \draw[->, gray] ({\i*0.33+0.165},4.98) -- (w\i.north);
    }
    \node[lbl, below=0pt of w0] {$W_0$};
    \node[lbl, below=0pt of w15] {$W_{15}$};
    \node[lbl, align=center] at (2.5,3.55) {\scriptsize\emph{16 seed words (from block)}};
    % ---- derived words W16 .. W63 ----
    \node[der, fill=lenfill] (w16) at (6.9,4.2) {};
    \node[lbl, below=0pt of w16] {$W_{16}$};
    \node[der] (w17) at (7.25,4.2) {};
    \node[lbl] at (7.75,4.2) {$\cdots$};
    \node[der] (w63) at (8.35,4.2) {};
    \node[lbl, below=0pt of w63] {$W_{63}$};
    \node[lbl, align=center] at (7.6,3.55) {\scriptsize\emph{48 derived}};
    % fan-in into W16 from W0,W1(sigma0),W9,W14(sigma1)
    \draw[->, projframe] (w0.north) ++(0,0.02) to[bend left=32] (w16.north west);
    \draw[->, projframe] (w1.north) to[bend left=30]
      node[above, lbl, pos=0.62] {$\sigma_0$} (w16.north);
    \draw[->, projframe] (w9.north) to[bend left=22] (w16.north);
    \draw[->, projframe] (w14.north) to[bend left=16]
      node[above, lbl, pos=0.5] {$\sigma_1$} (w16.north east);
    \node[lbl, align=center, anchor=south] at (3.2,5.95)
      {each derived word reaches back four:
       $W_t=\sigma_1(W_{t-2})\shaplus W_{t-7}\shaplus\sigma_0(W_{t-15})\shaplus W_{t-16}$};
    % ---- bridge arrow ----
    \node[lbl, anchor=east] at (9.9,3.35)
      {\emph{the $64$ words feed the $64$ rounds, one each}};
    \draw[->, thick] (5.0,3.55) -- (5.0,2.75);
    % ---- round pipeline ----
    \node[lbl, anchor=east] at (-0.2,1.9) {$\mathrm{IV}$};
    \draw[->, very thick, black!55] (-0.2,1.9) -- (0.28,1.9);
    \foreach \i/\lab/\x in {0/{$R_0$}/0.9, 1/{$R_1$}/2.35, 2/{$R_t$}/4.8,
                            3/{$R_{63}$}/7.4} {
      \node[rnd, minimum width=9mm] (r\i) at (\x,1.9) {\lab};
    }
    \node[lbl] at (3.575,1.9) {$\cdots$};
    \node[lbl] at (6.1,1.9) {$\cdots$};
    \draw[->, very thick, black!55] (r0.east) -- (r1.west);
    \draw[->, very thick, black!55] (r1.east) -- (3.2,1.9);
    \draw[->, very thick, black!55] (3.95,1.9) -- (r2.west);
    \draw[->, very thick, black!55] (r2.east) -- (5.75,1.9);
    \draw[->, very thick, black!55] (6.45,1.9) -- (r3.west);
    \draw[->, very thick, black!55] (r3.east) -- (8.15,1.9)
      node[lbl, anchor=west, black] {$h_{\mathrm{out}}$};
    \node[lbl, black!55, anchor=north] at (3.7,1.42)
      {the eight-word state $(a,\dots,h)$, threaded left to right};
    % W (wine) and K (navy) inputs into each shown round
    \foreach \r/\wl/\kl in {r0/{$W_0$}/{$K_0$}, r1/{$W_1$}/{$K_1$},
                            r2/{$W_t$}/{$K_t$}, r3/{$W_{63}$}/{$K_{63}$}} {
      \draw[->, citewine] ($(\r.north)+(-0.28,0.6)$) -- ($(\r.north)+(-0.28,0.02)$);
      \node[lbl, citewine, anchor=south] at ($(\r.north)+(-0.28,0.58)$) {\wl};
      \draw[->, linknavy] ($(\r.north)+(0.28,0.6)$) -- ($(\r.north)+(0.28,0.02)$);
      \node[lbl, linknavy, anchor=south] at ($(\r.north)+(0.28,0.58)$) {\kl};
    }
  \end{tikzpicture}
  \caption{How one block's data reaches every round. The block fills only
  the sixteen seed words $W_0,\dots,W_{15}$ (blue); the message-schedule
  recurrence then stretches them into $W_{16},\dots,W_{63}$ (each derived
  word, shown for $W_{16}$, built from four earlier ones via the
  $\Ftwo$-linear maps $\sigma_0,\sigma_1$ and modular addition). Each of
  the $64$ rounds $R_t$ consumes exactly two external inputs---its
  schedule word $W_t$ (wine) and its constant $K_t$ (navy)---and threads
  the eight-word state through from the incoming chaining value to
  $h_{\mathrm{out}}$. Arrows show influence: a bit of the block, via the
  schedule's reach-back, can steer arbitrarily late rounds.}
  \label{fig:schedule}
\end{figure}

The reach of that schedule is worth making quantitative, because it is a
clean computed pattern rather than a hand-drawn cartoon. Trace, for each
word $W_t$, the set of \emph{seed} words $W_0,\dots,W_{15}$---the only
ones carrying message data---that it transitively depends on, following
the recurrence's four back-pointers ($t-2,\,t-7,\,t-15,\,t-16$) down to
the data layer. The dependency graph
(\cref{fig:schedulegraph}) then answers ``how fast does a block's data
spread through the schedule?'' exactly: the count of seeds reached climbs
$4,4,7,7,10,10,13,14,15,15,\dots$ and \emph{saturates at $W_{26}$}---from
that word on, every schedule word depends on all sixteen data words. Ten
steps of the recurrence suffice to fully entangle the message, before a
single one of the sixty-four rounds has done its diffusion work. The
four coloured arc families in the figure are the four terms of the
recurrence, and their regular weave is the schedule's linear-feedback
skeleton laid bare; whether that visible regularity is a foothold or a
red herring is one more question small experiments can chew on.

\begin{figure}[ht]
  \centering
  \resizebox{\linewidth}{!}{% AUTO-GENERATED by tools/wschedule_graph.py -- do not edit by hand.
\begin{tikzpicture}[font=\scriptsize, >={Stealth[length=1.4mm]}]
  \def\dx{0.205}
  \draw[citewine, line width=0.2pt, opacity=0.4] (2.870,0) to[out=58,in=122] (3.280,0);
  \draw[projframe, line width=0.2pt, opacity=0.4] (1.845,0) to[out=58,in=122] (3.280,0);
  \draw[linknavy, line width=0.2pt, opacity=0.4] (0.205,0) to[out=58,in=122] (3.280,0);
  \draw[orange!80!black, line width=0.2pt, opacity=0.4] (0.000,0) to[out=58,in=122] (3.280,0);
  \draw[citewine, line width=0.2pt, opacity=0.4] (3.075,0) to[out=58,in=122] (3.485,0);
  \draw[projframe, line width=0.2pt, opacity=0.4] (2.050,0) to[out=58,in=122] (3.485,0);
  \draw[linknavy, line width=0.2pt, opacity=0.4] (0.410,0) to[out=58,in=122] (3.485,0);
  \draw[orange!80!black, line width=0.2pt, opacity=0.4] (0.205,0) to[out=58,in=122] (3.485,0);
  \draw[citewine, line width=0.2pt, opacity=0.4] (3.280,0) to[out=58,in=122] (3.690,0);
  \draw[projframe, line width=0.2pt, opacity=0.4] (2.255,0) to[out=58,in=122] (3.690,0);
  \draw[linknavy, line width=0.2pt, opacity=0.4] (0.615,0) to[out=58,in=122] (3.690,0);
  \draw[orange!80!black, line width=0.2pt, opacity=0.4] (0.410,0) to[out=58,in=122] (3.690,0);
  \draw[citewine, line width=0.2pt, opacity=0.4] (3.485,0) to[out=58,in=122] (3.895,0);
  \draw[projframe, line width=0.2pt, opacity=0.4] (2.460,0) to[out=58,in=122] (3.895,0);
  \draw[linknavy, line width=0.2pt, opacity=0.4] (0.820,0) to[out=58,in=122] (3.895,0);
  \draw[orange!80!black, line width=0.2pt, opacity=0.4] (0.615,0) to[out=58,in=122] (3.895,0);
  \draw[citewine, line width=0.2pt, opacity=0.4] (3.690,0) to[out=58,in=122] (4.100,0);
  \draw[projframe, line width=0.2pt, opacity=0.4] (2.665,0) to[out=58,in=122] (4.100,0);
  \draw[linknavy, line width=0.2pt, opacity=0.4] (1.025,0) to[out=58,in=122] (4.100,0);
  \draw[orange!80!black, line width=0.2pt, opacity=0.4] (0.820,0) to[out=58,in=122] (4.100,0);
  \draw[citewine, line width=0.2pt, opacity=0.4] (3.895,0) to[out=58,in=122] (4.305,0);
  \draw[projframe, line width=0.2pt, opacity=0.4] (2.870,0) to[out=58,in=122] (4.305,0);
  \draw[linknavy, line width=0.2pt, opacity=0.4] (1.230,0) to[out=58,in=122] (4.305,0);
  \draw[orange!80!black, line width=0.2pt, opacity=0.4] (1.025,0) to[out=58,in=122] (4.305,0);
  \draw[citewine, line width=0.2pt, opacity=0.4] (4.100,0) to[out=58,in=122] (4.510,0);
  \draw[projframe, line width=0.2pt, opacity=0.4] (3.075,0) to[out=58,in=122] (4.510,0);
  \draw[linknavy, line width=0.2pt, opacity=0.4] (1.435,0) to[out=58,in=122] (4.510,0);
  \draw[orange!80!black, line width=0.2pt, opacity=0.4] (1.230,0) to[out=58,in=122] (4.510,0);
  \draw[citewine, line width=0.2pt, opacity=0.4] (4.305,0) to[out=58,in=122] (4.715,0);
  \draw[projframe, line width=0.2pt, opacity=0.4] (3.280,0) to[out=58,in=122] (4.715,0);
  \draw[linknavy, line width=0.2pt, opacity=0.4] (1.640,0) to[out=58,in=122] (4.715,0);
  \draw[orange!80!black, line width=0.2pt, opacity=0.4] (1.435,0) to[out=58,in=122] (4.715,0);
  \draw[citewine, line width=0.2pt, opacity=0.4] (4.510,0) to[out=58,in=122] (4.920,0);
  \draw[projframe, line width=0.2pt, opacity=0.4] (3.485,0) to[out=58,in=122] (4.920,0);
  \draw[linknavy, line width=0.2pt, opacity=0.4] (1.845,0) to[out=58,in=122] (4.920,0);
  \draw[orange!80!black, line width=0.2pt, opacity=0.4] (1.640,0) to[out=58,in=122] (4.920,0);
  \draw[citewine, line width=0.2pt, opacity=0.4] (4.715,0) to[out=58,in=122] (5.125,0);
  \draw[projframe, line width=0.2pt, opacity=0.4] (3.690,0) to[out=58,in=122] (5.125,0);
  \draw[linknavy, line width=0.2pt, opacity=0.4] (2.050,0) to[out=58,in=122] (5.125,0);
  \draw[orange!80!black, line width=0.2pt, opacity=0.4] (1.845,0) to[out=58,in=122] (5.125,0);
  \draw[citewine, line width=0.2pt, opacity=0.4] (4.920,0) to[out=58,in=122] (5.330,0);
  \draw[projframe, line width=0.2pt, opacity=0.4] (3.895,0) to[out=58,in=122] (5.330,0);
  \draw[linknavy, line width=0.2pt, opacity=0.4] (2.255,0) to[out=58,in=122] (5.330,0);
  \draw[orange!80!black, line width=0.2pt, opacity=0.4] (2.050,0) to[out=58,in=122] (5.330,0);
  \draw[citewine, line width=0.2pt, opacity=0.4] (5.125,0) to[out=58,in=122] (5.535,0);
  \draw[projframe, line width=0.2pt, opacity=0.4] (4.100,0) to[out=58,in=122] (5.535,0);
  \draw[linknavy, line width=0.2pt, opacity=0.4] (2.460,0) to[out=58,in=122] (5.535,0);
  \draw[orange!80!black, line width=0.2pt, opacity=0.4] (2.255,0) to[out=58,in=122] (5.535,0);
  \draw[citewine, line width=0.2pt, opacity=0.4] (5.330,0) to[out=58,in=122] (5.740,0);
  \draw[projframe, line width=0.2pt, opacity=0.4] (4.305,0) to[out=58,in=122] (5.740,0);
  \draw[linknavy, line width=0.2pt, opacity=0.4] (2.665,0) to[out=58,in=122] (5.740,0);
  \draw[orange!80!black, line width=0.2pt, opacity=0.4] (2.460,0) to[out=58,in=122] (5.740,0);
  \draw[citewine, line width=0.2pt, opacity=0.4] (5.535,0) to[out=58,in=122] (5.945,0);
  \draw[projframe, line width=0.2pt, opacity=0.4] (4.510,0) to[out=58,in=122] (5.945,0);
  \draw[linknavy, line width=0.2pt, opacity=0.4] (2.870,0) to[out=58,in=122] (5.945,0);
  \draw[orange!80!black, line width=0.2pt, opacity=0.4] (2.665,0) to[out=58,in=122] (5.945,0);
  \draw[citewine, line width=0.2pt, opacity=0.4] (5.740,0) to[out=58,in=122] (6.150,0);
  \draw[projframe, line width=0.2pt, opacity=0.4] (4.715,0) to[out=58,in=122] (6.150,0);
  \draw[linknavy, line width=0.2pt, opacity=0.4] (3.075,0) to[out=58,in=122] (6.150,0);
  \draw[orange!80!black, line width=0.2pt, opacity=0.4] (2.870,0) to[out=58,in=122] (6.150,0);
  \draw[citewine, line width=0.2pt, opacity=0.4] (5.945,0) to[out=58,in=122] (6.355,0);
  \draw[projframe, line width=0.2pt, opacity=0.4] (4.920,0) to[out=58,in=122] (6.355,0);
  \draw[linknavy, line width=0.2pt, opacity=0.4] (3.280,0) to[out=58,in=122] (6.355,0);
  \draw[orange!80!black, line width=0.2pt, opacity=0.4] (3.075,0) to[out=58,in=122] (6.355,0);
  \draw[citewine, line width=0.2pt, opacity=0.4] (6.150,0) to[out=58,in=122] (6.560,0);
  \draw[projframe, line width=0.2pt, opacity=0.4] (5.125,0) to[out=58,in=122] (6.560,0);
  \draw[linknavy, line width=0.2pt, opacity=0.4] (3.485,0) to[out=58,in=122] (6.560,0);
  \draw[orange!80!black, line width=0.2pt, opacity=0.4] (3.280,0) to[out=58,in=122] (6.560,0);
  \draw[citewine, line width=0.2pt, opacity=0.4] (6.355,0) to[out=58,in=122] (6.765,0);
  \draw[projframe, line width=0.2pt, opacity=0.4] (5.330,0) to[out=58,in=122] (6.765,0);
  \draw[linknavy, line width=0.2pt, opacity=0.4] (3.690,0) to[out=58,in=122] (6.765,0);
  \draw[orange!80!black, line width=0.2pt, opacity=0.4] (3.485,0) to[out=58,in=122] (6.765,0);
  \draw[citewine, line width=0.2pt, opacity=0.4] (6.560,0) to[out=58,in=122] (6.970,0);
  \draw[projframe, line width=0.2pt, opacity=0.4] (5.535,0) to[out=58,in=122] (6.970,0);
  \draw[linknavy, line width=0.2pt, opacity=0.4] (3.895,0) to[out=58,in=122] (6.970,0);
  \draw[orange!80!black, line width=0.2pt, opacity=0.4] (3.690,0) to[out=58,in=122] (6.970,0);
  \draw[citewine, line width=0.2pt, opacity=0.4] (6.765,0) to[out=58,in=122] (7.175,0);
  \draw[projframe, line width=0.2pt, opacity=0.4] (5.740,0) to[out=58,in=122] (7.175,0);
  \draw[linknavy, line width=0.2pt, opacity=0.4] (4.100,0) to[out=58,in=122] (7.175,0);
  \draw[orange!80!black, line width=0.2pt, opacity=0.4] (3.895,0) to[out=58,in=122] (7.175,0);
  \draw[citewine, line width=0.2pt, opacity=0.4] (6.970,0) to[out=58,in=122] (7.380,0);
  \draw[projframe, line width=0.2pt, opacity=0.4] (5.945,0) to[out=58,in=122] (7.380,0);
  \draw[linknavy, line width=0.2pt, opacity=0.4] (4.305,0) to[out=58,in=122] (7.380,0);
  \draw[orange!80!black, line width=0.2pt, opacity=0.4] (4.100,0) to[out=58,in=122] (7.380,0);
  \draw[citewine, line width=0.2pt, opacity=0.4] (7.175,0) to[out=58,in=122] (7.585,0);
  \draw[projframe, line width=0.2pt, opacity=0.4] (6.150,0) to[out=58,in=122] (7.585,0);
  \draw[linknavy, line width=0.2pt, opacity=0.4] (4.510,0) to[out=58,in=122] (7.585,0);
  \draw[orange!80!black, line width=0.2pt, opacity=0.4] (4.305,0) to[out=58,in=122] (7.585,0);
  \draw[citewine, line width=0.2pt, opacity=0.4] (7.380,0) to[out=58,in=122] (7.790,0);
  \draw[projframe, line width=0.2pt, opacity=0.4] (6.355,0) to[out=58,in=122] (7.790,0);
  \draw[linknavy, line width=0.2pt, opacity=0.4] (4.715,0) to[out=58,in=122] (7.790,0);
  \draw[orange!80!black, line width=0.2pt, opacity=0.4] (4.510,0) to[out=58,in=122] (7.790,0);
  \draw[citewine, line width=0.2pt, opacity=0.4] (7.585,0) to[out=58,in=122] (7.995,0);
  \draw[projframe, line width=0.2pt, opacity=0.4] (6.560,0) to[out=58,in=122] (7.995,0);
  \draw[linknavy, line width=0.2pt, opacity=0.4] (4.920,0) to[out=58,in=122] (7.995,0);
  \draw[orange!80!black, line width=0.2pt, opacity=0.4] (4.715,0) to[out=58,in=122] (7.995,0);
  \draw[citewine, line width=0.2pt, opacity=0.4] (7.790,0) to[out=58,in=122] (8.200,0);
  \draw[projframe, line width=0.2pt, opacity=0.4] (6.765,0) to[out=58,in=122] (8.200,0);
  \draw[linknavy, line width=0.2pt, opacity=0.4] (5.125,0) to[out=58,in=122] (8.200,0);
  \draw[orange!80!black, line width=0.2pt, opacity=0.4] (4.920,0) to[out=58,in=122] (8.200,0);
  \draw[citewine, line width=0.2pt, opacity=0.4] (7.995,0) to[out=58,in=122] (8.405,0);
  \draw[projframe, line width=0.2pt, opacity=0.4] (6.970,0) to[out=58,in=122] (8.405,0);
  \draw[linknavy, line width=0.2pt, opacity=0.4] (5.330,0) to[out=58,in=122] (8.405,0);
  \draw[orange!80!black, line width=0.2pt, opacity=0.4] (5.125,0) to[out=58,in=122] (8.405,0);
  \draw[citewine, line width=0.2pt, opacity=0.4] (8.200,0) to[out=58,in=122] (8.610,0);
  \draw[projframe, line width=0.2pt, opacity=0.4] (7.175,0) to[out=58,in=122] (8.610,0);
  \draw[linknavy, line width=0.2pt, opacity=0.4] (5.535,0) to[out=58,in=122] (8.610,0);
  \draw[orange!80!black, line width=0.2pt, opacity=0.4] (5.330,0) to[out=58,in=122] (8.610,0);
  \draw[citewine, line width=0.2pt, opacity=0.4] (8.405,0) to[out=58,in=122] (8.815,0);
  \draw[projframe, line width=0.2pt, opacity=0.4] (7.380,0) to[out=58,in=122] (8.815,0);
  \draw[linknavy, line width=0.2pt, opacity=0.4] (5.740,0) to[out=58,in=122] (8.815,0);
  \draw[orange!80!black, line width=0.2pt, opacity=0.4] (5.535,0) to[out=58,in=122] (8.815,0);
  \draw[citewine, line width=0.2pt, opacity=0.4] (8.610,0) to[out=58,in=122] (9.020,0);
  \draw[projframe, line width=0.2pt, opacity=0.4] (7.585,0) to[out=58,in=122] (9.020,0);
  \draw[linknavy, line width=0.2pt, opacity=0.4] (5.945,0) to[out=58,in=122] (9.020,0);
  \draw[orange!80!black, line width=0.2pt, opacity=0.4] (5.740,0) to[out=58,in=122] (9.020,0);
  \draw[citewine, line width=0.2pt, opacity=0.4] (8.815,0) to[out=58,in=122] (9.225,0);
  \draw[projframe, line width=0.2pt, opacity=0.4] (7.790,0) to[out=58,in=122] (9.225,0);
  \draw[linknavy, line width=0.2pt, opacity=0.4] (6.150,0) to[out=58,in=122] (9.225,0);
  \draw[orange!80!black, line width=0.2pt, opacity=0.4] (5.945,0) to[out=58,in=122] (9.225,0);
  \draw[citewine, line width=0.2pt, opacity=0.4] (9.020,0) to[out=58,in=122] (9.430,0);
  \draw[projframe, line width=0.2pt, opacity=0.4] (7.995,0) to[out=58,in=122] (9.430,0);
  \draw[linknavy, line width=0.2pt, opacity=0.4] (6.355,0) to[out=58,in=122] (9.430,0);
  \draw[orange!80!black, line width=0.2pt, opacity=0.4] (6.150,0) to[out=58,in=122] (9.430,0);
  \draw[citewine, line width=0.2pt, opacity=0.4] (9.225,0) to[out=58,in=122] (9.635,0);
  \draw[projframe, line width=0.2pt, opacity=0.4] (8.200,0) to[out=58,in=122] (9.635,0);
  \draw[linknavy, line width=0.2pt, opacity=0.4] (6.560,0) to[out=58,in=122] (9.635,0);
  \draw[orange!80!black, line width=0.2pt, opacity=0.4] (6.355,0) to[out=58,in=122] (9.635,0);
  \draw[citewine, line width=0.2pt, opacity=0.4] (9.430,0) to[out=58,in=122] (9.840,0);
  \draw[projframe, line width=0.2pt, opacity=0.4] (8.405,0) to[out=58,in=122] (9.840,0);
  \draw[linknavy, line width=0.2pt, opacity=0.4] (6.765,0) to[out=58,in=122] (9.840,0);
  \draw[orange!80!black, line width=0.2pt, opacity=0.4] (6.560,0) to[out=58,in=122] (9.840,0);
  \draw[citewine, line width=0.2pt, opacity=0.4] (9.635,0) to[out=58,in=122] (10.045,0);
  \draw[projframe, line width=0.2pt, opacity=0.4] (8.610,0) to[out=58,in=122] (10.045,0);
  \draw[linknavy, line width=0.2pt, opacity=0.4] (6.970,0) to[out=58,in=122] (10.045,0);
  \draw[orange!80!black, line width=0.2pt, opacity=0.4] (6.765,0) to[out=58,in=122] (10.045,0);
  \draw[citewine, line width=0.2pt, opacity=0.4] (9.840,0) to[out=58,in=122] (10.250,0);
  \draw[projframe, line width=0.2pt, opacity=0.4] (8.815,0) to[out=58,in=122] (10.250,0);
  \draw[linknavy, line width=0.2pt, opacity=0.4] (7.175,0) to[out=58,in=122] (10.250,0);
  \draw[orange!80!black, line width=0.2pt, opacity=0.4] (6.970,0) to[out=58,in=122] (10.250,0);
  \draw[citewine, line width=0.2pt, opacity=0.4] (10.045,0) to[out=58,in=122] (10.455,0);
  \draw[projframe, line width=0.2pt, opacity=0.4] (9.020,0) to[out=58,in=122] (10.455,0);
  \draw[linknavy, line width=0.2pt, opacity=0.4] (7.380,0) to[out=58,in=122] (10.455,0);
  \draw[orange!80!black, line width=0.2pt, opacity=0.4] (7.175,0) to[out=58,in=122] (10.455,0);
  \draw[citewine, line width=0.2pt, opacity=0.4] (10.250,0) to[out=58,in=122] (10.660,0);
  \draw[projframe, line width=0.2pt, opacity=0.4] (9.225,0) to[out=58,in=122] (10.660,0);
  \draw[linknavy, line width=0.2pt, opacity=0.4] (7.585,0) to[out=58,in=122] (10.660,0);
  \draw[orange!80!black, line width=0.2pt, opacity=0.4] (7.380,0) to[out=58,in=122] (10.660,0);
  \draw[citewine, line width=0.2pt, opacity=0.4] (10.455,0) to[out=58,in=122] (10.865,0);
  \draw[projframe, line width=0.2pt, opacity=0.4] (9.430,0) to[out=58,in=122] (10.865,0);
  \draw[linknavy, line width=0.2pt, opacity=0.4] (7.790,0) to[out=58,in=122] (10.865,0);
  \draw[orange!80!black, line width=0.2pt, opacity=0.4] (7.585,0) to[out=58,in=122] (10.865,0);
  \draw[citewine, line width=0.2pt, opacity=0.4] (10.660,0) to[out=58,in=122] (11.070,0);
  \draw[projframe, line width=0.2pt, opacity=0.4] (9.635,0) to[out=58,in=122] (11.070,0);
  \draw[linknavy, line width=0.2pt, opacity=0.4] (7.995,0) to[out=58,in=122] (11.070,0);
  \draw[orange!80!black, line width=0.2pt, opacity=0.4] (7.790,0) to[out=58,in=122] (11.070,0);
  \draw[citewine, line width=0.2pt, opacity=0.4] (10.865,0) to[out=58,in=122] (11.275,0);
  \draw[projframe, line width=0.2pt, opacity=0.4] (9.840,0) to[out=58,in=122] (11.275,0);
  \draw[linknavy, line width=0.2pt, opacity=0.4] (8.200,0) to[out=58,in=122] (11.275,0);
  \draw[orange!80!black, line width=0.2pt, opacity=0.4] (7.995,0) to[out=58,in=122] (11.275,0);
  \draw[citewine, line width=0.2pt, opacity=0.4] (11.070,0) to[out=58,in=122] (11.480,0);
  \draw[projframe, line width=0.2pt, opacity=0.4] (10.045,0) to[out=58,in=122] (11.480,0);
  \draw[linknavy, line width=0.2pt, opacity=0.4] (8.405,0) to[out=58,in=122] (11.480,0);
  \draw[orange!80!black, line width=0.2pt, opacity=0.4] (8.200,0) to[out=58,in=122] (11.480,0);
  \draw[citewine, line width=0.2pt, opacity=0.4] (11.275,0) to[out=58,in=122] (11.685,0);
  \draw[projframe, line width=0.2pt, opacity=0.4] (10.250,0) to[out=58,in=122] (11.685,0);
  \draw[linknavy, line width=0.2pt, opacity=0.4] (8.610,0) to[out=58,in=122] (11.685,0);
  \draw[orange!80!black, line width=0.2pt, opacity=0.4] (8.405,0) to[out=58,in=122] (11.685,0);
  \draw[citewine, line width=0.2pt, opacity=0.4] (11.480,0) to[out=58,in=122] (11.890,0);
  \draw[projframe, line width=0.2pt, opacity=0.4] (10.455,0) to[out=58,in=122] (11.890,0);
  \draw[linknavy, line width=0.2pt, opacity=0.4] (8.815,0) to[out=58,in=122] (11.890,0);
  \draw[orange!80!black, line width=0.2pt, opacity=0.4] (8.610,0) to[out=58,in=122] (11.890,0);
  \draw[citewine, line width=0.2pt, opacity=0.4] (11.685,0) to[out=58,in=122] (12.095,0);
  \draw[projframe, line width=0.2pt, opacity=0.4] (10.660,0) to[out=58,in=122] (12.095,0);
  \draw[linknavy, line width=0.2pt, opacity=0.4] (9.020,0) to[out=58,in=122] (12.095,0);
  \draw[orange!80!black, line width=0.2pt, opacity=0.4] (8.815,0) to[out=58,in=122] (12.095,0);
  \draw[citewine, line width=0.2pt, opacity=0.4] (11.890,0) to[out=58,in=122] (12.300,0);
  \draw[projframe, line width=0.2pt, opacity=0.4] (10.865,0) to[out=58,in=122] (12.300,0);
  \draw[linknavy, line width=0.2pt, opacity=0.4] (9.225,0) to[out=58,in=122] (12.300,0);
  \draw[orange!80!black, line width=0.2pt, opacity=0.4] (9.020,0) to[out=58,in=122] (12.300,0);
  \draw[citewine, line width=0.2pt, opacity=0.4] (12.095,0) to[out=58,in=122] (12.505,0);
  \draw[projframe, line width=0.2pt, opacity=0.4] (11.070,0) to[out=58,in=122] (12.505,0);
  \draw[linknavy, line width=0.2pt, opacity=0.4] (9.430,0) to[out=58,in=122] (12.505,0);
  \draw[orange!80!black, line width=0.2pt, opacity=0.4] (9.225,0) to[out=58,in=122] (12.505,0);
  \draw[citewine, line width=0.2pt, opacity=0.4] (12.300,0) to[out=58,in=122] (12.710,0);
  \draw[projframe, line width=0.2pt, opacity=0.4] (11.275,0) to[out=58,in=122] (12.710,0);
  \draw[linknavy, line width=0.2pt, opacity=0.4] (9.635,0) to[out=58,in=122] (12.710,0);
  \draw[orange!80!black, line width=0.2pt, opacity=0.4] (9.430,0) to[out=58,in=122] (12.710,0);
  \draw[citewine, line width=0.2pt, opacity=0.4] (12.505,0) to[out=58,in=122] (12.915,0);
  \draw[projframe, line width=0.2pt, opacity=0.4] (11.480,0) to[out=58,in=122] (12.915,0);
  \draw[linknavy, line width=0.2pt, opacity=0.4] (9.840,0) to[out=58,in=122] (12.915,0);
  \draw[orange!80!black, line width=0.2pt, opacity=0.4] (9.635,0) to[out=58,in=122] (12.915,0);
  \fill[msgfill] (0.000,0) circle (1.5pt);
  \draw (0.000,0) circle (1.5pt);
  \fill[msgfill] (0.205,0) circle (1.5pt);
  \draw (0.205,0) circle (1.5pt);
  \fill[msgfill] (0.410,0) circle (1.5pt);
  \draw (0.410,0) circle (1.5pt);
  \fill[msgfill] (0.615,0) circle (1.5pt);
  \draw (0.615,0) circle (1.5pt);
  \fill[msgfill] (0.820,0) circle (1.5pt);
  \draw (0.820,0) circle (1.5pt);
  \fill[msgfill] (1.025,0) circle (1.5pt);
  \draw (1.025,0) circle (1.5pt);
  \fill[msgfill] (1.230,0) circle (1.5pt);
  \draw (1.230,0) circle (1.5pt);
  \fill[msgfill] (1.435,0) circle (1.5pt);
  \draw (1.435,0) circle (1.5pt);
  \fill[msgfill] (1.640,0) circle (1.5pt);
  \draw (1.640,0) circle (1.5pt);
  \fill[msgfill] (1.845,0) circle (1.5pt);
  \draw (1.845,0) circle (1.5pt);
  \fill[msgfill] (2.050,0) circle (1.5pt);
  \draw (2.050,0) circle (1.5pt);
  \fill[msgfill] (2.255,0) circle (1.5pt);
  \draw (2.255,0) circle (1.5pt);
  \fill[msgfill] (2.460,0) circle (1.5pt);
  \draw (2.460,0) circle (1.5pt);
  \fill[msgfill] (2.665,0) circle (1.5pt);
  \draw (2.665,0) circle (1.5pt);
  \fill[msgfill] (2.870,0) circle (1.5pt);
  \draw (2.870,0) circle (1.5pt);
  \fill[msgfill] (3.075,0) circle (1.5pt);
  \draw (3.075,0) circle (1.5pt);
  \fill[projframe!12] (3.280,0) circle (1.4pt);
  \fill[projframe!12] (3.485,0) circle (1.4pt);
  \fill[projframe!25] (3.690,0) circle (1.4pt);
  \fill[projframe!25] (3.895,0) circle (1.4pt);
  \fill[projframe!38] (4.100,0) circle (1.4pt);
  \fill[projframe!38] (4.305,0) circle (1.4pt);
  \fill[projframe!51] (4.510,0) circle (1.4pt);
  \fill[projframe!55] (4.715,0) circle (1.4pt);
  \fill[projframe!59] (4.920,0) circle (1.4pt);
  \fill[projframe!59] (5.125,0) circle (1.4pt);
  \fill[projframe!64] (5.330,0) circle (1.4pt);
  \fill[projframe!64] (5.535,0) circle (1.4pt);
  \fill[projframe!64] (5.740,0) circle (1.4pt);
  \fill[projframe!64] (5.945,0) circle (1.4pt);
  \fill[projframe!64] (6.150,0) circle (1.4pt);
  \fill[projframe!64] (6.355,0) circle (1.4pt);
  \fill[projframe!64] (6.560,0) circle (1.4pt);
  \fill[projframe!64] (6.765,0) circle (1.4pt);
  \fill[projframe!64] (6.970,0) circle (1.4pt);
  \fill[projframe!64] (7.175,0) circle (1.4pt);
  \fill[projframe!64] (7.380,0) circle (1.4pt);
  \fill[projframe!64] (7.585,0) circle (1.4pt);
  \fill[projframe!64] (7.790,0) circle (1.4pt);
  \fill[projframe!64] (7.995,0) circle (1.4pt);
  \fill[projframe!64] (8.200,0) circle (1.4pt);
  \fill[projframe!64] (8.405,0) circle (1.4pt);
  \fill[projframe!64] (8.610,0) circle (1.4pt);
  \fill[projframe!64] (8.815,0) circle (1.4pt);
  \fill[projframe!64] (9.020,0) circle (1.4pt);
  \fill[projframe!64] (9.225,0) circle (1.4pt);
  \fill[projframe!64] (9.430,0) circle (1.4pt);
  \fill[projframe!64] (9.635,0) circle (1.4pt);
  \fill[projframe!64] (9.840,0) circle (1.4pt);
  \fill[projframe!64] (10.045,0) circle (1.4pt);
  \fill[projframe!64] (10.250,0) circle (1.4pt);
  \fill[projframe!64] (10.455,0) circle (1.4pt);
  \fill[projframe!64] (10.660,0) circle (1.4pt);
  \fill[projframe!64] (10.865,0) circle (1.4pt);
  \fill[projframe!64] (11.070,0) circle (1.4pt);
  \fill[projframe!64] (11.275,0) circle (1.4pt);
  \fill[projframe!64] (11.480,0) circle (1.4pt);
  \fill[projframe!64] (11.685,0) circle (1.4pt);
  \fill[projframe!64] (11.890,0) circle (1.4pt);
  \fill[projframe!64] (12.095,0) circle (1.4pt);
  \fill[projframe!64] (12.300,0) circle (1.4pt);
  \fill[projframe!64] (12.505,0) circle (1.4pt);
  \fill[projframe!64] (12.710,0) circle (1.4pt);
  \fill[projframe!64] (12.915,0) circle (1.4pt);
  \draw[dashed, black!55] (5.330,-2.0) -- (5.330,-0.05);
  \node[black!70, anchor=west] at (5.410,-1.5) {\tiny $W_{26}$: all $16$ seeds reached};
  \draw[decorate, decoration={brace, mirror, amplitude=2.5pt}] (-0.03,-0.13) -- (3.075,-0.13) node[midway, below=1pt] {\tiny $16$ data seeds $W_0\dots W_{15}$};
  \draw[decorate, decoration={brace, mirror, amplitude=2.5pt}] (3.280,-0.13) -- (12.915,-0.13) node[midway, below=1pt] {\tiny $48$ derived $W_{16}\dots W_{63}$};
  \begin{scope}[yshift=-2.0cm]
    \draw[->, black!55] (-0.05,0) -- (13.120,0) node[right, black!70] {\tiny $t$};
    \draw[->, black!55] (-0.05,0) -- (-0.05,0.95) node[above, black!70, align=center] {\tiny seeds\\[-2pt]\tiny reached};
    \draw[black!30, dashed] (-0.05,0.8) -- (13.120,0.8) node[right, black!55] {\tiny $16$};
    \draw[projframe, thick] plot coordinates {(0.000,0.050) (0.205,0.050) (0.410,0.050) (0.615,0.050) (0.820,0.050) (1.025,0.050) (1.230,0.050) (1.435,0.050) (1.640,0.050) (1.845,0.050) (2.050,0.050) (2.255,0.050) (2.460,0.050) (2.665,0.050) (2.870,0.050) (3.075,0.050) (3.280,0.200) (3.485,0.200) (3.690,0.350) (3.895,0.350) (4.100,0.500) (4.305,0.500) (4.510,0.650) (4.715,0.700) (4.920,0.750) (5.125,0.750) (5.330,0.800) (5.535,0.800) (5.740,0.800) (5.945,0.800) (6.150,0.800) (6.355,0.800) (6.560,0.800) (6.765,0.800) (6.970,0.800) (7.175,0.800) (7.380,0.800) (7.585,0.800) (7.790,0.800) (7.995,0.800) (8.200,0.800) (8.405,0.800) (8.610,0.800) (8.815,0.800) (9.020,0.800) (9.225,0.800) (9.430,0.800) (9.635,0.800) (9.840,0.800) (10.045,0.800) (10.250,0.800) (10.455,0.800) (10.660,0.800) (10.865,0.800) (11.070,0.800) (11.275,0.800) (11.480,0.800) (11.685,0.800) (11.890,0.800) (12.095,0.800) (12.300,0.800) (12.505,0.800) (12.710,0.800) (12.915,0.800)};
  \end{scope}
\end{tikzpicture}
}
  \caption{The message schedule as a computed dependency graph
  (generated by \texttt{tools/wschedule\_graph.py}, which also emits a
  Graphviz \texttt{.dot} version). Nodes are the $64$ schedule words in
  index order; the sixteen data-carrying seeds $W_0,\dots,W_{15}$ are
  blue, and derived words darken with the number of seeds they
  transitively depend on. Each of the four arc colours is one term of
  $W_t=\sigma_1(W_{t-2})\shaplus W_{t-7}\shaplus\sigma_0(W_{t-15})\shaplus
  W_{t-16}$. The lower trace counts seeds reached versus $t$; it reaches
  all sixteen at $W_{26}$ (dashed line), after which the message is fully
  diffused through the schedule alone.}
  \label{fig:schedulegraph}
\end{figure}

One further structural reduction deserves emphasis, because it shrinks
the apparent dimensionality of the mystery. Only two of the eight
registers ever receive fresh values: each round computes one new word
entering the $a$-position and one entering the $e$-position, while the
other six merely shift along. The eight-word state is a \emph{delay
line}---a memoization of recent history, not eight independent degrees
of freedom. Writing $A_t$ and $E_t$ for the words produced at round $t$
(so that the state at that moment reads
$(A_t, A_{t-1}, A_{t-2}, A_{t-3},\, E_t, E_{t-1}, E_{t-2}, E_{t-3})$),
the sixty-four rounds are exactly the coupled two-track recurrence
\begin{align*}
  T_t &= E_{t-4} \shaplus \Sigma_1(E_{t-1})
        \shaplus \mathrm{Ch}(E_{t-1}, E_{t-2}, E_{t-3})
        \shaplus K_t \shaplus W_t,\\
  A_t &= T_t \shaplus \Sigma_0(A_{t-1})
        \shaplus \mathrm{Maj}(A_{t-1}, A_{t-2}, A_{t-3}),
  \qquad
  E_t \;=\; A_{t-4} \shaplus T_t,
\end{align*}
with the eight initial values $A_{-1}, \dots, A_{-4}, E_{-1}, \dots,
E_{-4}$ read off the incoming chaining state. So the ``$8$-dimensional''
iteration is really a pair of scalar recurrences over $\Ztwo$ with
lookback four---formally a nonlinear cousin of the lagged-Fibonacci and
linear-feedback recurrences whose theory is complete
(\cref{sec:prngs})---and the cryptanalytic literature does in fact work
in exactly this $A/E$ form. Nothing is lost in the reduction; the honest
shape of a round is two new words under an eight-word window of memory.

Notice what is \emph{absent}: no multiplications, no S-box tables, no
number-theoretic structure---only the three ARX primitives and two
quadratic Boolean functions, iterated $64$ times. The function is
degree-two nonlinearity compounded through $64$ rounds of geometric
whiplash between $\Ftwo^{256}$ and $(\Ztwo)^8$. That so spare a recipe
apparently produces a completely structureless map is the central mystery
this paper wants to sell.

\subsection{Two constructions that come with proofs}\label{sec:constructions}

It is worth pausing to locate the mystery precisely, because \SHA{} is
built in three nested layers and---remarkably---the two \emph{outer}
layers each carry a solid theorem. Only the innermost object is
theorem-free. Naming the layers from the outside in: the
\emph{mode} (Merkle--Damg{\aa}rd iteration, \cref{sec:merkledamgard})
wraps the \emph{compression function} $f$; the compression function is
itself built by the \emph{Davies--Meyer} construction around a block
cipher $E$; and $E$---the cipher SHACAL-2, which is just the round
machinery above run as a keyed permutation---is the finite map about
which nothing is proved. The provable security of \SHA{} is entirely a
statement about the two wrappers, conditional on the core behaving well.

\paragraph{The mode: Merkle--Damg{\aa}rd.} The outer layer's theorem is
already in hand: \cref{thm:md} says that a collision in the full hash
constructively yields a collision in $f$, so \emph{collision resistance
of $f$ is inherited by the whole function}---unconditionally, in the
standard model, needing only that the final block encode the message
length (the ``strengthening'' the proof consumes). This is as solid as
theorems in this subject get: no idealization, no unproven assumption, a
one-page reduction. What it does \emph{not} give is equally important and
sharply delimited---it transfers collision resistance and nothing else,
which is exactly why length extension and the multicollision family of
\cref{sec:mdstructure} slip through. The mode is a theorem about one
property, and an honest one.

\paragraph{The compression function: Davies--Meyer.} The middle layer is
the feedforward met in \cref{sec:roundshape}. Take a block cipher
$E_m(h)$---plaintext the chaining value $h\in\State$, key the message
block $m$---and set
\[
  f(h,m) \;=\; E_m(h) \,\shaplus\, h .
\]
The feedforward ``$\shaplus\,h$'' is the whole point. Without it,
$f = E_m$ would be a \emph{permutation} of $\State$ for each fixed $m$,
hence trivially invertible (decrypt with key $m$) and useless as a
one-way compression---the same observation that, read dynamically,
\cref{sec:unless} shows is what makes iterated \SHA{} lose entropy at
all. Adding the input back breaks invertibility, and does so in a way one
can prove is optimal, provided the cipher is idealized:

\begin{theorem}[Preneel--Govaerts--Vandewalle; Black--Rogaway--Shrimpton]
\label{thm:dm}
Model $E$ as an \emph{ideal cipher} (for each key an independent, uniformly
random permutation, accessible only as a black box in both directions).
Then the Davies--Meyer compression $f(h,m)=E_m(h)\shaplus h$ is optimally
collision- and preimage-resistant: any adversary asking $q$ cipher
queries finds a collision with probability at most $q(q+1)/2^{n}$, and a
preimage with probability of order $q/2^{n}$---matching the birthday
bound $2^{n/2}$ for collisions and the brute-force bound $2^{n}$ for
preimages, and both are tight.
\end{theorem}

\citet{PGV1993} isolated this by classifying all $64$ natural ways to
build a rate-one compression function from a block cipher and singling
out the $12$ sound ones (Davies--Meyer among them); \citet{BRS2002} then
supplied the rigorous black-box proofs and the exact bounds quoted above.
So the middle layer, too, has a theorem---but note the price on the
ticket. It holds in the \emph{ideal cipher model}, the block-cipher
cousin of the random-oracle model of \cref{sec:models}, and is therefore
a heuristic about \SHA{}, not a proof about it: SHACAL-2 is a fixed, public,
emphatically non-ideal cipher whose ``key schedule'' is literally the
message expansion of \cref{sec:roundshape}, so the assumption that
distinct message blocks act like independent random permutations is
exactly the leap of faith the model licenses and cannot justify. A
concrete reminder that the idealization is lossy: Davies--Meyer has
\emph{trivially computable fixed points}---solving $E_m(h)\shaplus h = h$
means $E_m(h)=0$, i.e.\ $h = E_m^{-1}(0)$, one short computation per
block---a harmless-looking feature that became a lever for the
long-message second-preimage attacks of \cref{sec:mdstructure}.

\paragraph{The moral, localized.} Two layers, two theorems: the mode
proved unconditionally, the compression proved in an idealized model, and
between them they reduce the security of the whole edifice to a single
demand---that the core map $f$ (equivalently the cipher $E$) have no
exploitable structure. That demand is \cref{sec:frameworks}'s gap, and
it is where every remaining page of this paper lives. The architecture is
solved mathematics; the mystery has been squeezed into one finite object,
and the constructions of this section are precisely the machinery that
did the squeezing.

\subsection{One honest theorem: invertible diffusion}\label{sec:rivest}

Amid so much that is unproved, it is worth exhibiting in full one of the
few genuine theorems about \SHA{}'s internals---partly because it is
lovely, partly because its proof lives in exactly the finite geometry of
\cref{sec:statespace}, and partly because it shows the ``middle
territory'' of \cref{sec:gap} is not empty by necessity.

Ask a basic anatomical question about the stirring maps of
\cref{sec:roundshape}: are $\Sigma_0, \Sigma_1, \sigma_0, \sigma_1$
\emph{bijections} of the word space $\Ftwo^{32}$? (They are
$\Ftwo$-linear, hence ``nonlinearities'' only through the other lens of
\cref{sec:statespace}---as maps of $\Ztwo$ they are wildly nonlinear.
The question is whether they discard information.) Rivest answered this
for the pure-rotation case in a four-page note~\citep{Rivest2011}:

\begin{theorem}[Rivest]\label{thm:rivest}
Let $w$ be a power of two. The exclusive-or of $t$ distinct rotations of
a $w$-bit word,
$x \mapsto \mathrm{ROTR}^{r_1}(x) \shaxor \cdots \shaxor \mathrm{ROTR}^{r_t}(x)$,
is invertible if and only if $t$ is odd.
\end{theorem}

\begin{proof}[Proof idea]
Identify $\Ftwo^{32}$ with the quotient ring
$R = \Ftwo[x]/(x^{32}-1)$, a bit vector becoming a polynomial and
rotation becoming multiplication by $x^{r}$. The map in question is
multiplication by $p(x) = x^{r_1} + \cdots + x^{r_t}$, invertible in $R$
exactly when $\gcd\bigl(p(x),\, x^{32}-1\bigr) = 1$. Over $\Ftwo$ the
modulus collapses---$x^{32} - 1 = (x-1)^{32}$, because $32$ is a power of
two and squaring is a ring homomorphism in characteristic $2$---so the
gcd is $1$ precisely when $(x-1) \nmid p$, i.e.\ when $p(1) = t \bmod 2$
is $1$.
\end{proof}

So $\Sigma_0 = \mathrm{ROTR}^{2} \shaxor \mathrm{ROTR}^{13} \shaxor
\mathrm{ROTR}^{22}$ and $\Sigma_1$, with their three rotations each, are
bijections: they stir without spilling. The theorem repays contemplation.
It is two lines of algebra, yet it turns on the arithmetic accident that
the word length is a power of two, which makes the cyclic algebra $R$
\emph{local} (one maximal ideal, generated by $x-1$): the only way an
xor-of-rotations can fail to be invertible is by having an even number of
terms. Had the designers used $33$-bit words, $x^{33}-1$ would factor
richly and the parity rule would be false. Small finite geometries have
loud opinions.

Because the proof is constructive, invertibility need not be left as an
existence statement: the inverses can be exhibited outright. Carrying out
the division in $R$---equivalently, inverting the corresponding
$32\times32$ circulant matrices over $\Ftwo$; \cref{app:sigmacode} gives
a complete reference implementation, which also confirms the results on
thousands of random words---produces closed forms that are again xors of
rotations:
\begin{equation}\label{eq:sigmainv}
\begin{aligned}
  \Sigma_0^{-1} \;&=\; \bigoplus_{r \in S_0} \mathrm{ROTR}^{r},
  & S_0 &= \{0,1,2,5,7,8,9,10,12,16,17,19,22,23,25,29,30\},\\
  \Sigma_1^{-1} \;&=\; \bigoplus_{r \in S_1} \mathrm{ROTR}^{r},
  & S_1 &= \{2,3,5,7,9,11,12,13,18,21,22,23,24,26,27,29,30\},
\end{aligned}
\end{equation}
each an xor of \emph{seventeen} rotations ($S_0$ includes $r = 0$, the
identity rotation). Seventeen is odd, as \cref{thm:rivest} demands of any
invertible xor-of-rotations---and as it must be on general grounds:
evaluating $p\,q \equiv 1$ at $x = 1$ forces
$|p| \cdot |q| \equiv 1 \pmod 2$, so odd-weight elements have odd-weight
inverses. The asymmetry of sparseness is the striking part: three
rotations to stir, seventeen to unstir---more than half of the
thirty-two possible offsets. Invertible is not the same as symmetrically
cheap, even inside the linear layer.

The ring picture yields more than the inverses; it computes the whole
multiplicative life of these maps, essentially by hand. Composition of
rotations adds offsets modulo $32$, so write an xor-of-rotations as a
polynomial with the \emph{offsets as exponents}: $\Sigma_0$ becomes
$p_0 = x^{2} + x^{13} + x^{22}$. Squaring is the Frobenius map of
characteristic $2$---it simply doubles every offset---so squaring $p_0$
three times gives, after reducing exponents modulo $32$,
$\Sigma_0^{8} = \mathrm{ROTR}^{8}$: \emph{eight applications of the
three-rotation stir collapse to one bare rotation}. Hence
$\Sigma_0^{32} = \mathrm{id}$, and the order of $\Sigma_0$ in
$\mathrm{GL}_{32}(\Ftwo)$ is exactly $32$; the same computation gives
$\Sigma_1^{16} = \mathrm{id}$ with order exactly $16$. In particular
$\Sigma_0^{-1} = \Sigma_0^{31}$ and $\Sigma_1^{-1} = \Sigma_1^{15}$, so
\cref{eq:sigmainv} could have been generated, offsets and all, by
repeated squaring. One more turn of the crank and the observation
becomes a theorem: for \emph{any} odd xor-of-rotations $u$ of $32$-bit
words, $u^{32} = u(x)^{32} = u(x^{32}) = u(1) = 1$, so every invertible
xor-of-rotations has order dividing $32$. The stirring maps of \SHA{},
for all their visual violence, are elements of small $2$-power order in
a $2$-group---one more exact structural fact living quietly next door to
the conjectured structurelessness (\cref{sec:hope}).

Note what the theorem does \emph{not} cover: $\sigma_0 =
\mathrm{ROTR}^{7} \shaxor \mathrm{ROTR}^{18} \shaxor \mathrm{SHR}^{3}$
and $\sigma_1$ each replace one rotation by a bit-discarding shift,
leaving the polynomial ring for plain linear algebra; and the genuinely
non-invertible ingredients of \SHA{} are elsewhere entirely---the
bitwise-quadratic $\mathrm{Ch}$ and $\mathrm{Maj}$ and the Davies--Meyer
feedforward, where all the designed information loss is concentrated. The
picture that emerges is architectural: \emph{lossless} stirring layers
(provably so, by \cref{thm:rivest}), \emph{lossy} mixing concentrated in
three auditable places. That so clean a structural statement was worth a
research note in 2011---twenty years into the design lineage---says
something about how sparsely theorized this territory is.

\begin{project}{The linear algebra of diffusion}{linear algebra;
  programming}{3--4 weeks}
  Write each of $\Sigma_0, \Sigma_1, \sigma_0, \sigma_1$ as an explicit
  $32 \times 32$ matrix over $\Ftwo$ and decide invertibility of the
  $\sigma$'s (settling what \cref{thm:rivest} leaves open). Compute the
  order of each invertible map in $\mathrm{GL}_{32}(\Ftwo)$, the cycle
  structure of its action on $\Ftwo^{32}$, and the subgroup the four maps
  generate. Then reprove \cref{thm:rivest} for general word length $w$:
  the answer is a pretty gcd condition (worked out in Rivest's note), and
  comparing $w = 32$ with $w = 33$ makes the power-of-two accident vivid.
\end{project}

\begin{project}{Deriving the exact inverses of $\Sigma_0$ and $\Sigma_1$}
  {polynomial arithmetic over $\Ftwo$; programming}{2--3 weeks}
  \Cref{thm:rivest} guarantees that $\Sigma_0$ and $\Sigma_1$ are
  bijections, and \cref{eq:sigmainv} displays their inverses; this
  project derives that display from scratch, by two independent methods,
  and decides what \cref{eq:sigmainv} deliberately leaves out. The steps
  are fully determined:
  \begin{enumerate}
    \item \textbf{Move to the ring.} Identify a word with a polynomial in
      $R = \Ftwo[x]/(x^{32}-1)$, bit $i$ becoming the coefficient of
      $x^i$, so that a rotation is multiplication by a power of $x$.
      Write $\Sigma_0$ as multiplication by
      $p_0(x) = x^{2} + x^{13} + x^{22}$ and $\Sigma_1$ by
      $p_1(x) = x^{6} + x^{11} + x^{25}$ (the \SHA{} rotation offsets).
    \item \textbf{Invert by linear algebra.} Multiplication by $p$ is an
      $\Ftwo$-linear circulant map; build its $32\times32$ matrix (columns
      are $p, xp, x^2p, \dots$) and invert it over $\Ftwo$ by Gaussian
      elimination. The last column of the inverse, read as a polynomial
      $q$, satisfies $p\,q \equiv 1 \pmod{x^{32}-1}$.
    \item \textbf{Read off the answer.} You should recover exactly the
      seventeen-term rotation sets of \cref{eq:sigmainv}. Verify
      $p\,q = 1$ in $R$ directly, and check your derivation against the
      reference implementation of \cref{app:sigmacode}. Then reproduce
      \cref{eq:sigmainv} a third way, with no linear algebra at all: by
      repeated squaring, using $\Sigma_0^{-1} = \Sigma_0^{31}$ and
      $\Sigma_1^{-1} = \Sigma_1^{15}$.
    \item \textbf{Cross-check by Euclid.} Recompute $q$ by the extended
      Euclidean algorithm on $p(x)$ and $x^{32}-1 = (x-1)^{32}$ in
      $\Ftwo[x]$; the two methods must agree.
    \item \textbf{Contrast with the $\sigma$'s.} Repeat for
      $\sigma_0 = \mathrm{ROTR}^{7}\shaxor\mathrm{ROTR}^{18}\shaxor
      \mathrm{SHR}^{3}$: the shift $\mathrm{SHR}^{3}$ leaves the cyclic
      ring (it is not a unit multiple of $x$), so the polynomial trick
      fails and only the matrix method applies. Decide invertibility of
      $\sigma_0, \sigma_1$ this way---settling what \cref{thm:rivest}
      does not cover.
  \end{enumerate}
  Deliverable: the two explicit inverse maps, a proof they are correct,
  and a one-paragraph account of why the $\sigma$'s need the heavier tool.
\end{project}

\subsection{Nothing up my sleeve}\label{sec:nums}

One question a mathematician should ask of \cref{sec:roundshape}: where do
the constants come from? The initial state $\mathrm{IV}$ and the sixty-four
round constants $K_t$ could, in principle, hide a trapdoor---constants
engineered so their designer knows structure nobody else can see. The
history is not hypothetical: the unexplained S-box constants of the 1970s
DES cipher fed decades of suspicion (later resolved in the designers'
favor), and post-Snowden, an actual back-doored NIST standard (the
Dual\_EC random generator) was withdrawn. The SHA family itself supplies
the sharpest demonstration that its designer's choices are informed
rather than arbitrary: the NSA's unexplained one-bit revision of SHA-0
into SHA-1 (see the footnote in \cref{sec:history}) anticipated by four
years the differential attack later found in the open
literature~\citep{ChabaudJoux1998}. The SHA-2 designers therefore
committed to constants they visibly could not have chosen:
\[
  \mathrm{IV}_j = \text{first 32 fractional bits of } \sqrt{p_j},
  \qquad
  K_t = \text{first 32 fractional bits of } \sqrt[3]{p_t},
\]
where $p_1 < p_2 < \cdots$ are the primes $2, 3, 5, \dots$~\citep{FIPS180-4}.
(So $\mathrm{IV}_1 = \mathtt{6a09e667}$ encodes $\sqrt{2}-1$ in binary.)
Such choices are called \emph{nothing-up-my-sleeve} constants: points of
the state space selected by a rule so canonical that no freedom remains in
which to hide intent.

There is a pleasing irony here that deserves a mathematician's smile: the
constants are trusted \emph{because} they are digits of irrational
numbers, yet essentially nothing is proved about those digits either---the
normality of $\sqrt{2}$ is an open problem. The construction does not
replace an unproved assumption with a proved one; it replaces an
\emph{unauditable} choice with an \emph{auditable} one. Epistemology, not
theorems, is doing the work---a theme by now familiar
(\cref{sec:epistemology}).

\subsection{Reasons for hope: exact solutions next door}\label{sec:hope}

Almost everything after this section will be statistical:
random-mapping silhouettes, test batteries, mixing arguments---the
approximate machinery one reaches for when exact structure is
unavailable. Before conceding that this is the only possible register,
it is worth assembling the evidence that exact structure keeps turning
up \emph{next door}, because the history of discrete iteration is
littered with problems that looked hopelessly procedural until someone
found the right coordinates.

The Josephus process (\cref{sec:josephus}) is this paper's opening
specimen: an $(n-1)$-step elimination dynamics collapsed to a single
bit rotation. \Cref{thm:rivest} is a second, living inside \SHA{}
itself: a question about its diffusion maps settled by two lines of
algebra in a local ring. The genre is much larger. Additive cellular
automata---iterative, visually turbulent, every bit as ``chaotic'' in
appearance as a hash round---were solved outright by
\citet{MartinOdlyzkoWolfram1984}: represent configurations as
polynomials over $\Ftwo$ and time evolution becomes multiplication,
with global transient and cycle structure falling out as exact
divisibility statements. \citet{Sutner1989} dispatched the Lights-Out
puzzle the same way (the $\sigma$-game as $\Ftwo$-linear algebra). The
pattern is older than computing: the logistic map at parameter $4$, the
textbook ``chaotic'' system, has the closed form
$x_t = \sin^2\!\bigl(2^t \arcsin\sqrt{x_0}\,\bigr)$---an observation
going back to Ulam and von Neumann, and the simplest instance of
Schr{\"o}der's program of solving iteration by
conjugation~\citep{Schroeder1871}: find $\psi$ making
$\psi \circ f \circ \psi^{-1}$ linear, and the $t$-th iterate is
computed in one stroke. That program grew into modern complex
dynamics~\citep{Milnor2006}. ``Chaotic'' and ``exactly solvable'' are
not antonyms; they are frequently the same object described before and
after the right change of variables.

Within the hash family the lesson recurs. SHA-1's message schedule is
$\Ftwo$-linear, hence \emph{exactly} analyzable---its difference
propagation is linear algebra---and that handle is precisely what the
collision attacks of \cref{sec:falling} gripped. The professional
response to such objects has been the cryptanalyst's panoply,
differential attacks in the tradition of \citet{BihamShamir1991} and
their descendants: a \emph{statistical} theory of near-structure,
probability-weighted trails threaded through the rounds
(\cref{sec:differential} is devoted to it). The examples
above argue for keeping a different question alive: whether some change
of coordinates---a direct finite-geometry path, in the spirit of the
polynomial representations that solved cellular automata, with no
mixing theory, no ergodicity, and nothing approximate---renders some
nontrivial slice of the iteration exactly computable. The designers'
explicit goal was to foreclose this (\cref{sec:statespace}), and the
smart money says they succeeded. But the foreclosure is a conjecture,
not a theorem; the components keep falling one at a time to exact
methods (\cref{thm:rivest}, the T-function
theory~\citep{KlimovShamir2003, Anashin2006}, the two-track reduction
of \cref{sec:roundshape}); and the hunting grounds are tiny finite
objects, where undergraduate computation is a genuine research
instrument (\cref{sec:projects}).

% ------------------------------------------------------------
% §4  Random mappings and Flajolet's lens
% ------------------------------------------------------------
\section{Random mappings: Flajolet's telescope}\label{sec:randommaps}

If the working conjecture is ``\SHA{} behaves like a random function,''
then the right response is the scientist's: \emph{a conjecture that
mentions randomness is a bundle of quantitative predictions}. What does a
random function actually look like? Remarkably, this question has an
exact, beautiful answer, developed by Harris in the probabilistic
era~\citep{Harris1960} and perfected by Flajolet and Odlyzko with
generating-function methods~\citep{FlajoletOdlyzko1990}. The answer turns
the vague conjecture into a table of measurable constants.

\subsection{The functional graph of a random map}\label{sec:fungraph}

Let $\varphi$ be a function from a finite set of size $N$ to itself---for
us, eventually, $\varphi = $ (some restriction of) \SHA{}, with
$N = 2^n$. Draw the \emph{functional graph}: a node for each point, an
arrow $x \to \varphi(x)$. Because every node has out-degree exactly one,
the shape is forced: each connected component consists of exactly one
directed cycle, with trees hanging off the cycle nodes
(\cref{fig:rho}). Iterating $\varphi$ from any seed $x_0$ walks a
$\rho$-shaped path: a \emph{tail} of $\lambda$ new points, then a
\emph{cycle} of length $\mu$ repeated forever.

\begin{figure}[ht]
  \centering
  \begin{tikzpicture}[>={Stealth[length=1.8mm]},
      pt/.style={circle, fill, inner sep=1.4pt},
      lbl/.style={font=\scriptsize}]
    % cycle of 8, label inside
    \foreach \i in {0,...,7} {
      \pgfmathsetmacro{\ang}{\i*45}
      \node[pt] (c\i) at ($(4.55,0)+(\ang:1.2)$) {};
    }
    \foreach \i [evaluate=\i as \j using {int(mod(\i+1,8))}] in {0,...,7}
      \draw[->] (c\i) to[bend right=18] (c\j);
    \node[lbl, align=center] at (4.55,0) {cycle\\length $\mu$};
    % horizontal tail entering the leftmost cycle node c4
    \node[pt] (t0) at (0.4,0) {};
    \node[pt] (t1) at (1.4,0) {};
    \node[pt] (t2) at (2.4,0) {};
    \node[lbl, left=1mm of t0] {$x_0$};
    \draw[->] (t0) -- (t1);
    \draw[->] (t1) -- (t2);
    \draw[->] (t2) -- (c4);
    % small trees hanging off cycle nodes, kept clear of everything
    \node[pt] (a1) at ($(c2)+(-0.45,0.75)$) {};
    \node[pt] (a2) at ($(c2)+(0.5,0.8)$) {};
    \node[pt] (a3) at ($(a2)+(0.55,0.65)$) {};
    \draw[->] (a1) -- (c2); \draw[->] (a2) -- (c2); \draw[->] (a3) -- (a2);
    \node[pt] (b1) at ($(c7)+(0.65,-0.6)$) {};
    \node[pt] (b2) at ($(b1)+(0.55,-0.6)$) {};
    \draw[->] (b1) -- (c7); \draw[->] (b2) -- (b1);
    % tail brace
    \draw[decorate, decoration={brace, amplitude=4pt, mirror}]
      (0.4,-0.35) -- (3.1,-0.35)
      node[midway, below=5pt, lbl] {tail, length $\lambda$};
  \end{tikzpicture}
  \caption{One component of a functional graph. Iteration from any seed
  traces a $\rho$: a tail of length $\lambda$ into a cycle of length
  $\mu$. Trees hang off the cycle; components with their unique cycles
  tile the whole space.}
  \label{fig:rho}
\end{figure}

Now let $\varphi$ be chosen \emph{uniformly at random} among all $N^N$
self-maps. Flajolet and Odlyzko computed the expected value of essentially
every natural parameter of this picture; a sample, asymptotically as
$N \to \infty$~\citep{FlajoletOdlyzko1990}:

\begin{center}
\begin{tabular}{@{}llll@{}}
  \toprule
  parameter && expectation & at $N = 2^{256}$ \\
  \midrule
  tail length $\lambda$ (random seed)   && $\sqrt{\pi N / 8}$        & $\approx 2^{127.3}$ \\
  cycle length $\mu$ (random seed)      && $\sqrt{\pi N / 8}$        & $\approx 2^{127.3}$ \\
  rho length $\lambda + \mu$            && $\sqrt{\pi N / 2}$        & $\approx 2^{128.3}$ \\
  number of cyclic points               && $\sqrt{\pi N / 2}$        & $\approx 2^{128.3}$ \\
  number of components                  && $\tfrac12 \ln N$          & $\approx 89$ \\
  image size $|\varphi(\State)|$        && $(1 - e^{-1})\,N$         & $\approx 0.632\,N$ \\
  points with no preimage (leaves)      && $e^{-1}\, N$              & $\approx 0.368\,N$ \\
  largest component                     && $\approx 0.7582\,N$       & --- \\
  largest tree                          && $\approx 0.48\,N$         & --- \\
  \bottomrule
\end{tabular}
\end{center}

\noindent
Pause on the flavor of these results. A ``typical'' random map is nothing
like a random permutation: it is lossy (a third of the space is never hit
at all), it has one giant component swallowing three quarters of
everything, its cycles are tiny compared to $N$ (order $\sqrt{N}$), and
the number of components grows only logarithmically. This is a sharp,
peculiar, recognizable statistical silhouette---a fingerprint of
structurelessness itself.

\subsection{Where the constants come from}\label{sec:singularity}

The route to that table is pure Flajolet, and worth a paragraph even in an
overview, because it connects our discrete object to complex analysis. A
tree rooted at a node whose arrow enters a cycle is a labelled rooted
tree; its exponential generating function $T(z)$ satisfies the functional
equation
\[
  T(z) = z\, e^{T(z)}
\]
(a root, plus a set of subtrees). Components are cycles of trees, and
mappings are sets of components, which by the symbolic calculus gives the
mapping EGF
\[
  M(z) \;=\; \frac{1}{1 - T(z)} .
\]
All the parameters in the table are obtained by marking features in these
equations and then reading off coefficient asymptotics. The engine for the
last step: $T$ has a square-root singularity at $z = e^{-1}$, where
$T(e^{-1}) = 1$---so $M$ blows up there---and \emph{singularity analysis}
transfers the local behavior of a generating function at its dominant
singularity directly into asymptotics of its coefficients. The
half-integer exponents in the singular expansions are what produce all
those $\sqrt{N}$'s; the method is developed at book length
in~\citet{FlajoletSedgewick2009}, with random mappings as a recurring
worked example. For a mathematician, this is the pleasant surprise of the
whole subject: questions about iterating a hash function route through
the local behavior of an analytic function at a branch point.

\subsection{The telescope, pointed at \SHA{}}\label{sec:telescope}

Here is the epistemological payoff. The random-function conjecture of
\cref{sec:history} predicts that if we build from \SHA{} a self-map of an
$N$-point space---most simply
\[
  \varphi_n : \bitstrings^n \to \bitstrings^n, \qquad
  \varphi_n(x) = \text{first } n \text{ bits of } \SHA(x),
\]
for moderate $n$---then the functional graph of $\varphi_n$, an object
\SHA{}'s designers never contemplated and certainly never tuned, should
match every line of the Flajolet--Odlyzko table, not merely in expectation
but in distribution (the limiting distributions are also known). Each
parameter is a falsifiable prediction; jointly they are a
\emph{coordinate system} in which the position of \SHA{} among
pseudorandom objects can be measured rather than argued about. For
contrast, structured maps land visibly elsewhere in these coordinates:
$x \mapsto x \shaplus c$ decomposes into cycles only (no trees: it is a
bijection), $x \mapsto x^2 \bmod p$ has its graph dictated by quadratic
reciprocity, and a one-round toy \SHA{} should sit measurably off the
random silhouette. Where exactly it rejoins the silhouette as rounds are
added is, as far as we know, unmeasured---and eminently measurable
(\cref{sec:projects}).

At full scale ($n = 256$, $N = 2^{256}$) nobody will ever traverse a rho:
the expected rho length is about $2^{128}$ steps. The table is then the
\emph{only} available description of the global geometry of the object the
whole world is using---a telescope pointed at a structure we can predict
in detail but never visit. And the predictions are not idle: they are the
engineering basis of real algorithms. Pollard's rho
method~\citep{Pollard1978} finds collisions and discrete logarithms
precisely by betting that its iteration behaves like a random map (cycle
detection in $O(\sqrt{N})$ steps with $O(1)$ memory); van Oorschot and
Wiener's parallel collision search~\citep{vanOorschotWiener1999} turned
that bet into an industrial discipline (distinguished points, linear
parallel speedup), and the SHAttered attack's budget of $2^{63}$
(\cref{sec:falling}) was costed with exactly this calculus. Every time
such an algorithm's observed running time matches its predicted constant,
a Flajolet statistic of a real hash function has been confirmed to one
more decimal place.

So far the telescope has resolved a \emph{single} application of the
function. Because \SHA{} is so often composed with itself---in password
stretching, hash chains, and Bitcoin's double hash---the dynamics of
\emph{iterating} $\varphi$ form their own coordinate system, and
\cref{sec:iterated} takes it up next: how the image shrinks, where
orbits settle, and the precise sense in which iteration almost, but not
quite, induces a random permutation.

\begin{project}{A functional-graph atlas of truncated \SHA{}}
  {programming; a first probability course}{6--8 weeks (flagship)}
  For $n = 16, 18, \dots, 30$: exhaustively build the functional graph of
  $\varphi_n$, measure every parameter in the table above (and the
  distributions, not just means), and compare with the
  Flajolet--Odlyzko asymptotics \emph{and} with genuinely random maps of
  the same size as a control. Deliverable: an ``atlas'' of plots---theory
  curve, random-map control cloud, \SHA{} data points. Then repeat with
  reduced-round \SHA{} ($1, 2, 4, \dots, 64$ rounds) and watch the data
  walk onto the random silhouette as rounds accumulate. To our knowledge
  no published systematic atlas of this kind exists; undergraduates can
  own it.
\end{project}

\begin{project}{Collision hunting, the van Oorschot--Wiener way}
  {programming; probability}{4--6 weeks}
  Implement Pollard rho with Brent or Floyd cycle detection on
  $\varphi_n$ for $n$ up to ${\sim}48$, then the van Oorschot--Wiener
  distinguished-points parallel method on a lab's worth of machines. Measure the distribution of time-to-collision against the
  theory of \cref{sec:fungraph}. Students reproduce, at toy scale, the
  exact methodology of the SHAttered attack---and learn why $256$ bits is
  chosen so that $\sqrt{N}$ is out of humanity's reach.
\end{project}

\subsection{Coda: counting to three, and the dials of \SHA{}}
\label{sec:josephus3}

Before putting the telescope away, we settle a debt incurred in the
opening paragraph of this paper (\cref{sec:josephus}). The Josephus
circle there counted to \emph{two}, and the survivor map collapsed to a
one-bit rotation. What happens when the circle counts to three? The
answer is known, and it comes---pleasingly---from the same Odlyzko whose
lens this section has been borrowing~\citep{OdlyzkoWilf1991}. The
survivor $J_3(n)$ is delivered by a crisp algorithm~\citep[\S 3.3]{GKP1994}:
iterate $D_0 = 1$, $D_k = \bigl\lceil \tfrac{3}{2} D_{k-1} \bigr\rceil$,
stop at the first $D_k > 2n$, and take $J_3(n) = 3n + 1 - D_k$. Odlyzko
and Wilf proved that these iterates admit an exact closed form,
\[
  D_k \;=\; \Bigl\lfloor K \bigl(\tfrac{3}{2}\bigr)^{k} \Bigr\rfloor
  \quad (k = 0, 1, \dots), \qquad
  K \;=\; 1.62227\,05028\,84767\ldots,
\]
so $J_3$, like $J_2$, ``has a formula.'' But inspect the formula. The
constant $K$ is assembled from the error terms of the very recurrence it
is supposed to shortcut; no independent characterization of it is known,
so the only way anyone can evaluate the closed form is to run the
iteration it was meant to abolish---as the authors say plainly, the
explicit formula ``is not an improvement over the algorithm.'' Nor is
this an artifact of technique. For the general iteration
$f_{n+1} = \lceil \alpha f_n \rceil$ the analogous constant $c(\alpha)$
is a jagged function of $\alpha$, with jump discontinuities at every
integer and at every Josephus point $\alpha = q/(q-1)$; and the paper
closes with a conjecture on the fine distribution of the error terms
that its authors liken to proving that a given real number is normal in
a given base---``likely to be very difficult.'' Inside this family the
solved case is a degenerate point: at $q = 2$ the error terms vanish
identically, $K(2) = 1$, and the machinery collapses back to the bit
rotation of \cref{eq:josephus}.

Dwell on what just happened. We did not change the problem; we turned
one dial by one click, from every-second to every-third, and
analyzability fell off a cliff---from a closed form a child can verify
to a constant nobody can compute independently and a conjecture
calibrated against normality. The collapse that opened this paper was
not a property of decimation in general. It was an accident of the
parameter $q = 2$. In the family, wildness is generic; solvability is
the exception.

\SHA{} deserves to be examined in exactly this light, because its
definition is also a single point in an evident parameter family, with
many more dials: the rotation offsets $(2,13,22)$ and $(6,11,25)$ of
$\Sigma_0$ and $\Sigma_1$, the rotation-and-shift offsets $(7,18;\,3)$
and $(17,19;\,10)$ of $\sigma_0$ and $\sigma_1$, the word length $32$,
the round count $64$, the schedule taps $(2,7,15,16)$, the sixty-four
additive constants $K_t$. Each could have been set otherwise, and asking
what happens as one is perturbed---the analogue of moving $q$---is a
legitimate and largely unexplored experimental axis. In one direction
the family provably degenerates: set the three rotation offsets of
$\Sigma_0$ equal and the map collapses to a single rotation, as
transparent as the Josephus shift; drop to an even number of rotation
summands and invertibility itself dies (\cref{thm:rivest}); move the
word length off a power of two and the parity law is repealed. The
published constants are believed to sit in the opposite, generic regime,
where nothing collapses---but between the degenerate points and the
conjecturally generic bulk, nobody has mapped the terrain. The
functional-graph atlas proposed above thus gains a second axis: measure
the Flajolet coordinates of reduced-round variants not only as rounds
accumulate but as the rotation and shift amounts are turned away from
the standard's values, and ask whether the chosen point is
distinguishable, in any coordinate of the table, from a randomly chosen
neighbor.

% ------------------------------------------------------------
% §5  Iterating the function: the dynamics of SHA composed with itself
% ------------------------------------------------------------
\section{Iterating the function}\label{sec:iterated}

\Cref{sec:randommaps} pointed Flajolet's telescope at a single
application of \SHA{}. But the function is very often applied to its own
output---%
\[
  \varphi(x) \;=\; \SHA(x), \qquad
  \varphi^{k} \;=\; \underbrace{\SHA \circ \SHA \circ \cdots \circ
  \SHA}_{k},
\]
read as a self-map of the $256$-bit strings (feed the digest back as a
one-block message). This is no idle variation. Password stretching runs
$\varphi^{k}$ for $k$ in the hundreds of thousands; hash chains for
one-time passwords and hash-based signatures are iteration made into a
protocol; Bitcoin's core primitive is literally $\varphi^{2}$, the
double hash $\SHA(\SHA(\cdot))$. And iteration is its own coordinate
system for the random-function conjecture: the \emph{dynamics} of
$\varphi$---how its image shrinks, where its orbits settle, whether it
hides fixed points---are a fresh battery of falsifiable predictions,
sharper in some ways than the static silhouette of \cref{sec:randommaps}
because they probe the function's interaction with \emph{itself}. This
section reads those predictions off the random-map model and asks, in
each case, what an anomaly would mean. The user's instinct is the right
place to start: one might imagine iteration eventually shuffling the
state space like a random permutation---and seeing \emph{exactly why
that almost happens, and exactly what stops it}, turns out to expose the
single design decision that makes \SHA{} one-way.

\subsection{Iteration leaks: the shrinking image}\label{sec:shrinking}

A permutation loses nothing: applied repeatedly, it merely relabels. A
random-looking map is different, and the difference compounds. Recall
from \cref{sec:fungraph} that one application of a random map covers only
$(1-e^{-1})N \approx 0.63\,N$ of the space; the leaves---the $37\%$ of
points with no preimage---are gone after one step and never return.
Apply $\varphi$ again and the survivors thin again, by the same logic
applied to the image. Writing $\beta_{k}$ for the expected fraction of
the space still reachable after $k$ steps, the recurrence is exactly
\[
  \beta_{0} = 1, \qquad \beta_{k+1} = 1 - e^{-\beta_{k}},
\]
(``a point survives to step $k+1$ iff at least one of its
$\varphi$-preimages survived to step $k$''), and its solution decays
with a clean asymptotic law~\citep{FlajoletOdlyzko1990}:
\[
  \beta_{k} \;\sim\; \frac{2}{k} \qquad (k \to \infty).
\]
Iteration is a funnel (\cref{fig:funnel}). After $k$ rounds the reachable
image has shrunk to about $2N/k$ points, so the process has quietly
discarded roughly $\log_{2}(k/2)$ bits of range. The number is small but
not zero and the moral is exact: \emph{composing a hash with itself can
only lose entropy, never create it}. For a password-stretching count of
$k = 10^{6}$, the naive iterate $\varphi^{k}$ confines its outputs to a
$2^{-19}$ fraction of the digest space---nineteen bits of the advertised
$256$ silently surrendered to the funnel.

\begin{figure}[ht]
  \centering
  \begin{tikzpicture}[font=\small, >={Stealth[length=1.6mm]}]
    % funnel: nested horizontal bars shrinking downward
    \foreach \y/\w/\lab in {0/4.4/{$\State$: all $N=2^{256}$},
        -0.72/2.78/{$\varphi(\State)\approx 0.63\,N$},
        -1.34/2.06/{$\varphi^{2}(\State)\approx 0.47\,N$},
        -1.86/1.30/{$\varphi^{4}(\State)\approx 0.37\,N$}} {
      \draw[fill=projfill, draw=projframe]
        (-\w/2,\y) rectangle (\w/2,\y-0.42);
    }
    \node[font=\scriptsize] at (0,-0.21) {$N$};
    \node[right, font=\scriptsize] at (2.3,-0.21) {all strings};
    \node[right, font=\scriptsize] at (2.3,-0.93) {leaves gone (37\%)};
    \node[right, font=\scriptsize] at (2.3,-2.07)
      {$\approx 2N/k$ after $k$ steps};
    % dots
    \node at (0,-2.42) {$\vdots$};
    % recurrent set at the bottom
    \draw[fill=lenfill, draw=orange!70!black]
      (-0.62,-3.2) rectangle (0.62,-2.85);
    \node[font=\scriptsize] at (0,-3.02) {$R$};
    \node[right, font=\scriptsize, align=left] at (2.3,-3.02)
      {recurrent set $R=\bigcap_{k}\varphi^{k}(\State)$:\\
       the cycles, $\approx\!\sqrt{\pi N/2}=2^{128.3}$ points};
    \draw[->] (0,-2.5) -- (0,-2.83);
    % downward funnel guide lines
    \draw[gray, dashed] (-2.2,0) -- (-0.62,-2.85);
    \draw[gray, dashed] (2.2,0) -- (0.62,-2.85);
    \node[left, font=\scriptsize, align=right] at (-2.3,-1.5)
      {\emph{entropy}\\ \emph{funnel}\\ $\beta_{k}\!\sim\!2/k$};
  \end{tikzpicture}
  \caption{Iterating a random-looking map is a one-way funnel. Each
  application discards the current leaves, so the reachable image shrinks
  as $\beta_{k}\sim 2/k$ toward the \emph{recurrent set} $R$---the union
  of all cycles, about $\sqrt{\pi N/2}\approx 2^{128.3}$ points. On $R$,
  and only there, $\varphi$ acts as a permutation (\cref{sec:permutation}).
  The funnel is why iterated hashing loses (a little) entropy, and why
  real key-stretching functions are \emph{not} pure iteration.}
  \label{fig:funnel}
\end{figure}

Two consequences deserve flagging. First, the prediction is
\emph{measurable}: the shrinkage law $\beta_{k}\sim 2/k$ is a specific
curve that truncated \SHA{} must trace if it is random-like, and a
reduced-round or structured map need not (\cref{sec:projects}). Second,
the practical fix is illuminating precisely because it \emph{breaks} the
composition. Standardized password-based key derivation does not compute
$\varphi^{k}$; it re-injects the password and a salt at every step,
iterating a map like $x \mapsto \SHA(x \shaxor \text{password})$ that
changes each round~\citep{NIST_SP800132}. That is no longer a single
function composed with itself---it is a \emph{driven} iteration, formally
the random walk of \cref{sec:prngs} rather than a functional graph---and
the salt-dependent driving is exactly what defeats the funnel (and the
precomputation attack of \cref{sec:itervuln}). The theory tells the
engineer why the naive design leaks and the standard one does not.

\subsection{The permutation at the bottom of the funnel}
\label{sec:permutation}

Where does the funnel end? Not at nothing: the image cannot shrink below
the points that lie on cycles, since a cyclic point is its own distant
iterate and survives forever. These \emph{recurrent} points---the set
$R=\bigcap_{k\ge 0}\varphi^{k}(\State)$, the union of all the cycles of
\cref{fig:rho}---are the funnel's floor, and here the user's intuition
lands exactly. \emph{On $R$, iterated \SHA{} is a genuine permutation.}
This is not a probabilistic near-miss; it is a theorem about every
finite self-map: restricted to its recurrent set a function is a
bijection, because each cyclic point has exactly one cyclic preimage
(the point before it on its cycle). \SHA{} permutes its cycles, and
$\varphi$ carried to the limit is a permutation of the roughly
$\sqrt{\pi N/2}\approx 2^{128.3}$ recurrent points---conjecturally a
\emph{uniformly random} permutation of a random such subset, which is
the strongest form of the random-map hypothesis.

So the imagined ``random permutation induced by iteration'' is real---%
but confined to an unreachable floor. The recurrent set is a
random-looking subset of size $2^{128.3}$, with no short description and
no way to enumerate it short of traversing orbits of length $2^{128}$;
one can prove it hosts a permutation without ever exhibiting a single
one of its elements. This is the same epistemic shape as the rest of the
paper (\cref{sec:telescope}): a structure we can characterize in full
and visit never.

Now the sharp contrast that answers ``why astronomically unlikely?''
Being a permutation on the \emph{recurrent set} is automatic; being a
permutation on the \emph{whole} space $\State$ is the rarest thing a map
can be. The fraction of self-maps that are permutations is
\[
  \frac{N!}{N^{N}} \;\approx\; \sqrt{2\pi N}\; e^{-N},
\]
so a uniformly random map is a permutation with probability of order
$e^{-N}$: not astronomically but \emph{doubly}-exponentially small, about
$2^{-1.44\cdot 2^{256}}$. Equivalently, injectivity anywhere is
expensive---a random map already has about $0.37 N$ points sharing their
images with others. Hence if \SHA{} were ever found to be injective on
its full domain, or even to cover meaningfully more than the predicted
$0.63N$ on the first step, that would not be a lucky coincidence: it
would be a gross structural signal, a deviation of the exact kind
\cref{sec:gap} is built to detect. The permutation lives at the bottom
of the funnel by necessity; a permutation at the \emph{top} would be a
five-alarm anomaly.

\subsection{``Unless\dots'': the one addition that breaks the bijection}
\label{sec:unless}

Why is \SHA{} not that anomaly---why does it funnel at all, rather than
permute? The answer is a single line of its construction, and it is worth
seeing because it identifies one-wayness with lossiness as the \emph{same
design act}. Recall from \cref{sec:roundshape} that the compression
function is a block cipher run in Davies--Meyer mode:
\[
  f(h, m) \;=\; E_{m}(h) \,\shaplus\, h,
\]
where $E_{m}$---the core of \SHA{}, essentially the cipher SHACAL-2---is,
for each fixed message block $m$, a \emph{permutation} of the state
(every round is invertible, so the whole cipher is a bijection one could
run backward). If $f$ were just $E_{m}$, iteration would be the iteration
of a permutation: reversible, lossless, its functional graph all cycles
and no trees, injective on the whole space---precisely the
doubly-exponentially unlikely object of \cref{sec:permutation}. The
feedforward ``$\shaplus\,h$'' is what destroys this
(\cref{fig:feedforward}). Adding the input back turns a bijection into a
generically many-to-one map---the sum of a permutation and the identity
need not be injective---and it is this one modular addition that makes
$f$ one-way, non-invertible, and lossy all at once. The tree-and-funnel
geometry of \cref{sec:fungraph,sec:shrinking}, the $37\%$ of leaves, the
$2/k$ shrinkage, the very existence of preimage-resistance: all of it is
manufactured by the single ``$\shaplus\,h$'' that the designers added
exactly so the rounds could not be run backward.

\begin{figure}[ht]
  \centering
  \begin{tikzpicture}[font=\small, >={Stealth[length=1.7mm]},
      pt/.style={circle, fill, inner sep=1.3pt}]
    % LEFT: E_m a permutation (bijection) — parallel arrows, no merging
    \begin{scope}
    \node[font=\scriptsize] at (0.9,1.35) {$E_{m}$ alone: bijection};
    \foreach \i in {0,1,2,3} {
      \node[pt] (l\i) at (0,-\i*0.5) {};
      \node[pt] (r\i) at (1.8,-\i*0.5) {};
    }
    \draw[->] (l0) -- (r1); \draw[->] (l1) -- (r3);
    \draw[->] (l2) -- (r0); \draw[->] (l3) -- (r2);
    \node[font=\scriptsize, align=center] at (0.9,-2.15)
      {every target hit once\\ \emph{reversible, lossless}};
    \end{scope}
    % RIGHT: f = E_m(h) + h — collisions and leaves appear
    \begin{scope}[xshift=5.4cm]
    \node[font=\scriptsize] at (0.9,1.35)
      {$f=E_{m}(h)\shaplus h$: collapses};
    \foreach \i in {0,1,2,3} {
      \node[pt] (a\i) at (0,-\i*0.5) {};
      \node[pt] (b\i) at (1.8,-\i*0.5) {};
    }
    \draw[->] (a0) -- (b1); \draw[->] (a1) -- (b1); % collision into b1
    \draw[->] (a2) -- (b3); \draw[->] (a3) -- (b3); % collision into b3
    \node[orange!80!black, font=\scriptsize, right] at (1.95,-0.5) {2 preimages};
    \node[orange!80!black, font=\scriptsize, right] at (1.95,-1.5) {2 preimages};
    \node[citewine, font=\scriptsize, right] at (1.95,0) {leaf};
    \node[citewine, font=\scriptsize, right] at (1.95,-1.0) {leaf};
    \node[font=\scriptsize, align=center] at (0.9,-2.15)
      {targets missed \& doubled\\ \emph{one-way, lossy}};
    \end{scope}
    \draw[->, thick] (2.55,-0.75) -- (3.35,-0.75)
      node[midway, above, font=\scriptsize] {$+\,h$};
  \end{tikzpicture}
  \caption{The ``unless.'' The cipher $E_{m}$ at the heart of \SHA{} is a
  permutation (left): iterate it and nothing is lost. The Davies--Meyer
  feedforward $f=E_{m}(h)\shaplus h$ (right) adds the input back, turning
  the bijection into a many-to-one map with leaves and collisions---the
  source of every tree, every lost bit, and all one-wayness. Iterated
  \SHA{} would be a random permutation on its whole domain
  (\cref{sec:permutation}) were it not for this one modular addition.}
  \label{fig:feedforward}
\end{figure}

This is the resolution of the opening puzzle. Iterated \SHA{} does induce
a random permutation---but on its recurrent floor, by necessity, not on
the whole space, which the feedforward forbids. The ``unless'' is
literal: \emph{unless one deletes the ``$\shaplus\,h$,''} in which case
\SHA{} degenerates into its underlying block cipher, iteration becomes
reversible, and the funnel disappears along with the security. The
lossiness the user intuited as a near-impossibility is, seen correctly,
the deliberate and load-bearing content of the design.

\subsection{Vulnerabilities specific to iteration}\label{sec:itervuln}

Iteration is not merely a lens; it is an attack surface and a building
block, and three concrete phenomena round out the picture.

\paragraph{Hash chains, and why they are only as strong as one step.}
Applying $\varphi$ repeatedly and revealing the chain \emph{backward} is
a classic authentication primitive: Lamport's one-time
passwords~\citep{Lamport1981} store $\varphi^{n}(s)$ on the server and
release $\varphi^{n-1}(s), \varphi^{n-2}(s), \dots$ one login at a time,
each value verified by one forward application and unforgeable ahead of
time by one-wayness. The same skeleton---iterate a one-way step, reveal
in reverse---underlies hash-based signatures (\cref{sec:uses}). What the
functional-graph view contributes is a caution: a chain is exactly as
strong as the one-wayness of a \emph{single} step, and if a short cycle
or a cheap partial inversion existed for the toy scales of
\cref{sec:projects}, the whole chain would unravel at once. Iteration
concentrates risk onto the per-step preimage problem.

\paragraph{Time--memory trade-offs: iteration as the attacker's
weapon.} The most striking vulnerability is that iterated hashing is the
data structure an adversary uses to \emph{invert} a one-way function by
precomputation. Hellman's time--memory trade-off~\citep{Hellman1980}
covers the domain with long chains $x, \varphi(x), \varphi^{2}(x),
\dots$, stores only the endpoints, and thereby inverts $\varphi$ on a
fresh target in online time $T$ and memory $M$ obeying
\[
  T\,M^{2} \;\approx\; N^{2},
\]
so that a one-time precomputation buys balanced online cost
$T = M = N^{2/3}$---far below the $N$ of blind brute force. Oechslin's
\emph{rainbow tables}~\citep{Oechslin2003} sharpen the constants with
per-column reduction functions and made the technique an everyday
password-cracking tool. At full \SHA{} scale $N^{2/3}=2^{170.7}$ is
safely unreachable, but against low-entropy inputs (short passwords) the
trade-off is devastating---and its defeat is once again the
\emph{driving} of \cref{sec:shrinking}: a per-user salt makes every
target's chain independent, so no shared precomputation
applies~\citep{NIST_SP800132}. The same iteration that stretches a key
for the defender builds the table for the attacker; salt is what breaks
the symmetry.

\paragraph{Fixed points and short cycles: cheap, sharp predictions.}
Finally, iteration poses questions with crisp random-model answers that
are eminently testable at toy scale. How many fixed points---strings
with $\SHA(x)=x$---should exist? For a random map the count is
$\mathrm{Poisson}(1)$: expected value exactly $1$, with probability
$e^{-1}\approx 0.37$ of \emph{none}. So \SHA{} very probably has a
handful (likely one) genuine $256$-bit fixed point that no one will ever
exhibit, a pleasant cousin of the ``does $\SHA(x)=x$ have a solution?''
puzzle. More generally the expected number of cycles of length $\ell$ is
$1/\ell$, a harmonic profile summing to the $\sqrt{\pi N/2}$ recurrent
points of \cref{sec:permutation}. Each of these is a falsifiable
prediction about a structureless map; a truncated \SHA{} that showed too
many fixed points, or cycles of anomalous length, would be exhibiting
structure. That nobody has looked, at model-organism scale, is the
recurring opportunity of this prospectus.

\begin{project}{The dynamics of iteration}{programming; a first
  probability course}{6--8 weeks (flagship)}
  Extend the functional-graph atlas of \cref{sec:telescope} from static
  statistics to \emph{dynamics}. For $\varphi_{n}(x) = $ first $n$ bits
  of $\SHA(x)$, $n = 16, \dots, 30$: (i) measure the image-shrinkage
  curve $\beta_{k}=|\varphi_{n}^{k}(\State)|/N$ and fit it against the
  $\beta_{k+1}=1-e^{-\beta_{k}}$ law and its $2/k$
  tail~\citep{FlajoletOdlyzko1990}; (ii) count fixed points and short
  cycles and compare with the $\mathrm{Poisson}(1)$ and $1/\ell$
  predictions; (iii) locate the recurrent set and confirm $\varphi$
  permutes it. Then repeat with reduced-round \SHA{} and with the
  feedforward \emph{removed} (pure $E_{m}$, a permutation) as a control,
  watching the funnel switch off. As a capstone, implement a small
  Hellman table~\citep{Hellman1980} on $\varphi_{n}$ and measure the
  $TM^{2}=N^{2}$ trade-off directly, with and without salt. Every
  quantity is exactly computable for $n \le 30$, and the ``remove the
  feedforward'' control makes the mechanism of \cref{sec:unless} visible
  as data.
\end{project}

% ------------------------------------------------------------
% §5  Frameworks for analysis, and the gap
% ------------------------------------------------------------
\section{Frameworks, and the size of the gap}\label{sec:frameworks}

We now assemble the existing theoretical apparatus around \SHA{} and take
honest stock of how far it reaches. The summary in advance: there are
mature frameworks on \emph{both} sides of the object---idealized models
above it, empirical and cryptanalytic measurements below it---and a
canyon in the middle where the function itself lives.

\subsection{What ``pseudorandom'' officially means}
\label{sec:definitions}

The foundational definitions of pseudorandomness are about
\emph{families}, not single objects. A pseudorandom function family in
the sense of Goldreich, Goldwasser, and Micali~\citep{GGM1986} is a keyed
collection $\{F_k\}$ such that no efficient algorithm, allowed to query a
black box, can tell a randomly keyed $F_k$ from a truly random function
with nonnegligible advantage. This definition is a triumph---it supports
composition theorems, equivalences, a whole complexity theory of
randomness---but notice that it \emph{cannot even be stated} for the
fixed, keyless, public function \SHA{}. For any fixed function, a trivial
\emph{distinguisher} exists. A distinguisher is the formal adversary in
these definitions: an efficient algorithm that queries a black box
containing either the function under test or a truly random function,
emits a one-bit verdict, and is scored by its \emph{advantage}, the gap
between its acceptance probabilities in the two worlds. The notion
descends from Yao's theory of computational
indistinguishability~\citep{Yao1982} and receives its textbook treatment
in Chapter~3 of~\citet{Goldreich2001}. In our context the trivial
distinguisher simply hard-codes one known value: query the box at
$\mathsf{0}$ and accept if the answer is the published
$\SHA(\mathsf{0})$. Against \SHA{} it accepts always; against a random
function, with probability $2^{-256}$; its advantage is essentially $1$,
and no keyless public function escapes it. The same
foundational awkwardness afflicts collision resistance: since collisions
exist by pigeonhole, ``no efficient algorithm finds one'' is technically
false---some two-line program prints a collision; humanity just cannot
exhibit it. Rogaway's ``human ignorance'' formalism~\citep{Rogaway2006}
repairs the semantics by making theorems constructive (explicit
reductions transforming any breaker of a protocol into an explicit
collision-finder), and it is the honest logical home of every security
claim about \SHA{}. But note what has happened: \emph{the definition of
success has been moved from mathematics into history}---``no one has
written down a collision, so far.''

\subsection{Idealized models: reasoning above the object}
\label{sec:models}

Protocol designers escape upward. In the random-oracle
model~\citep{BellareRogaway1993} one simply \emph{postulates} a public
uniformly random function and proves protocols secure relative to it; the
leap back to \SHA{} is heuristic, frankly acknowledged, and empirically
successful. The model even has an internal quality-control theory:
\emph{indifferentiability}~\citep{MRH2004} asks whether a construction
built from an ideal component is safely interchangeable with an ideal
whole, and by this standard plain Merkle--Damg{\aa}rd \emph{fails}---the
length-extension property (from $\SHA(m)$ anyone computes
$\SHA(m \Vert \mathrm{pad} \Vert m')$ without knowing $m$) is a genuine,
exploitable deviation from random-oracle behavior, identified precisely
and repaired by variants with security proofs~\citep{CDMP2005}. This is
the theory at its best: it cannot certify $f$ itself, but given an
idealized $f$ it classifies \emph{constructions} completely. The
architecture above the compression function is, in a real sense, solved
mathematics---the two-layer picture of \cref{sec:constructions}, where
the Merkle--Damg{\aa}rd mode (\cref{thm:md}) and the Davies--Meyer
compression (\cref{thm:dm}) each carry a theorem---and all residual
mystery is compressed into the finite map $f$.

There is a sharper way to say that the model is a heuristic, and it is a
theorem rather than a caveat. Canetti, Goldreich, and Halevi
constructed signature and encryption schemes that are \emph{provably
secure} in the random-oracle model and \emph{provably insecure} the
moment the oracle is replaced by \emph{any} concrete, publicly specified
function whatsoever, \SHA{} or otherwise~\citep{CGH2004}. The
construction is diagonal---the scheme tests whether the code it is given
is a short program and betrays its secret if so, which a true random
oracle never is but every real hash is---so it is contrived, and no
natural protocol is known to suffer. But it settles the logic: a
random-oracle proof is not a proof about \SHA{} with an admitted gap; it
is a proof about an idealized object that \emph{no} function can fully
become. The heuristic is not merely unproven, it is in general false,
and its continued empirical success (\cref{sec:below}) is therefore
itself part of the natural history---a regularity of the world, not a
lemma.

\subsection{The construction is not the oracle: provable structure
above \texorpdfstring{$f$}{f}}\label{sec:mdstructure}

Here is a genuine surprise, and it cuts against the drumbeat of this
paper. We keep asserting that \SHA{} exhibits \emph{no} observable
deviation from a random function. That is true of the compression
function $f$; it is \emph{false} of the Merkle--Damg{\aa}rd wrapper
around it. The iteration mode has provable, nameable, exploitable
structure---deviations from random-oracle behavior that hold
\emph{even if $f$ is a perfect random oracle}---and they are among the
most elegant results in the subject. They locate the paper's mystery
more precisely than we have so far: the residue of the unknown lives in
$f$, because the part outside $f$ is not unknown at all. Its
non-randomness is a set of theorems.

The keystone is Joux's multicollision attack~\citep{Joux2004}
(\cref{fig:multicoll}). For a true $256$-bit random function, finding
$2^k$ messages that all share one digest should cost about
$2^{256(2^k-1)/2^k}$ work---for eight-way agreement, near $2^{224}$,
astronomically beyond a single collision's $2^{128}$. In an iterated
hash it costs only $k$ times a single collision. The trick is pure
construction-level reasoning: birthday-find a colliding \emph{block
pair} at the first chaining value, then \emph{from the merged state}
find another colliding block pair, and so on for $k$ stages; freely
mixing the two choices at each stage yields $2^k$ full messages all
landing on the same final digest, at total cost $k \cdot 2^{128}$
instead of $2^{128 \cdot (2^k-1)/2^k}$. Cascading is the immediate
casualty: hashing with $H_1 \Vert H_2$, the ``belt and braces'' habit
of concatenating two hashes for safety, is no stronger than the better
single one---build a $2^{n_1/2}$-multicollision in $H_1$, then by
pigeonhole two of its exponentially many members collide in $H_2$ as
well. A folklore engineering practice was refuted by one structural
idea, no cryptanalysis of any $f$ required.

\begin{figure}[ht]
  \centering
  \begin{tikzpicture}[font=\small, >={Stealth[length=1.7mm]},
      cv/.style={circle, draw, minimum size=6.5mm, inner sep=0pt,
                 fill=projfill, font=\scriptsize},
      lbl/.style={font=\scriptsize}]
    \node[cv] (h0) at (0,0) {$h_0$};
    \node[cv] (h1) at (2.6,0) {$h_1$};
    \node[cv] (h2) at (5.2,0) {$h_2$};
    \node[cv] (h3) at (7.8,0) {$h_3$};
    \node[lbl] at (9.1,0) {$\cdots$};
    % each stage: two parallel colliding blocks
    \foreach \a/\b/\i in {h0/h1/1, h1/h2/2, h2/h3/3} {
      \draw[->] (\a) to[bend left=42]
        node[above, lbl] {$b_{\i}$} (\b);
      \draw[->] (\a) to[bend right=42]
        node[below, lbl] {$b_{\i}'$} (\b);
    }
    \node[lbl, align=center, projframe] at (3.9,-1.55)
      {each stage: one birthday collision $f(h,b_i)=f(h,b_i')$, cost
       $2^{128}$};
    \node[lbl, align=center] at (3.9,1.7)
      {choose $b_i$ \emph{or} $b_i'$ freely at every stage
       $\Rightarrow$ $2^{k}$ messages, one digest};
    \draw[decorate, decoration={brace, amplitude=4pt}]
      (0,2.05) -- (7.8,2.05);
    \node[lbl] at (10.4,0) {total $k\cdot 2^{128}$};
  \end{tikzpicture}
  \caption{Joux multicollisions~\citep{Joux2004}. Because the state is
  a narrow \emph{delay line} passed from block to block
  (\cref{sec:roundshape}), a chain of $k$ independent single collisions
  multiplies into $2^k$ colliding messages at only $k$ times the cost
  of one---a provable, non-statistical deviation from random-oracle
  behavior that lives entirely in the iteration, not in $f$. The
  narrow internal state is the culprit; a wide-pipe or sponge design
  (the SHA-3 sponge of \cref{sec:merkledamgard}) closes exactly this
  door.}
  \label{fig:multicoll}
\end{figure}

Two further theorems mine the same seam, and both convert the
\emph{narrowness of the internal state}---the eight-word delay line of
\cref{sec:roundshape}, no wider than the output---into an attack.
Kelsey and Schneier showed that a \emph{second preimage} for a long
message is far cheaper than the ideal $2^{256}$: given a target message
of $2^\ell$ blocks, ``expandable messages'' (a multicollision between
message lengths) let one splice in a forged prefix that rejoins the
victim's chain at some interior state, at cost about
$2^{256-\ell}$~\citep{KelseySchneier2005}. Long messages are
\emph{provably} weaker than short ones---a sentence with no meaning at
all for a random oracle, where every input is independent.
Kelsey and Kohno pushed the genre from collision to \emph{commitment
fraud}: their ``herding'' attack precomputes a branching ``diamond''
structure of chaining values, publishes a single digest as a
soothsayer's sealed prophecy, and then, given any later-revealed
prefix, herds it into the committed digest by appending a short
suffix~\citep{KelseyKohno2006}. The hash's advertised binding
property---the entire basis of the commitment use of
\cref{sec:uses}---is quantitatively weakened by construction-level
structure alone.

The moral reframes the whole paper, so we state it flatly. Every one of
these attacks treats $f$ as an ideal random oracle and still breaks the
advertised behavior of \SHA{}-the-construction; each is a
\emph{theorem}, not a measurement, and each traces to one visible
design fact---the internal state is no wider than the digest, so the
chaining value is a bottleneck that a random function would not have
(\cref{sec:roundshape}). This is why the SHA-3 competition selected a
sponge with a \emph{wide} internal
state~\citep{BDPV2007} that structurally forecloses all three attacks,
and why the
honest slogan is not ``\SHA{} behaves like a random oracle'' but the
sharper \emph{``$f$ behaves like a random function, and the wrapper's
deviations from a random oracle are completely known.''} The mystery of
this paper is now pinned exactly where it belongs: not in the
architecture, which is theorems all the way down, but in the single
finite map $f$, about which there are none.

\subsection{Empirical and adversarial measurement: attacking from below}
\label{sec:below}

From below, two very different measurement cultures apply.

\emph{Statistical testing} treats the function as a black-box PRNG and
runs batteries of hypothesis tests: NIST's SP\,800-22
suite~\citep{NIST2010} and the far more demanding TestU01
``BigCrush''~\citep{LEcuyerSimard2007}. \SHA{}-based generators pass
everything, which is informative but bounded: these batteries famously
demolish the classical linear generators (a Mersenne Twister fails
BigCrush's linear-complexity tests; lattice structure betrays linear
congruential generators---the pathologies catalogued since
Knuth~\citep{Knuth1997}), so passing places \SHA{} strictly above the
traditional PRNG canon; but a battery is a fixed, finite set of null
hypotheses, and structure invisible to all of them may be gross. Knuth's
own moral stands: statistical testing can refute, never confirm. For
\SHA{} the concrete record is easily stated: the counter sequence
$\SHA(1), \SHA(2), \SHA(3), \dots$ is not merely testable but
\emph{standardized}---NIST's Hash\_DRBG is essentially this
generator~\citep{SP80090A}---and the same construction over SHA-1
passed every one of BigCrush's tests in the paper that introduced the
battery. Passes of this kind are so thoroughly expected that the
literature scarcely bothers to record them. To see what that
expectation is calibrated against---and what a \emph{finished} theory
of pseudorandomness looks like---we pause (\cref{sec:prngs}).

The record deserves tabulating battery by battery, because the honest
picture distinguishes three things the phrase ``passes every test''
usually blurs: what has actually been \emph{run} and published, what is
merely \emph{expected} and therefore never written down, and what is
\emph{provably out of reach} of any feasible test at all
(\cref{tab:batteries}).

\begin{table}[H]
  \centering\small
  \renewcommand{\arraystretch}{1.25}
  \begin{tabular}{@{}p{0.15\linewidth} p{0.235\linewidth}
                     p{0.35\linewidth} p{0.115\linewidth}@{}}
    \toprule
    \textbf{Battery (size)} & \textbf{What it would catch} &
    \textbf{Status on the SHA family} & \textbf{Source} \\
    \midrule
    NIST SP\,800-22 \newline (15 tests) &
    monobit \& block frequency, runs, matrix rank, spectral, template
    matching, Maurer universal, linear complexity, serial, approximate
    entropy, cusums, excursions &
    The standard instrument for cryptographic RNGs. The \SHA{} counter
    generator is standardized as \textsf{Hash\_DRBG} and passes; a
    failure would be reportable and none is reported. Passing is
    \emph{expected}, so rarely published as such. &
    \citet{NIST2010}; \citet{SP80090A} \\
    TestU01 Crush / BigCrush \newline (${\sim}96$/$106$ tests, $160$ stats) &
    birthday spacings, collisions, gaps, poker, rank, linear complexity,
    Hamming-weight, random-walk statistics &
    \textbf{Documented run.} SHA-1 (OFB and CTR modes) passes
    \emph{all}; AES likewise, bar one $p$-value that cleared on
    replication. \SHA{} itself was not in that study, but is expected to
    behave identically. &
    \citet{LEcuyerSimard2007} \\
    DIEHARD / Dieharder \newline (${\sim}30$ tests) &
    birthday spacings, overlapping permutations, binary rank, ``monkey''
    (bitstream) tests &
    Routinely applied to cryptographic generators and passed; we know of
    no \SHA{}-specific published study---passing is folklore, exactly as
    the text describes. &
    Marsaglia \newline (software) \\
    PractRand \newline (streams to ${\gtrsim}$\,TB) &
    bit-pattern and transition anomalies visible only at very large
    sample volume &
    Applied to \SHA{}-based counter generators and passed to the volumes
    tested. A software tool, with no formal paper to cite. &
    --- (software) \\
    \addlinespace[1pt]
    \emph{Global $256$-bit uniformity / closeness} &
    \emph{any} constant-$\varepsilon$ deviation of the \emph{full} output
    distribution from uniform &
    \textbf{Not feasible.} Requires ${\sim}2^{128}$ samples---the
    birthday wall as an information-theoretic lower bound; provably
    undecidable at any affordable cost (\cref{sec:disttest}). &
    \citet{Canonne2020} \\
    \bottomrule
  \end{tabular}
  \caption{Standard statistical batteries and their verdict on the
  \SHA{} family. Only the TestU01 row is a \emph{documented, published}
  run---and it tested SHA-1, not \SHA{}. The SP\,800-22 row rests on
  standardization and the absence of reported failure; the DIEHARD and
  PractRand rows on routine practice that the literature, expecting
  success, does not record; and the last row is not a gap in effort but
  a theorem (\cref{sec:disttest}) that no feasible sample can close.
  Every battery that \emph{has} been run reports uniform; the one that
  would see the whole distribution cannot be run.}
  \label{tab:batteries}
\end{table}

It is worth being concrete about \emph{which} regularities are absent
and which are present, since ``passes the batteries'' is too coarse to
convey the texture of the thing. A ledger.

\emph{Regularities \SHA{} does \textbf{not} exhibit} (all empirical, none
proven absent):
\begin{itemize}
  \item \textbf{No lattice structure.} Successive output tuples show no
  Marsaglia effect: where the $d$-tuples of a linear congruential
  generator lie on at most $(d!\,m)^{1/d}$ hyperplanes
  (\cref{sec:prngs};~\citealp{Marsaglia1968}), the points
  $(\SHA(n), \SHA(n{+}1), \dots)$ display no detectable hyperplane
  concentration, no lattice, no low-dimensional alignment, in any
  projection anyone has computed.
  \item \textbf{No observable period.} The counter stream $\SHA(1),
  \SHA(2), \dots$ has no detected cycle; viewed as iteration
  (\cref{sec:randommaps}) the expected recurrence time is $\sim 2^{128}$,
  operationally unreachable. Contrast the classical generators, whose
  periods are exact theorems ($2^{19937}-1$ for the Mersenne
  Twister;~\citealp{MatsumotoNishimura1998}).
  \item \textbf{No linear or low-degree structure.} The $\Ftwo$-linear
  complexity of the output stream is empirically near-maximal---no low
  linear complexity, the exact defect that sinks the Twister on
  BigCrush---and no correlation of any output bit with a
  low-degree Boolean function of the input has been found. Each output
  bit, as a function of the input, appears to have a flat
  Walsh--Hadamard spectrum: no good linear approximation, the property
  linear cryptanalysis would exploit.
  \item \textbf{No first-order statistical bias.} Output bit frequencies,
  runs, gaps, block frequencies, and autocorrelations at every measured
  lag sit within sampling noise of the uniform expectation---the content
  of ``passes SP\,800-22 and BigCrush.''
\end{itemize}

\emph{Regularities \SHA{} \textbf{does} exhibit} (these are the true
structure, and they are exactly the harmless kind):
\begin{itemize}
  \item \textbf{Determinism.} It is a fixed function: $\SHA(x)$ is always
  the same, the one perfectly reliable ``pattern,'' and the sole basis of
  every use in \cref{sec:uses}.
  \item \textbf{The avalanche property.} Flipping a single input bit
  flips each of the $256$ output bits with probability indistinguishable
  from $\tfrac12$, and near-independently, so the number of changed
  output bits is essentially $\mathrm{Binomial}(256, \tfrac12)$, sharply
  peaked at $128$. This \emph{strict avalanche criterion}, a deliberate
  design target since~\citet{WebsterTavares1985}, is a regularity in the
  sense that it holds tightly and measurably---a property of
  \emph{structurelessness itself}, and the subject of a project below.
  \item \textbf{Completeness/diffusion.} After the full round count every
  output bit depends on every input bit; there is no input coordinate any
  output ignores.
\end{itemize}

The avalanche property is the one to dwell on, because it dictates the
self-correlation structure---or rather its total absence. Consecutive
counter values $n$ and $n+1$ differ, as bit strings, in just a few
low-order bits; avalanche then forces $\SHA(n)$ and $\SHA(n+1)$ to differ
in about half of \emph{their} bits, with no residual correlation between
them. So the output stream has essentially zero autocorrelation at every
lag and no continuity in any metric on the inputs: nearness of arguments
carries no information about nearness of values. This is precisely the
opposite of the individual ARX primitives, each of which is $2$-adically
$1$-Lipschitz---continuous, correlation-preserving
(\cref{sec:statespace})---and it is the composition of $64$ rounds that
annihilates that continuity. The designed regularity (avalanche) is the
mechanism that manufactures the absence of every other regularity a
statistical test could name.

\emph{Cryptanalysis} is the adversarial culture: bespoke,
design-specific, and quantitative in a different way---its currency is
the number of rounds that can be broken and at what cost.
\Cref{sec:differential} anatomizes the method behind these numbers; the
headline state of the art on \SHA{}: practical collisions for the
compression function reduced to 28 of its 64 rounds and theoretical
ones at 31~\citep{MendelNadSchlaffer2013}, ``semi-free-start''
collisions (the attacker also chooses the chaining value) at
39~\citep{LiLiuWang2024}, and preimage attacks via bicliques reaching
about 45 rounds at cost barely below brute force~\citep{KRS2012}. Two readings of these numbers
coexist. Comforting: after twenty years of world-class effort, half the
rounds stand untouched, and nothing approaches the full function.
Discomfiting: the same community's tools went, for SHA-1, from academic
paper to industrial collision in one decade (\cref{sec:falling}), and
there is no theory predicting which trajectory \SHA{} is on. The round
margin is a measurement of \emph{our} progress, not of \emph{its}
strength.

\subsection{How much can statistics even see? The sublinear ceiling}
\label{sec:disttest}

Knuth's moral---testing refutes, never confirms---has a quantitative
refinement worth stating, because it says precisely \emph{which}
questions about \SHA{}'s output distribution are answerable with a
feasible number of samples and which are not. The relevant theory is
\emph{distribution testing}, the sublinear-algorithms study of how many
independent samples from an unknown distribution over a domain of size
$N$ are needed to decide a property of it; the modern survey is
\citet{Canonne2020}.

The foundational result is the sample complexity of \emph{uniformity
testing}: to distinguish ``the distribution is uniform on $N$ points''
from ``it is $\varepsilon$-far from uniform in total-variation
distance,'' one needs, and suffices with, $\Theta(\sqrt{N}/\varepsilon^{2})$
samples. The upper bound is achieved by a startlingly simple statistic---%
\emph{count collisions}: draw $m$ samples, count coinciding pairs, and
compare against the $\binom{m}{2}/N$ expected under uniformity, a test
going back to Goldreich and Ron and made sample-optimal by Paninski
(both in~\citealp{Canonne2020}). The reader will recognize the birthday
calculus of \cref{sec:randommaps}: collisions among $\sqrt{N}$ samples
are exactly the signal, and $\sqrt{N}$ is exactly the threshold.

Point this at \SHA{}'s $256$-bit output distribution ($N = 2^{256}$) and
the ceiling is stark: detecting \emph{any} constant-$\varepsilon$
deviation of the global output distribution from uniform requires on the
order of $2^{128}$ samples---the birthday wall again, now wearing the
hat of an \emph{information-theoretic lower bound} on testing rather than
an algorithmic cost. No battery, however clever, that draws fewer than
$\sim 2^{128}$ outputs can certify or refute uniformity of the full
distribution; the question is not hard, it is \emph{unanswerable} at
feasible sample sizes. So whatever SP\,800-22 and BigCrush are doing when
they ``pass \SHA{},'' they are emphatically \emph{not} testing the
$256$-bit distribution.

What they \emph{are} testing is the reconciling point, and it is exactly
the sublinear-feasible part. Every test in those batteries is a
low-dimensional or low-complexity statistic---the frequency of single
bits, of short blocks, of runs; linear complexity; a fixed template
count---each of which lives on a domain small enough ($2^{k}$ for modest
$k$, or a real-valued summary) that $\sqrt{\,\cdot\,}$ samples suffice
and a laptop has them. Passing means: on every \emph{marginal} anyone has
thought to compute and can afford to test, \SHA{} matches uniform. The
$2^{128}$ ceiling says the joint, high-dimensional distribution is
untestable; the batteries probe its cheap shadows, and the shadows are
clean. That precise divorce---feasible low-dimensional marginals all
uniform, infeasible global distribution forever unknown---is the honest
content of ``passes every test,'' and it is a sharper statement than the
phrase usually carries.

Two further refinements are worth flagging because each squeezes signal
from limited samples rather than brute count. First, \emph{closeness}
(two-sample) testing: deciding whether \SHA{}'s output distribution
\emph{equals} a reference distribution, rather than testing it against
the explicit uniform, has its own sublinear theory~\citep{Batu2013}, and
it is the right frame for the comparative program of \cref{sec:randommaps}---%
one does not test \SHA{} against ``uniform'' in the abstract but against
a measured random-map \emph{control}, in the low-dimensional Flajolet
coordinates where $\sqrt{\,\cdot\,}$ samples reach. Second, the
\emph{$p$-value-of-$p$-values} trick that SP\,800-22 actually
uses~\citep{NIST2010}: run one weak test on many independent streams,
then test whether the resulting $p$-values are themselves uniform on
$[0,1]$. A bias too faint for any single run to flag becomes visible in
the second-order distribution of the $p$-values---a genuine way to
amplify a subtle distributional signal without paying for
prohibitively many samples per test. Against \SHA{} these meta-tests, too,
report uniform. None of this proves anything; but it locates, to the
sample, the boundary between what the empirical program can decide and
what it provably cannot.

\begin{project}{Collision testing at the birthday threshold}{probability;
  programming}{4--6 weeks}
  Implement the collision-count uniformity tester on the $n$-bit output
  of truncated \SHA{} for $n$ up to ${\sim}48$, and measure its
  \emph{sample complexity} empirically: how many outputs are needed to
  reject uniformity for reduced-round or deliberately-biased variants,
  and how does the number scale with $n$? Confirm the
  $\Theta(\sqrt{2^{n}})$ law~\citep{Canonne2020} and watch it become the
  same $2^{n/2}$ wall the collision-hunting project of
  \cref{sec:telescope} runs into---one barrier, met from the
  testing side and the attacking side at once. Then add the
  $p$-value-of-$p$-values layer and see how much fainter a bias it can
  catch at fixed budget.
\end{project}

\subsection{Interlude: what a finished theory looks like}\label{sec:prngs}

Two mature bodies of mathematics sit directly adjacent to our object,
and both show what ``understanding a pseudorandom process'' means when
understanding is actually achieved.

\paragraph{Pseudorandom numbers and finite fields.}
Pseudorandom number generators were born with the stored-program
computer: Monte Carlo simulation demanded rivers of reproducible random
digits, physical sources were slow and unrepeatable, and by mid-century
von Neumann's middle-square method and Lehmer's linear congruential
generator $x_{n+1} = a x_n + c \bmod m$ had founded the subject (von
Neumann also supplied its founding joke, about the ``state of sin'' of
anyone producing random digits by arithmetic). The subject's demands
are precise and \emph{non-adversarial}: a simulation needs
equidistribution of the statistics it consumes, a long period, and
above all speed---a physics run may consume $10^{12}$ variates, so a
generator must cost a few word operations per output. The discipline
was codified by \citet[Ch.~3]{Knuth1997}, and its mature mathematical
home is the arithmetic of finite fields, as systematized by Lidl and
Niederreiter~\citep{LidlNiederreiter1997, Niederreiter1992}. And here
the theorems are \emph{theorems}. A congruential generator attains full
period exactly when three elementary congruence conditions on
$(a, c, m)$ hold \citep[the Hull--Dobell criterion;][\S 3.2.1]{Knuth1997}.
A linear recurrence over $\Ftwo$ with characteristic polynomial $f$ has
period equal to the order of $f$, maximal $2^k - 1$ precisely when $f$
is primitive~\citep[Ch.~8]{LidlNiederreiter1997}---which is how
Matsumoto and Nishimura could \emph{prove} that their Mersenne Twister
has period $2^{19937}-1$ and equidistribution in $623$
dimensions~\citep{MatsumotoNishimura1998}. Marsaglia \emph{proved} that
successive $d$-tuples from any congruential generator fall on at most
$(d!\,m)^{1/d}$ hyperplanes~\citep{Marsaglia1968}---for $d = 10$ and
$m = 2^{32}$, a few dozen planes carrying every point the generator
will ever produce. For nonlinear congruential generators,
equidistribution is controlled through exponential-sum estimates, so
that Weil's Riemann hypothesis for curves over finite fields quietly
underwrites the uniformity of practical
generators~\citep{Niederreiter1992}. Classical pseudorandomness, in
short, is a satellite of the arithmetic of finite fields.

Now the double-edged contrast. These theorems exist \emph{because} the
generators are algebraically transparent---one clean structure, fully
exposed---and the same transparency is what makes them cryptographically
worthless: Marsaglia's planes and the Twister's $\Ftwo$-linearity are
simultaneously the content of the theorems and the substance of the
distinguishers (\cref{sec:below}). Efficiency tells the story from the
cost side. A Twister step is a handful of word operations; \SHA{}
spends on the order of a few thousand word operations per $256$-bit
output---one to two orders of magnitude more work per bit in
software---and the entire surcharge purchases exactly one property,
resistance to an intelligent adversary, which no simulation needs.
Where the classical theory lives inside a single algebraic structure,
\SHA{} deliberately straddles two incompatible ones
(\cref{sec:statespace}) precisely so that no analogue of the
Lidl--Niederreiter machinery can grip it. That is the trade: the
theorems stop where the adversary begins.

\paragraph{Random permutations and the symmetric group.}
The second finished theory concerns a different pseudorandom object: a
\emph{shuffle} is a pseudorandom number generator valued in the
symmetric group, and ``how many riffle shuffles randomize a deck?'' is
the question ``how many rounds suffice?'' asked of cards. Here the
theory reached a summit. Modeling repeated riffle shuffles as a random
walk on $S_{52}$, \citet{BayerDiaconis1992} computed the total-variation
distance to uniformity essentially exactly, exhibiting the famous
answer (seven shuffles) and, more importantly, the \emph{cutoff
phenomenon}. Their theorem is worth quoting in the exact form that makes
the abruptness visible. Riffle-shuffling a deck of $n$ cards
$k = \tfrac{3}{2}\log_2 n + \theta$ times, the total-variation distance
to the uniform distribution obeys
\[
  \bigl\| Q^{\ast k} - U \bigr\|_{\mathrm{TV}}
  \;=\; 1 - 2\,\Phi\!\Bigl(\tfrac{-2^{-\theta}}{4\sqrt{3}}\Bigr) + o(1),
  \qquad n \to \infty,
\]
where $\Phi$ is the standard normal cumulative distribution function
\[
  \Phi(z) \;=\; \frac{1}{\sqrt{2\pi}} \int_{-\infty}^{z}
                 e^{-t^{2}/2}\, \mathrm{d}t
\]
\citep[Thm.~1]{BayerDiaconis1992}, so that the approach to uniformity is
governed by the Gaussian tail $\int e^{-t^2/2}\,\mathrm{d}t$ evaluated at
the exponentially small argument $-2^{-\theta}/(4\sqrt 3)$. The scaling parameter
$\theta$ measures shuffles past the threshold $\tfrac{3}{2}\log_2 n$;
because the right-hand side collapses from $1$ to $0$ as $\theta$ runs
through an $O(1)$ window while the threshold itself grows like
$\log_2 n$, the transition is a true cutoff. For $n = 52$ this puts the
knee at $k \approx 8.55$ and yields distances dropping from $1.000$ at
$k=5$ to $0.334$ at $k=7$ and $0.167$ at $k=8$: the origin of ``seven
shuffles.'' The technical engine behind such exact statements is the
representation theory of finite groups: for the walk generated by random
transpositions, \citet{DiaconisShahshahani1981} proved cutoff at
$\tfrac12 n \log n$ by Fourier analysis on $S_n$, with the sharp
exponential rate
$\| Q^{\ast k} - U \|_{\mathrm{TV}} \le \tfrac12\, e^{-2c}$ once
$k = \tfrac12 n(\log n + c)$---the irreducible representations, indexed
by Young diagrams, diagonalize the walk, and a single upper-bound lemma
converts character sums into sharp convergence inequalities. The method is expounded in Diaconis's
monograph~\citep{Diaconis1988}, and it is the standing model of what a
\emph{hard convergence theorem} for an iterated mixing process looks
like: not ``eventually close to uniform,'' but \emph{exactly when},
with matching bounds.

\paragraph{Random walks on the hypercube, and the cutoff phenomenon.}
Between the card-shuffling summit and our object sits the solved mixing
problem closest to \SHA{}'s own geometry: the random walk on the
hypercube $\Ftwo^n$---the very space of \cref{sec:statespace}, walked
one coordinate at a time. At each step, pick one of the $n$ bit
positions uniformly and replace that bit by a fresh coin flip (this is
the Ehrenfest urn of statistical mechanics, wearing bits). How many
steps until the position is uniform on all $2^n$ points? The answer is
a second exact cutoff theorem: the total-variation distance stays
pinned near $1$ until $\tfrac14\, n \ln n$ steps, then collapses to $0$
inside a window of width $O(n)$---asymptotically instantaneous
(\cref{fig:cutoff}; the phenomenon, its history, and this example are
surveyed in~\citealp{Diaconis1996}). The threshold is pure coupon
collecting: a coordinate never yet touched is a coordinate whose bit is
still the initial one, and $\tfrac14 n \ln n$ is when the last
stragglers have been visited; the logarithm is the price of the
\emph{slowest} coordinate, not the average one---a moral worth
remembering whenever ``rounds needed to mix'' is discussed.

Better still, the machinery is transparent enough to expose a
coincidence this paper finds genuinely striking. On the abelian group
$\Ftwo^n$ the Fourier analysis that diagonalized card shuffles becomes
elementary: the characters are the Walsh functions
$x \mapsto (-1)^{\langle u, x\rangle}$, every translation-invariant
walk has them as eigenvectors, and \emph{the eigenvalues are exactly
the Walsh--Hadamard coefficients of the one-step distribution}. Mixing
is slow precisely in the directions $u$ where that coefficient is
large. Now reread the empirical ledger of \cref{sec:below}: ``flat
Walsh--Hadamard spectrum, no good linear approximation'' is the
property \emph{linear cryptanalysis} probes. The quantity the
cryptanalyst hunts and the quantity that governs mixing are the same
number: a linear distinguisher is a slow eigenvector, caught in the
act. Mixing analysis and linear cryptanalysis are, in this abelian
setting, one computation performed by two communities that rarely cite
each other---and the projects of \cref{sec:gap} amount to measuring
\SHA{}'s spectral gap experimentally, round by round. For scale: at
$n = 256$ the theorem's threshold is $\tfrac14\, n \ln n \approx 355$
single-coordinate updates, while one \SHA{} round nominally touches all
$256$ state bits at once---but in correlated fashion, which is exactly
why no theorem applies and the measured curve is the interesting
object.

\begin{figure}[ht]
  \centering
  \begin{tikzpicture}[font=\small, xscale=3.1, yscale=2.3]
    % axes
    \draw[->] (0,0) -- (2.25,0) node[below left=0.5mm and -6mm]
      {steps $\div$ $\tfrac14 n \ln n$};
    \draw[->] (0,0) -- (0,1.18) node[rotate=90, above left=1mm and 3mm]
      {distance to uniform};
    \draw (0.03,1) -- (-0.03,1) node[left, font=\scriptsize] {$1$};
    \draw (1,0.03) -- (1,-0.03) node[below, font=\scriptsize] {$1$};
    \draw[dashed, gray] (1,0) -- (1,1.05);
    % cutoff curves, sharpening with n (precomputed: 0.5*(1-tanh(a(x-1))))
    \draw[thick, msgfill!60!black] plot[smooth] coordinates {
      (0.00,0.982) (0.20,0.961) (0.40,0.917) (0.60,0.832) (0.80,0.690)
      (1.00,0.500) (1.20,0.310) (1.40,0.168) (1.60,0.083) (1.80,0.039)
      (2.00,0.018) (2.20,0.008)};
    \draw[thick, projframe] plot[smooth] coordinates {
      (0.00,1.000) (0.20,0.998) (0.40,0.992) (0.60,0.961) (0.70,0.917)
      (0.80,0.832) (0.90,0.690) (1.00,0.500) (1.10,0.310) (1.20,0.168)
      (1.30,0.083) (1.50,0.018) (1.70,0.004) (2.00,0.000) (2.20,0.000)};
    \draw[very thick, citewine] plot[smooth] coordinates {
      (0.00,1.000) (0.40,1.000) (0.60,0.999) (0.70,0.996) (0.80,0.973)
      (0.90,0.858) (1.00,0.500) (1.10,0.142) (1.20,0.027) (1.30,0.004)
      (1.50,0.000) (1.80,0.000) (2.20,0.000)};
    \node[font=\scriptsize, msgfill!60!black, anchor=west] at (1.32,0.30)
      {$n$ small};
    \node[font=\scriptsize, citewine, anchor=west] at (1.13,0.10)
      {$n$ large};
    \node[font=\scriptsize, gray, anchor=south, align=center]
      at (1.72,0.62) {window shrinks\\ relative to threshold};
    \draw[->, gray, font=\scriptsize] (1.45,0.60) -- (1.12,0.45);
  \end{tikzpicture}
  \caption{The cutoff phenomenon, schematically: distance to
  uniformity for the hypercube walk stays near $1$, then collapses
  within a window that is asymptotically negligible compared to the
  $\tfrac14 n \ln n$ threshold. ``How many rounds
  are enough?'' has, for this toy relative of \SHA{}, an answer sharp
  to within lower-order terms---and the toy-compression-walk project
  of this section asks whether \SHA{}-family walks show the same
  knee.}
  \label{fig:cutoff}
\end{figure}

The bridge to our object is not an analogy but a change of alphabet.
Every keyed block cipher is a family of permutations, and \SHA{}'s own
compression function is one in disguise (its core is a block cipher,
run in Davies--Meyer mode, \cref{sec:roundshape}). More directly: the
chaining recursion $h_i = f(h_{i-1}, m_i)$ of \cref{sec:merkledamgard},
fed \emph{random} message blocks, is literally a random walk on the
state space $\State$, driven by $2^{512}$ generators. Diaconis's
questions transfer verbatim---how many blocks until the distribution of
$h_i$ is uniform? is there cutoff? what is the spectral gap?---and, to
our knowledge, none has been answered, or seriously asked, for any
member of the SHA family at any scale. The obstruction is instructive:
the representation-theoretic engine wants a group acting on itself, and
the walk lives merely on a set---unless one exploits the two group
structures the set actually carries (\cref{sec:statespace}). Here is a
place where the paper's finite geometry and a fully developed classical
theory face each other across a gap that small experiments can begin to
survey.

\begin{project}{Cutoff for toy compression walks}{probability; linear
  algebra; programming}{4--6 weeks}
  Shrink the compression function (\cref{sec:roundshape}) to an $8$- or
  $16$-bit state. Random message blocks then make chaining a random
  walk on $2^{8}$ or $2^{16}$ points: build its transition matrix,
  compute the spectral gap, plot total-variation distance to uniformity
  against the number of blocks, and look for a cutoff profile. Vary the
  rounds-per-block and watch the gap and the cutoff window move. Every
  quantity here is exactly computable at toy scale, and the questions
  are verbatim those of the shuffling literature.
\end{project}

The interlude's moral for \cref{sec:gap}: on one flank of \SHA{} stands
a complete algebraic theory of linear generators, on the other a
complete probabilistic theory of mixing walks on groups; the function
itself was engineered into the mathematical no-man's-land between
them---touching complex analysis through \cref{sec:randommaps}, finite
fields and $2$-adic analysis through \cref{sec:statespace},
representation theory through the walk just described, and complexity
theory through \cref{sec:definitions}---and nearly every one of those
contacts is an unfollowed lead rather than a developed theory.

\subsection{The gap, stated plainly}\label{sec:gap}

Now place the object among its relatives, the pseudorandom number
generators, along two informal axes: how fast it is, and how much theory
supports its randomness claims (\cref{fig:prngspace}).

\begin{figure}[ht]
  \centering
  \begin{tikzpicture}[scale=1.0, font=\small]
    \draw[->] (0,0) -- (8.8,0) node[below left=1mm and 0mm]
      {analytical understanding};
    \draw[->] (0,0) -- (0,4.8) node[rotate=90, above left=1mm and 14mm]
      {speed / practicality};
    % fully understood, weak
    \node[align=center] (lcg) at (7.1,4.2)
      {LCG\\ \scriptsize lattice structure: solved, weak};
    \node[align=center] (mt) at (6.3,3.1)
      {Mersenne Twister\\ \scriptsize $\Ftwo$-linear: solved, distinguishable};
    % provable, slow
    \node[align=center] (bbs) at (6.7,0.9)
      {number-theoretic PRGs\\ \scriptsize provable from factoring-type\\
       \scriptsize assumptions, slow};
    % SHA
    \node[draw, very thick, circle, inner sep=1.5pt, fill=projfill,
          align=center] (sha) at (1.2,4.05) {\SHA};
    \node[gray, font=\scriptsize, align=left, anchor=west]
      at (1.95,4.05) {passes every test,\\ no theory at all};
    % the gap: brace bulging toward the empty lower-left region
    \draw[decorate, decoration={brace, mirror, amplitude=6pt}]
      (2.2,3.5) -- (5.0,1.55);
    \node[rotate=-35, font=\footnotesize] at (2.75,1.7)
      {\emph{the gap}: no intermediate landmarks};
  \end{tikzpicture}
  \caption{A cartoon of PRNG space. Classical generators are fast and
  fully understood---which is exactly how one knows they are weak.
  Number-theoretic generators carry proofs (conditional on hardness
  assumptions) and are slow. \SHA{} sits alone: fastest, empirically
  strongest, and analytically blank.}
  \label{fig:prngspace}
\end{figure}

The classical generators are \emph{transparent}: linear congruential
generators are so well understood that their defects (Marsaglia's
lattice planes) are theorems; twisters are $\Ftwo$-linear and hence
completely analyzable, and completely distinguishable. At the other
extreme, number-theoretic generators (of Blum--Micali type) come with
genuine reductions---distinguishing their output is provably as hard as
factoring-type problems---purchased at a price of thousands of arithmetic
operations per output bit. \SHA{} belongs to neither culture: it has
\emph{no} reduction to any studied hard problem, \emph{no} algebraic
model that survives even a handful of rounds, and \emph{no} distinguisher
after twenty-five years. It is simultaneously the best-tested and least
theorized object in the entire space.

Why should the gap be so large? Part of the answer is a legitimate
excuse: an unconditional proof that \SHA{} (or anything) is a one-way or
pseudorandom function would imply $\mathrm{P} \ne \mathrm{NP}$ and more;
in Impagliazzo's celebrated taxonomy of possible computational worlds,
whether we live in a world with one-way functions at all
(``Minicrypt''/``Cryptomania'' versus ``Algorithmica''/``Pessiland'')
is precisely the great open question~\citep{Impagliazzo1995}. Nobody
should demand a proof. But the excuse covers less than the gap. Between
``unconditional proof'' and ``nothing'' there is an enormous unexplored
middle: reductions to \emph{any} clean assumption; partial structure
theorems for few rounds; quantitative mixing results in either of the two
geometries of \cref{sec:statespace}; distributional results in Flajolet's
coordinates for scaled-down families (\cref{sec:randommaps}). For linear
generators that middle territory is dense with theorems; for ARX
functions it is nearly empty---not, as far as we can tell, because
theorems are known to be impossible, but because the object fell into a
disciplinary crack: too applied for one community, too structureless for
another, and (a hypothesis this prospectus exists to test) never yet
attacked with the analytic-combinatorics instrument at the appropriate
model-organism scale.

That is the research program: not ``prove \SHA{} secure''---nobody sane
proposes that---but \emph{populate the gap with measurements}: an atlas
of scaled-down \SHA{}-family maps (fewer rounds, narrower words, fewer
words), located in Flajolet's coordinate system, with the two-geometry
anatomy of \cref{sec:statespace} guiding which scaled families to
compare. Every panel of that atlas is a small, finishable, publishable
piece of natural history; \cref{sec:projects} cuts it into
undergraduate-sized pieces.

\begin{project}{Where do the tests start failing?}{programming;
  statistics}{4--6 weeks}
  Reduced-round \SHA{} must fail statistical batteries---at one round it
  is nearly raw message. Find the phase transition: run
  SP\,800-22~\citep{NIST2010} and TestU01~\citep{LEcuyerSimard2007}
  against generators built from $r$-round \SHA{} for $r = 1, 2, \dots$
  and chart which tests fail at which $r$. Compare with the round counts
  where cryptanalysis stalls (\cref{sec:below}). The two ``difficulty
  frontiers'' have, to our knowledge, never been plotted on one axis.
\end{project}

\begin{project}{Avalanche cartography}{programming; linear
  algebra}{4--6 weeks}
  Flip one input bit; which of the 256 output bits flip? For the true
  random-function silhouette each flips independently with probability
  $\tfrac12$~\citep{WebsterTavares1985}. Measure the full $512 \times 256$
  influence matrix round by round, render it as heatmaps, and quantify
  convergence to $\tfrac12$ (in what norm? after how many rounds?
  uniformly over bit positions?). Then close the loop to the ledger of
  \cref{sec:below}: confirm that full-round avalanche forces the counter
  stream's lag-$1$ autocorrelation to zero, and measure how many rounds
  are needed before that self-correlation vanishes. Relate the observed
  diffusion speed to the rotation constants of
  \cref{sec:roundshape}---the designers' choices become visible as
  geography.
\end{project}

\subsection{The road not taken: hashes with theorems}\label{sec:provable}

One region of \cref{fig:prngspace} deserves its own visit before we
leave the map: the lower right, where proofs live and speed dies. It is
not empty of hash functions. \emph{Provably secure} hash functions
exist, several are beautiful mathematics, and the world declined every
one of them; both halves of that sentence are instructive.

Exhibit one: VSH, ``very smooth hash''~\citep{ContiniLenstraSteinfeld2006}.
Its collision resistance is an honest theorem, conditional on a single
clean number-theoretic assumption (finding congruences of a certain
smooth-square kind modulo a hard-to-factor $n$), and it is fast by the
standards of its genre---only $25$ times or so slower than the SHA
family, where earlier provable designs had been slower by thousands.
But the same multiplicative structure that powers the proof is
\emph{visible in the outputs}: VSH is not remotely a random-oracle
stand-in (short relations among hash values exist by construction), so
of the whole catalogue of \cref{sec:uses} it can serve collision-bound
uses only. The proof purchases one property by installing exactly the
kind of algebraic transparency that \cref{sec:prngs} showed to be fatal
to all the others.

Exhibit two is closer to this paper's geometry. \emph{Cayley hashes}
walk a graph: fix a group $G$ and two generators $g_0, g_1$; hash a bit
string by reading it as directions,
$m_1 m_2 \cdots m_k \mapsto g_{m_1} g_{m_2} \cdots g_{m_k}$. Tillich
and Z{\'e}mor proposed $G = \mathrm{SL}_2(\mathbb{F}_{2^n})$ in
1994~\citep{TillichZemor1994}; Charles, Goren, and Lauter modernized
the idea with Lubotzky--Phillips--Sarnak Ramanujan graphs and, in a
second variant, supersingular isogeny graphs of elliptic
curves~\citep{CharlesGorenLauter2009}. Two theorems come for free.
A collision is a pair of distinct words with the same product---so
collision resistance \emph{is} a group-theoretic statement (no short
relations among the generators are findable), inherited whole from a
recognizable mathematical problem. And output quality is spectral: if
the Cayley graph is an \emph{expander}---nontrivial eigenvalues
uniformly bounded away from the degree---then the distribution of
hashes of length-$k$ messages converges to uniform at an exponential
rate governed by the spectral gap~\citep{HooryLinialWigderson2006}.
The mixing mathematics of \cref{sec:prngs} returns, no longer as
analysis but as a \emph{design principle}: these are hash functions
whose avalanche property is a theorem about eigenvalues.

Then the pothole. In 2008 the LPS instantiation was broken---collisions
computed outright---by lifting the problem from the finite group to a
characteristic-zero arithmetic where Euclidean-style algorithms find
short relations; the authors of the break were Tillich and Z{\'e}mor,
the designers of 1994, and their own original $\mathrm{SL}_2$ proposal
later fell, for its published parameters, to a similar
lift~\citep{TillichZemor2008}. Note carefully what died. The theorems
were and remain true; the \emph{assumptions} were false---the
underlying group problems, young and algebraically rich, simply were
not hard. (The isogeny-graph sibling's assumption has so far survived,
and seeded a whole branch of post-quantum cryptography.) The moral is
the sober one: a security proof does not abolish risk, it relocates
risk from the object to the assumption, and an assumption chosen for
its provability can be weaker than ignorance. Set beside
\cref{sec:falling}: MD5 fell \emph{with} twenty years of massive use
and no theorem; LPS fell \emph{with} a theorem and almost no use. The
fossil record refuses to pick a side.

Why, then, does the world run on the theorem-free object? Speed is the
visible reason---two orders of magnitude---but the load-bearing one is
now easy to state with this section's vocabulary: every use of
\cref{sec:uses} beyond bare collision resistance needs the
random-oracle \emph{pretension}, behaving like structurelessness
itself, and a provable design is structured \emph{by construction};
its proof and its disqualification are the same fact. The market
equilibrium of \cref{fig:prngspace} is thus not an accident of
laziness: given the choice between an object about which something is
proved and something else is visibly false, and an object about which
nothing is proved and nothing is visibly false, every deployed system
chose the second, with Rogaway's human-ignorance calculus
(\cref{sec:definitions}) as the honest license. The gap is not a hole
in the market; it is the market.

\begin{project}{Cayley hashes by hand}{group theory; linear algebra;
  programming}{3--4 weeks}
  Implement the Tillich--Z{\'e}mor hash~\citep{TillichZemor1994} at toy
  scale: $\mathrm{SL}_2(\mathbb{F}_{2^n})$ for $n = 8, \dots, 16$, two
  generator matrices, hash $=$ matrix product read off the message
  bits. (i) Measure output equidistribution against message length and
  match the exponential rate to the spectral gap of the Cayley graph,
  computable by eigenvalue methods at this
  scale~\citep{HooryLinialWigderson2006}---the mixing theory of
  \cref{sec:prngs} as a bench experiment. (ii) Hunt collisions two
  ways, generic birthday search versus structure (palindromic words,
  elements of small order, subfield membership), and compare costs.
  (iii) As a stretch goal, read the lifting attack~\citep{TillichZemor2008}
  and implement its Euclidean core for tiny parameters. Deliverable: a
  working Cayley hash, gap-versus-equidistribution plots, and a
  collision gallery---provable cryptography's rise and fall on one
  laptop.
\end{project}

% ------------------------------------------------------------
% §6  Differential cryptanalysis: the adversary's calculus
% ------------------------------------------------------------
\section{The adversary's calculus: differential cryptanalysis}
\label{sec:differential}

\Cref{sec:below} reported the adversarial state of the art as bare round
counts. This section opens the instrument that produces those numbers.
Differential cryptanalysis is the only mathematical technology that has
ever killed a member of this design family---it felled MD5, SHA-0, and
SHA-1 (\cref{sec:falling})---and its core idea is elementary enough to
state in one line: \emph{since the function cannot be understood, study
its finite differences instead}. What follows is the calculus itself,
its startling secret history, the story of its mechanization, and an
honest account of where, on \SHA{}, it is currently stuck. Readers who
absorb this section will understand precisely what the ``round counts''
of \cref{sec:below} measure.

\subsection{A calculus of differences}\label{sec:diffcalc}

For a map $F$ on bit strings and a fixed \emph{input difference}
$\Delta$, consider the discrete derivative
\[
  D_{\Delta}F(x) \;=\; F(x \shaxor \Delta) \,\shaxor\, F(x),
\]
the exact finite-difference analogue of differentiation, with $\shaxor$
playing subtraction in the $\Ftwo$-geometry of \cref{sec:statespace}.
The adversary's question is distributional: for random $x$, what does
$D_{\Delta}F(x)$ look like?

Two observations power everything. First, \emph{differencing kills
constants and linear structure}. If $L$ is $\Ftwo$-linear---a rotation,
a shift, one of the $\Sigma/\sigma$ maps of \cref{sec:roundshape}---then
\[
  D_{\Delta}L(x) \;=\; L(\Delta) \qquad \text{for every } x:
\]
the derivative is a \emph{constant}, independent of the unknown input,
exactly as $\mathrm{d}/\mathrm{d}x$ annihilates constants and
linearizes. All the values the attacker does not know---the other
message words, the chaining state, the round constants $K_t$---drop out
of the linear part of every round. This is why one differentiates: a
difference propagates through the transparent parts of the circuit
deterministically, leaving a residue of uncertainty only at the
nonlinear gates.

Second, at those nonlinear gates---the modular additions $\shaplus$ and
the quadratic functions $\mathrm{Ch}, \mathrm{Maj}$---a difference does
not propagate deterministically but \emph{branches into a probability
distribution over output differences}. The fundamental case is a single
addition (\cref{fig:diffadd}). Give one addend a one-bit difference at
position $i$ and add the same unknown $y$ to both members of the pair:
the two sums differ at bit $i$ certainly, but whether the difference
stops there is decided by the \emph{carries}, which the $\Ftwo$-lens
cannot see. If the carries out of position $i$ agree---this happens for
exactly half of all $(x,y)$---the output difference is again the single
bit $e_i$; otherwise a carry-difference runs leftward, one further bit
at a time, each extra bit halving the probability. One position is
exceptional: a difference in the top bit ($i = 31$) propagates with
probability $1$, because the disagreeing carry falls off the end of the
word and is discarded by the reduction mod $2^{32}$. Attackers
therefore park differences in the most significant bit whenever they
can---a free move in an economy where everything else costs. The
complete theory of differential propagation through addition (arbitrary
difference patterns, exact probabilities, all computable in linear
time) was worked out by \citet{LipmaaMoriai2001}, and its headline is a
pleasing cost principle: \emph{the probability of the straight-through
transition is $2^{-(\text{number of active bits, msb excluded})}$}. To
first order, difference weight is money, and the exchange rate is a
factor of two per active bit. The bitwise functions obey the same
economy: at each bit where its inputs differ, $\mathrm{Ch}$ or
$\mathrm{Maj}$ passes or absorbs the difference depending on the values
of the other, unknown operands---a coin flip per active bit, one more
factor of $\tfrac12$ each.

\begin{figure}[ht]
  \centering
  \begin{tikzpicture}[font=\small,
      cell/.style={draw, minimum width=5.2mm, minimum height=5.2mm,
                   inner sep=0pt},
      dbit/.style={cell, fill=red!25},
      qbit/.style={cell, fill=lenfill},
      lbl/.style={font=\scriptsize}]
    % ---------------- left panel: difference at interior bit 3 ------
    \foreach \i [evaluate=\i as \x using {(7-\i)*0.52}] in {0,...,7} {
      \ifnum\i=3
        \node[dbit] (a\i) at (\x,0) {$1$};
      \else
        \node[cell] (a\i) at (\x,0) {};
      \fi
      \node[lbl, above=0.5mm of a\i] {\i};
    }
    \node[lbl, anchor=east] at (-0.55,0) {$\Delta x$};
    \draw[->] (1.82,-0.45) -- (1.82,-1.25)
      node[midway, right=1mm, lbl, align=left]
      {$\shaplus\, y$ (same $y$ on both\\ members of the pair)};
    \foreach \i [evaluate=\i as \x using {(7-\i)*0.52}] in {0,...,7} {
      \ifnum\i=3
        \node[dbit] (b\i) at (\x,-1.7) {$1$};
      \else
        \ifnum\i>3 \ifnum\i<7
          \node[qbit] (b\i) at (\x,-1.7) {?};
        \else
          \node[cell] (b\i) at (\x,-1.7) {};
        \fi
        \else
          \node[cell] (b\i) at (\x,-1.7) {};
        \fi
      \fi
    }
    \node[lbl, anchor=east] at (-0.55,-1.7) {$\Delta(x \shaplus y)$};
    \draw[->, thick, orange!80!black]
      (b3.north west) ++(-0.05,0.06) to[bend right=20] (-0.15,-1.05);
    \node[lbl, orange!80!black, align=right, anchor=east] at (-0.55,-0.85)
      {carry difference\\ runs leftward};
    \node[lbl, align=left, anchor=north west] at (-0.6,-2.45)
      {stops immediately: probability $\tfrac12$\\
       each further bit: another $\tfrac12$};
    % ---------------- right panel: msb difference --------------------
    \begin{scope}[xshift=6.4cm]
    \foreach \i [evaluate=\i as \x using {(7-\i)*0.52}] in {0,...,7} {
      \ifnum\i=7
        \node[dbit] (c\i) at (\x,0) {$1$};
      \else
        \node[cell] (c\i) at (\x,0) {};
      \fi
      \node[lbl, above=0.5mm of c\i] {\i};
    }
    \node[lbl, anchor=east] at (-0.55,0) {$\Delta x$};
    \draw[->] (1.82,-0.45) -- (1.82,-1.25)
      node[midway, right=1mm, lbl] {$\shaplus\, y$};
    \foreach \i [evaluate=\i as \x using {(7-\i)*0.52}] in {0,...,7} {
      \ifnum\i=7
        \node[dbit] (d\i) at (\x,-1.7) {$1$};
      \else
        \node[cell] (d\i) at (\x,-1.7) {};
      \fi
    }
    \node[lbl, anchor=east] at (-0.55,-1.7) {$\Delta(x \shaplus y)$};
    \draw[->, thick, gray]
      (d7.north west) ++(-0.02,0.08) to[bend right=20] ++(-0.5,0.35);
    \node[lbl, gray, align=left] at (0.6,-0.75) {carry falls\\ off the word};
    \node[lbl, align=left, anchor=north west] at (-0.6,-2.45)
      {top bit: probability $1$\\ (the attacker's free move)};
    \end{scope}
  \end{tikzpicture}
  \caption{How a one-bit difference passes through modular addition
  (bit $7$ = most significant, in this $8$-bit cartoon of a $32$-bit
  word). \emph{Left:} an interior difference survives unchanged only if
  the carries agree---probability $\tfrac12$---and every extra bit of
  carry-difference costs another factor $\tfrac12$---the Lipmaa--Moriai
  cost formula of the text. \emph{Right:} a difference in the top bit
  always passes intact, since the disagreeing carry is discarded mod
  $2^{32}$. Differential cryptanalysis views the whole of \SHA{}
  through this lens: linear pieces are free wires; every active bit at
  a nonlinear gate levies a toll of $\tfrac12$.}
  \label{fig:diffadd}
\end{figure}

Notice what has happened in the language of \cref{sec:statespace}: the
attacker has \emph{chosen the $\Ftwo$-geometry} and agreed to pay, in
probability, at every point where the design invokes the other one. All
of the analytical cost is concentrated exactly at the carry chains and
quadratic gates---the terms $\Ftwo$-linear algebra cannot see. There is
a dual calculus using arithmetic differences $x' - x \bmod 2^{32}$,
which tames the additions and instead pays at the xors and rotations;
the craft of the great hash attacks (\cref{sec:craft}) was, in large
part, the discovery that one should track \emph{both} lenses
simultaneously, bit by bit, in the form of signed differences. The
two-geometry tension that \cref{sec:statespace} presented as the
design's defense is thus also, precisely, the structure of the attack:
differential cryptanalysis is what the two-lens problem looks like when
an adversary refuses to give up.

\subsection{Trails, tolls, and free rounds}\label{sec:craft}

A single round is only the first step. A \emph{differential
characteristic} (or \emph{trail}) prescribes a difference for every
intermediate quantity through all the rounds---a full flight plan for
the perturbation---and, under a heuristic independence assumption, its
probability is the \emph{product} of the per-gate transition
probabilities along the way. An attacker who wants a collision needs a
trail beginning at a nonzero message difference and ending, after the
feedforward, at state difference \emph{zero}; if the trail has
probability $p$, then about $1/p$ random conforming attempts are
expected before one message pair follows the entire flight plan and the
hashes agree. The attack economy is thus brutally simple: \emph{total
toll} $=$ sum, over all rounds, of active bits at nonlinear gates;
\emph{cost} $\approx 2^{\text{toll}}$; the design's job is to force
every possible flight plan to accumulate an unpayable toll.

\begin{figure}[ht]
  \centering
  \begin{tikzpicture}[font=\small,
      inj/.style={draw, fill=red!25, minimum width=8mm, minimum height=5mm,
                  font=\scriptsize},
      cor/.style={draw, fill=msgfill, minimum width=8mm, minimum height=5mm,
                  font=\scriptsize},
      non/.style={draw, fill=white, minimum width=8mm, minimum height=5mm},
      lbl/.style={font=\scriptsize}]
    % round axis
    \foreach \r/\x in {t/0, t+1/1.15, t+2/2.3, t+3/3.45, t+4/4.6, t+5/5.75}
      \node[lbl] at (\x, 2.15) {$\r$};
    \node[lbl, anchor=east] at (-0.9, 1.55) {message diff.\ $\delta W$};
    \node[inj] at (0,1.55) {seed};
    \foreach \x in {1.15, 2.3, 3.45, 4.6, 5.75}
      \node[cor] at (\x,1.55) {corr.};
    % state difference weight bars
    \node[lbl, anchor=east] at (-0.9, 0.6) {state diff.\ weight};
    \draw[->] (-0.55,0.05) -- (7.0,0.05);
    \fill[red!50]    (-0.2,0.05) rectangle (0.25,0.65);
    \fill[red!35]    (0.95,0.05) rectangle (1.4,0.5);
    \fill[red!35]    (2.1,0.05) rectangle (2.55,0.5);
    \fill[red!35]    (3.25,0.05) rectangle (3.7,0.42);
    \fill[red!25]    (4.4,0.05) rectangle (4.85,0.3);
    \draw[very thick, projframe] (5.55,0.05) -- (7.0,0.05)
      node[pos=0.6, above, lbl, projframe] {zero again};
    % toll ledger
    \node[lbl, anchor=east] at (-0.9,-0.55) {toll this round};
    \foreach \x/\p in {0/2^{-1}, 1.15/2^{-3}, 2.3/2^{-3}, 3.45/2^{-2},
                       4.6/2^{-2}, 5.75/2^{-1}}
      \node[lbl] at (\x,-0.55) {$\p$};
    \node[lbl, anchor=west] at (6.35,-0.55) {$\Rightarrow\ \prod = 2^{-12}$};
    \draw[decorate, decoration={brace, mirror, amplitude=5pt}]
      (-0.35,-0.95) -- (6.1,-0.95)
      node[midway, below=6pt, lbl]
      {one \emph{local collision}: inject, shepherd, cancel};
  \end{tikzpicture}
  \caption{Cartoon of a \emph{local collision}, the modular building
  block of every attack on this family since Chabaud and Joux's 1998
  analysis of SHA-0 (cited in the text): a seed difference is injected
  through one
  message word, then five successive message-word corrections shepherd
  and finally cancel the disturbance, leaving the state clean. Each
  round levies its toll (\cref{fig:diffadd}); the trail's probability
  is the product. A full attack chains many local collisions along a
  pattern compatible with the message schedule---and the toll of the
  best full trail is what decides whether the design lives or dies.}
  \label{fig:localcollision}
\end{figure}

Three refinements turn this arithmetic into the attacks of
\cref{sec:falling}, and all three are conceptually simple. First,
\emph{local collisions} (\cref{fig:localcollision}): rather than one
monolithic trail, the attacker builds a short pattern---inject a
difference through one message word, cancel it a few rounds later
through others---and then chains copies of the pattern wherever the
message schedule permits~\citep{ChabaudJoux1998}. Second, \emph{message
freedom}: for the first sixteen rounds the words $W_t$ are the message
itself (\cref{sec:roundshape}), so the attacker does not \emph{hope}
the early tolls are paid---she \emph{solves} for message bits that pay
them deterministically. Wang's celebrated breaks of MD5 and
SHA-1~\citep{WangYu2005, WangYinYu2005} rested on pushing this
``message modification'' astonishingly deep; the practical effect is
that the toll meter starts running only around round $17$, and the
first quarter of the function contributes nothing to its own defense.
Third, \emph{automation}: what Wang's school did by virtuosic hand
calculation---juggling xor- and arithmetic-differences as signed,
bit-level conditions---was recast by \citet{DeCanniereRechberger2006}
as a constraint-propagation search over ``generalized conditions,''
mechanically discovering trails no human would find. The current
records on \SHA{} are held by exactly this lineage of tools, now
built on MILP, \textsc{sat}, and \textsc{smt}
solvers~\citep{MendelNadSchlaffer2013, LiLiuWang2024}, closing the loop
with \cref{sec:npsearch}: the adversary's calculus has itself become a
branch of solver engineering.

\subsection{A method twenty years secret}\label{sec:secrethistory}

The published history begins in 1990, when Biham and Shamir announced
differential cryptanalysis and used it to attack reduced variants of
DES, the American encryption standard of
1977~\citep{BihamShamir1991}. A curiosity surfaced immediately: DES,
designed in the early 1970s, was \emph{strangely well prepared} for the
new attack. Its nonlinear tables were, component for component, close
to optimally resistant; the method that broke every academic variant
of DES bounced off the real thing.

The explanation arrived in 1994, and it is one of the most remarkable
documents in the literature. Coppersmith, of the original IBM design
team, was cleared to reveal the design criteria and wrote them down
plainly: the IBM team had discovered differential cryptanalysis
themselves---internally, the ``T-attack''---\emph{in 1974}, and
hardened DES against it before standardization; the NSA, consulted on
the design, asked that the rationale be kept secret, since publishing
the criteria ``would reveal the technique of differential
cryptanalysis, a powerful technique that can be used against many
ciphers''~\citep{Coppersmith1994}. The central tool of public
cryptanalysis was thus twenty years old on the day it was discovered.
For sixteen of those years, every open researcher who studied DES was
probing defenses calibrated against a weapon whose existence they did
not suspect.

This is not merely a good story; it is a load-bearing datum for the
epistemology of \cref{sec:epistemology}, because the pattern
\emph{repeats within our own object's lineage}. Recall the footnote of
\cref{sec:history}: in 1994--95 the NSA withdrew SHA-0 and changed a
single rotation, citing an undisclosed weakness; the public rediscovery
of what that tweak defends against---a differential collision
attack---took until 1998~\citep{ChabaudJoux1998}. Twice measured, the
gap between the classified and the public understanding of this exact
attack has been twenty years and four years. The honest inference for a
naturalist: the map of \cref{sec:below} is the \emph{public} frontier,
a lower bound on the true state of knowledge, drawn by a community that
has historically been behind by years to decades---and \SHA{}'s design
parameters were set by the party that was ahead. The comfort and the
discomfort of that fact are left to the reader to apportion.

\subsection{The siege of \SHA{}, and where it stands}\label{sec:siege}

Why did the calculus shatter MD5 and SHA-1 while \SHA{}, made of the
same primitives, holds? The anatomy of \cref{sec:roundshape} contains
the visible answers. SHA-1's message schedule is entirely
$\Ftwo$-linear, so schedule differences propagate by exact linear
algebra and \emph{sparse} schedules---low total weight, hence low
toll---can be found by searching a linear code for low-weight words;
this is precisely the handle the collision attacks
gripped (\cref{sec:hope}). \SHA{} interleaves modular additions into
the schedule itself and taxes every round twice over: each of
$\Sigma_0, \Sigma_1$ fans any single active bit into three, both tracks
of the two-track recurrence pass through a quadratic gate every round,
and the $\sigma$-mixing makes a sparse schedule difference bloom
within a few steps. The result, empirically: low-toll trails simply
stop existing a dozen-odd rounds past the free zone, and two decades of
search---human, then mechanical---have moved the frontier only from
about $19$ rounds (2006) to $31$ rounds (collisions) today. The precise
current records: practical colliding pairs for $28$ rounds and a
$2^{65.5}$ collision attack on $31$ rounds of the $64$, both from the
automated-search school~\citep{MendelNadSchlaffer2013}; practical
\emph{semi-free-start} collisions (the attacker may choose the chaining
value, removing the Merkle--Damg{\aa}rd anchor of \cref{thm:md}) for
$38$ rounds in 2013 and $39$ rounds in 2024, the latter found with a
\textsc{sat}/\textsc{smt} trail search~\citep{LiLiuWang2024,
AlamgirNejatiBright2024}; and, on the inversion side, preimages to $45$
rounds at a cost indistinguishable from brute
force~\citep{KRS2012}. Li, Liu, and Wang open their record-setting 2024
paper with the flat sentence that ``there is almost no progress in
collision attacks on SHA-2 after ASIACRYPT 2015''---a decade of solver
revolutions purchasing one round.

The shape of the stall is worth internalizing, because it is the
quantitative content of the phrase ``security margin.'' Rounds $1$--$16$
are free (message freedom); rounds $17$ to roughly $40$ are the
contested zone, where automated searches still find trails but each
round multiplies the toll; beyond that, \emph{no usable trail has ever
been exhibited}, not because a theorem forbids one but because every
known search method drowns. The design's $64$ rounds thus stand at
about twice the depth the best siege engines have ever reached with
unlimited choice of starting position. Whether that factor of two is a
fortress wall or a picket fence is exactly the question
\cref{sec:falling} showed nobody can answer from first principles: no
theory predicts the frontier's future velocity, and the SHA-1
precedent---from academic $2^{69}$ to industrial collision in twelve
years---is the reason nobody treats the margin as a moat. (One further
credential: the compression function's core, extracted and run as a
block cipher under the name SHACAL-2, survived the open European
NESSIE evaluation and was selected for its 2003 portfolio---the same
engine vetted from a second angle, as a cipher rather than a hash.)

For this paper's program the moral is double. The adversary's calculus
is the most refined body of exact knowledge about \emph{reduced}
members of the family---every trail is a theorem-with-certificate about
a specific truncation, machine-checkable pair in hand---and at the same
time it is purely \emph{local} knowledge: a trail certifies one orbit
of one difference, and says nothing global about the map. The
statistical coordinates of \cref{sec:randommaps} and the trail calculus
of this section are complementary instruments, and to our knowledge no
one has systematically pointed them at the same scaled-down family to
see whether the round where trails die is the round where the Flajolet
silhouette locks in (\cref{sec:projects}).

\begin{project}{Differential trails in miniature}{probability;
  programming; patience with bit manipulation}{4--6 weeks}
  Build the toll calculus for a toy \SHA{} on $8$-bit words. First
  verify \cref{fig:diffadd} empirically: tabulate the exact
  distribution of $\Delta(x \shaplus y)$ over all pairs for every
  one-bit $\Delta x$, and check the Lipmaa--Moriai
  cost formula~\citep{LipmaaMoriai2001}. Then implement a
  branch-and-bound search for the lowest-toll trail through $r$ rounds
  of the toy compression function and plot $\log_2(\text{best trail
  probability})$ against $r$. Two landmarks make the plot speak: the
  birthday line ($-n/2$ for an $n$-bit toy state), below which a trail
  is useless to an attacker, and the round at which your search itself
  starts to drown. You will have reproduced, at fruit-fly scale, both
  the designer's round-count decision and the attacker's wall---the two
  sides of \cref{sec:siege}---and the same code instruments the
  comparative questions of \cref{sec:projects}.
\end{project}

\subsection{Coda: can a machine learn the difference?}\label{sec:mlcrypt}

One newer thread deserves flagging, because it closes a loop with the
proof-barrier of \cref{sec:barriers} in a way that is either profound or
a curiosity, and nobody yet knows which. In 2019 Gohr trained a neural
network to distinguish the outputs of a round-reduced ARX cipher
(Speck, a near-relative of \SHA{}'s primitives) from random, and it
\emph{beat} the best hand-built differential distinguishers---and, on
inspection, had learned features \emph{beyond} the difference
distribution table, correlations the human theory of \cref{sec:diffcalc}
does not track~\citep{Gohr2019}. Read against Razborov--Rudich
(\cref{sec:barriers}), this is quietly startling. A trained
distinguisher \emph{is} a natural property in their precise sense---an
efficiently computable, large predicate that separates the function
from random---so a machine that reliably learns such properties on
\SHA{}-family maps would be walking, empirically, along the exact edge
the natural-proofs barrier says cannot be crossed at full scale.
Whether the network finds real, nameable structure (extractable back
into human mathematics, in the spirit of \cref{sec:hope}) or merely a
more efficient encoding of the same statistical residue is, at present,
an open and very testable question---and a natural apex to the
undergraduate program, where a student can point a modern learning tool
at a scaled-down family and see, concretely, what a machine can and
cannot pull out of manufactured randomness.

\begin{project}{A neural distinguisher at toy scale}{programming;
  a machine-learning course}{6--8 weeks (flagship)}
  Reproduce Gohr's experiment~\citep{Gohr2019} on a \SHA{}-family toy:
  train a small network to distinguish $r$-round outputs (or output
  differences) from random, and chart its advantage against $r$
  alongside the trail-toll curve of the miniature-trails project above
  and the statistical-battery frontier of \cref{sec:gap}. Then
  ask the hard question: \emph{interpret} the trained model---does its
  decision reduce to a differential-distribution feature the human
  theory already knows (\cref{sec:diffcalc}), or to something new? A
  negative result is a clean measurement; a positive one is the first
  thread of a natural property in a real hash family, and worth far
  more than a course grade.
\end{project}

% ------------------------------------------------------------
% §7  Computability, complexity, and barriers
% ------------------------------------------------------------
\section{The view from computability}\label{sec:computability}

The gap of \cref{sec:gap} is a statement about missing \emph{positive}
knowledge. It is worth closing the overview by asking the complementary,
sharper question: what is \emph{provably} true about the difficulty of
breaking \SHA{}? Here the news is a genuine mix of hard theorems and hard
barriers, and the boundary between them is itself part of the object's
significance.

\subsection{Breaking \SHA{} is an \textsc{np} search problem}
\label{sec:npsearch}

Fix a target digest $y$ and ask for a preimage: some $x$ with
$\SHA(x) = y$. Because \SHA{} is computable by a small circuit
(\cref{sec:history}), a candidate $x$ can be \emph{checked} in linear
time, so preimage-finding is an \textsc{np} search problem in the exact
sense of Cook and Levin~\citep{CookLevin1971}: the ``inversion''
relation is polynomial-time decidable, and the only question is the cost
of \emph{search}. Encode one evaluation of the compression function as a
Boolean circuit---about $2^{16}$ gates---and the equation
$\SHA(x)=y$ becomes a concrete instance of the satisfiability problem
\textsc{sat}, the canonical \textsc{np}-complete language.

The translation is worth seeing once, because it is almost embarrassingly
direct (\cref{fig:tseitin}). Give every wire of the circuit its own
Boolean variable; each gate then constrains its three wires by a handful
of clauses, satisfied exactly by the gate's truth table; pin the $256$
output variables to the bits of $y$. The conjunction of all clauses is
satisfiable \emph{if and only if} a preimage exists, and any satisfying
assignment, read off the input wires, \emph{is} a preimage. The whole of
\SHA{}'s round structure---$\Sigma$'s, carries, message schedule---flattens
into a few hundred thousand such clauses, syntactically indistinguishable
from any other \textsc{sat} instance. A modern conflict-driven
(``\textsc{cdcl}'') solver then does something almost mathematical with
it: it explores partial assignments, and each dead end is analyzed into
a \emph{learned clause}---a small lemma, added to the formula, recording
why that region of the search space is empty. A solver run is a
proof-by-cases of astronomical width, conducted by a machine that
invents its own case analysis.

\begin{figure}[ht]
  \centering
  \begin{tikzpicture}[font=\small,
      wire/.style={thick},
      lbl/.style={font=\scriptsize}]
    % --- little AND gate ---
    \draw[wire] (-1.0,0.55) node[left, lbl] {$u$} -- (0,0.55);
    \draw[wire] (-1.0,-0.15) node[left, lbl] {$v$} -- (0,-0.15);
    % gate body (D shape)
    \draw[thick, fill=projfill]
      (0,-0.5) -- (0.55,-0.5) arc[start angle=-90, end angle=90,
      radius=0.7] -- (0,0.9) -- cycle;
    \node at (0.42,0.2) {$\wedge$};
    \draw[wire] (1.25,0.2) -- (2.1,0.2) node[right, lbl] {$w$};
    % --- arrow to clauses ---
    \draw[->, thick] (2.9,0.2) -- (3.6,0.2);
    \node[anchor=west, align=left, font=\small] at (3.75,0.2)
      {$(\lnot u \vee \lnot v \vee w)$\\
       $(u \vee \lnot w)$\\
       $(v \vee \lnot w)$};
    % --- caption line under both ---
    \node[lbl, align=center, anchor=north] at (3.1,-1.0)
      {one fresh variable per wire, a few clauses per gate;\\
       ${\sim}2^{16}$ gates $\Rightarrow$ a few $\times 10^{5}$ clauses;
       pin $256$ output literals to $y$ and solve.};
  \end{tikzpicture}
  \caption{The reduction that makes ``invert \SHA{}'' a \textsc{sat}
  instance: each gate becomes clauses that are satisfied exactly by its
  truth table (shown: $w = u \wedge v$). Nothing about the construction
  is clever, and that is the point---membership in \textsc{np} is a
  statement about \emph{checkability}, and it holds for \SHA{} as
  mechanically as for any puzzle whose solutions can be verified.}
  \label{fig:tseitin}
\end{figure}

This is not a metaphor: the reduction is run in practice, and by now the
solvers are genuine instruments of record. As early as 2006, Mironov and
Zhang reproduced Wang-style MD5 collisions with an off-the-shelf
\textsc{sat} solver, no differential expertise wired in---the solver
rediscovering, clause by learned clause, constraints it had never been
told~\citep{MironovZhang2006}. Today the traffic runs both directions:
the current record semi-free-start collision on $39$-round \SHA{} was
found with a \textsc{sat}/\textsc{smt} search for the differential
trails of \cref{sec:craft}~\citep{LiLiuWang2024}, and a programmatic
\textsc{sat} solver---one that injects cryptanalytic knowledge as it
runs---produced the first practical $38$-round semi-free-start collision
independently of the classical toolchain~\citep{AlamgirNejatiBright2024}.
The empirical shape of all this effort is a cliff: instances for few
rounds collapse in milliseconds, then somewhere past round $20$ the
running times turn vertical, and no amount of solver engineering has
moved the wall more than a round or two per decade. The cliff's very
existence is one more coordinate in the natural history of the
function---and, unlike most such coordinates, it is cheap to measure.

\begin{project}{The \textsc{sat} cliff}{programming; propositional
  logic}{4--6 weeks}
  Build (or generate with existing tools) \textsc{cnf} encodings of
  $r$-round \SHA{} preimage instances (\cref{fig:tseitin}) and feed
  them to two or three modern solvers. Chart solving time against $r$
  until timeout; then vary what is pinned (full digest, partial digest,
  collision instead of preimage) and watch the cliff move. Plot, on the
  same axis, where the statistical batteries start failing
  (\cref{sec:gap}'s project) and where the differential trails of
  \cref{sec:siege} die: three instruments' frontiers, never to our
  knowledge drawn on one chart. The deliverable is that chart---a
  measured portrait of ``hardness turning on'' round by round.
\end{project}

But note precisely what \textsc{np}-completeness does and does not say.
It says the \emph{worst case} of \textsc{sat} is hard unless
$\mathrm{P}=\mathrm{NP}$; it says almost nothing about the specific,
highly structured instances that hashing produces. Cryptography needs
\emph{average-case} and \emph{one-way} hardness---hard instances that are
easy to \emph{generate} with a solution---and this is a strictly stronger
and more delicate demand than worst-case \textsc{np}-hardness. No
reduction shows that inverting \SHA{} is \textsc{np}-hard, and this is
not merely an unfilled cell in the table: there is a theorem-shaped
reason to expect no such reduction. Akavia, Goldreich, Goldwasser, and
Moshkovitz showed that a black-box reduction (of the standard,
non-adaptive kind) from an \textsc{np}-hard problem to \emph{inverting a
one-way function} would itself trigger complexity-theoretic collapses
nobody believes~\citep{AkaviaGGM2006}. The missing bridge between
$\mathrm{P} \ne \mathrm{NP}$ and \SHA{}'s security is not awaiting a
clever graduate student; by the best current evidence, the standard
bridge-building technology provably cannot span that river.

\subsection{Five worlds, one function}\label{sec:worlds}

The distinction just drawn---worst-case hardness versus generatable
hardness---organizes the deepest open question in the area, and
Impagliazzo gave it an unforgettable form: five possible worlds, exactly
one of which is real~\citep{Impagliazzo1995}. \Cref{fig:worlds} draws
the ladder and scores our function in each world. The point of the
exercise for this paper: \emph{whether an object like \SHA{} can exist
at all} is not a detail of cryptography but the exact content of the
rung between Pessiland and Minicrypt---the existence of one-way
functions---and it is wide open. Everything this paper treats as
mysterious about one specific function is, one level up, the mystery on
which the classification of computational reality turns.

\begin{figure}[ht]
  \centering
  \begin{tikzpicture}[font=\small,
      wname/.style={font=\bfseries\small, anchor=west},
      wdesc/.style={font=\scriptsize, align=left, anchor=west,
                    text width=82mm, inner sep=0pt},
      lbl/.style={font=\scriptsize}]
    % strata, bottom = algorithmica; each row 1.2cm tall, 11.8cm wide
    \fill[padfill]  (0,4.8) rectangle (11.8,6.0);
    \fill[projfill] (0,3.6) rectangle (11.8,4.8);
    \fill[lenfill]  (0,2.4) rectangle (11.8,3.6);
    \fill[red!7]    (0,1.2) rectangle (11.8,2.4);
    \fill[red!14]   (0,0)   rectangle (11.8,1.2);
    \draw (0,0) rectangle (11.8,6.0);
    \foreach \y in {1.2,2.4,3.6,4.8} \draw (0,\y) -- (11.8,\y);
    \node[wname] at (0.2,5.4) {Cryptomania};
    \node[wdesc] at (2.85,5.4)
      {trapdoors exist: public-key cryptography, the internet as
       deployed (\SHA{} still does the hashing)};
    \node[wname] at (0.2,4.2) {Minicrypt};
    \node[wdesc] at (2.85,4.2)
      {one-way functions exist: \SHA{} can be what it appears---PRGs,
       signatures, proof-of-work all sound};
    \node[wname] at (0.2,3.0) {Pessiland};
    \node[wdesc] at (2.85,3.0)
      {hard random instances abound, but none can be generated
       \emph{with its solution}: no one-wayness; \SHA{} invertible};
    \node[wname] at (0.2,1.8) {Heuristica};
    \node[wdesc] at (2.85,1.8)
      {\textsc{np} hard only in the worst case; typical instances
       easy---\SHA{} broken in practice};
    \node[wname] at (0.2,0.6) {Algorithmica};
    \node[wdesc] at (2.85,0.6)
      {$\mathrm{P}=\mathrm{NP}$ (in spirit): every preimage findable
       fast; cryptography impossible, algorithmics paradise};
    % the one-way-function line
    \draw[very thick, dashed, citewine] (0,3.6) -- (12.35,3.6);
    \node[lbl, citewine, anchor=north west, align=left, text width=21mm]
      at (12.0,3.45) {the one-way-\\function line};
    % brace: SHA's conjectured home
    \draw[decorate, decoration={brace, amplitude=6pt}]
      (12.0,3.7) -- (12.0,5.95)
      node[midway, right=8pt, lbl, align=left, text width=21mm]
      {if \SHA{} is what it appears to be, the real world is here};
  \end{tikzpicture}
  \caption{Impagliazzo's five worlds: mutually
  exclusive possibilities for how computation works, exactly one of
  them actual, none of them ruled out. \SHA{}'s conjectured properties
  assert that we live above the one-way-function line (at least
  Minicrypt)---a claim strictly \emph{weaker} than what public-key
  cryptography bets on, and strictly \emph{stronger} than
  $\mathrm{P} \ne \mathrm{NP}$. Note where the quantum threat falls:
  Shor-type algorithms attack the classical inhabitants of Cryptomania,
  while against the generic hashing of Minicrypt the provable quantum
  gain is only the square root of \cref{sec:barriers}.}
  \label{fig:worlds}
\end{figure}

Two theorems furnish that observation with content, and a third---the
one this section builds toward---turns it into an exact equivalence.

First: \emph{Minicrypt has a canonical witness}. Levin constructed an
explicit, fully specified ``universal one-way function''---a concrete
program one can write down today---which is one-way \emph{if any one-way
function exists at all}~\citep[the tale, and the construction, are told
in][]{Levin2003}. Existence of the whole kingdom is thus equivalent to
the hardness of one specific function, in print since the 1980s. Why,
then, does anyone bother with \SHA{}? Because Levin's function is a
diagonalization---roughly, it runs all short programs and combines
their outputs---and its guarantee is bought at constants and exponents
that make it a philosophical instrument, not an engineering material.
Once any one-way function exists, moreover, the whole private-key
toolbox follows by explicit constructions: H{\aa}stad, Impagliazzo,
Levin, and Luby showed pseudorandom generators can be built from
\emph{any} one-way function~\citep{HILL1999}, with signatures and
message authentication downstream. Read against \cref{sec:gap}: theory
already possesses, with proofs, the entire architecture the world
needs---resting on an unproven ground floor and running too slowly to
use. \SHA{} is the engineering shadow of Levin's function: the object
you deploy when you want the universal guarantee at gigabytes per
second, and are willing to trade the theorem away for it.

Second, the same result read as a boundary: the Akavia--Goldreich--%
Goldwasser--Moshkovitz theorem above says the ground floor cannot be
poured from worst-case \textsc{np}-hardness by standard reductions. The
two faces together frame the epistemic situation exactly: everything
above the one-way-function line is solved mathematics; the line itself
is untethered to the rest of complexity theory.

\subsection{An exact address for the one-way line}\label{sec:liupass}

The third result is recent (2020) and, to our taste, the most
beautiful in the whole subject, because it identifies the one-way-%
function line of \cref{fig:worlds} with the hardness of a problem this
paper has been circling from the start: \emph{detecting structure in
apparently random strings}. Making that precise takes four definitions,
each elementary; we give them carefully, because the payoff is a genuine
``if and only if'' and it is worth being able to read it exactly.

\paragraph{Kolmogorov complexity.} Fix once and for all a universal
Turing machine $U$ (a fixed interpreter: it runs a program $\Pi$ on an
input and prints the result). The \emph{Kolmogorov complexity} of a
string $x \in \bitstrings^{*}$ is the length of the shortest program
that outputs it,
\[
  K(x) \;=\; \min\bigl\{\, |\Pi| \;:\; U(\Pi) = x \,\bigr\}.
\]
This is the honest measure of ``information content'': a patterned
string like $(01)^{n}$ has a tiny generator (``print \texttt{01}
$n$ times''), so $K \approx \log_2 n$, while a structureless string has
no description shorter than itself, $K(x) \approx |x|$. The trouble, for
our purposes, is that $K$ is \emph{uncomputable}---pinning it down
solves the halting problem---so it cannot be the hard problem
cryptography stands on. Something must give the program a deadline.

\paragraph{Time-bounded Kolmogorov complexity.} Fix a polynomial time
bound $t(\cdot)$. The \emph{$t$-time-bounded} complexity restricts the
minimization to programs that print $x$ \emph{within $t(|x|)$ steps}:
\[
  K^{t}(x) \;=\; \min\bigl\{\, |\Pi| \;:\;
      U(\Pi) = x \text{ in at most } t(|x|) \text{ steps} \,\bigr\}.
\]
Now the quantity \emph{is} computable---enumerate all programs of length
$\le |x|$, run each for $t(|x|)$ steps, keep the shortest that prints
$x$---but that brute force costs about $2^{|x|}$. The entire question is
whether one can do better than brute force, and specifically whether one
can do better \emph{on average}, over a random string $x$.

\paragraph{Why almost every string is incompressible.} A counting fact
fixes the landscape (\cref{fig:kolmogorov}). There are only
$2^{n-d}-1 < 2^{n-d}$ programs shorter than $n-d$ bits, so at most a
$2^{-d}$ fraction of the $2^{n}$ strings of length $n$ can have
$K^{t}(x) < n-d$. Compressible strings---those with real structure a
fast program can exploit---are \emph{exponentially rare}; the mass of
$K^{t}$ piles up just below the diagonal $K^{t}(x)=n$. In particular the
\emph{trivial estimate} ``$K^{t}(x) = |x|$''---equivalently, the
one-line algorithm that ignores $x$ and outputs its length---is correct
to within $d$ bits on all but a $2^{-d}$ fraction of inputs. Estimating
$K^{t}$ better than this trivial baseline is precisely the act of
\emph{detecting the rare structure}, and that is the task whose
difficulty the theorem measures.

\begin{figure}[ht]
  \centering
  \begin{tikzpicture}[font=\small]
    % axes
    \draw[->] (0,0) -- (7.4,0) node[below left=0.5mm and -8mm]
      {$K^{t}(x)$};
    \draw[->] (0,0) -- (0,3.5) node[above, align=center, font=\scriptsize]
      {\# of $n$-bit\\ strings};
    \draw[dashed, gray] (6,0) -- (6,3.2) node[above, font=\scriptsize,
      black] {$n$};
    % exponential pile-up near n: bars doubling toward n
    \foreach \k/\h in {1.2/0.10, 1.8/0.14, 2.4/0.2, 3.0/0.28, 3.6/0.42,
                       4.2/0.62, 4.8/0.95, 5.4/1.5, 5.85/2.7} {
      \fill[projframe!75] (\k-0.11,0) rectangle (\k+0.11,\h);
    }
    \node[font=\scriptsize, align=center, projframe] at (2.5,2.35)
      {structured strings:\\ a $2^{-d}$ sliver\\ with $K^{t}<n-d$};
    \draw[->, projframe, font=\scriptsize] (2.5,1.85) -- (3.1,0.45);
    \node[font=\scriptsize, align=center, anchor=east] at (7.3,2.4)
      {almost all mass:\\ $K^{t}(x)\approx n$\\ (incompressible)};
    \draw[decorate, decoration={brace, amplitude=4pt}]
      (0.15,-0.35) -- (4.2,-0.35)
      node[midway, below=4pt, font=\scriptsize]
      {beating ``output $|x|$'' here $=$ detecting structure};
  \end{tikzpicture}
  \caption{The distribution of time-bounded Kolmogorov complexity over
  random $n$-bit strings. By counting, at most a $2^{-d}$ fraction can
  be compressed below $n-d$; the histogram is crushed against the
  incompressible edge $K^{t}\!=\!n$. The one-line ``output $|x|$''
  algorithm is therefore right for nearly everyone---and Liu--Pass say
  that noticeably out-performing it, on average, is exactly as hard as
  inverting every one-way function.}
  \label{fig:kolmogorov}
\end{figure}

\paragraph{Mild average-case hardness, and one-wayness.} Two more
definitions, now the load-bearing ones. Say $K^{t}$ is
\emph{mildly hard-on-average} if some polynomial $p$ makes it
uncomputable on a noticeable slice: for every probabilistic
polynomial-time algorithm $A$ and all large $n$,
\[
  \Pr_{x \,\gets\, \bitstrings^{n}}
     \bigl[\, A(x) \ne K^{t}(x) \,\bigr] \;\ge\; \frac{1}{p(n)} .
\]
The word to dwell on is \emph{mild}: $A$ is permitted to be right on a
$1 - 1/p(n)$ super-majority of strings; hardness on just an
inverse-polynomial sliver suffices. And a \emph{one-way function}---the
one-wayness informally invoked back in \cref{sec:npsearch}---is a
polynomial-time computable $f$ that no probabilistic polynomial-time $A$
can invert except with vanishing probability: for every such $A$, every
polynomial $p$, and all large $n$,
\[
  \Pr_{x \,\gets\, \bitstrings^{n}}
     \bigl[\, f(A(f(x))) = f(x) \,\bigr] \;<\; \frac{1}{p(n)} .
\]
Existence of any such $f$ is the Minicrypt hypothesis---the rung of
\cref{fig:worlds} on which \SHA{}'s conjectured life depends.

\begin{theorem}[Liu--Pass 2020]\label{thm:liupass}
For every polynomial $t(n) \ge (1+\varepsilon)n$ with $\varepsilon>0$
constant, the following are equivalent:
\begin{enumerate}
  \item One-way functions exist (equivalently: secure private-key
  encryption, pseudorandom generators, pseudorandom functions, digital
  signatures, and commitments all exist).
  \item $K^{t}$ is mildly hard-on-average.
\end{enumerate}
Moreover the equivalence is robust to approximation: if one-way
functions exist then no efficient algorithm can even
$(d\log n)$-approximate $K^{t}$ on meaningfully more than the trivial
fraction, for any constant $d$~\citep{LiuPass2020}.
\end{theorem}

The approximation clause is the sharpest way to feel the result. The
do-nothing algorithm ``output $|x|$'' already $(d\log n)$-approximates
$K^{t}$ with probability $\ge 1 - n^{-d}$ (the counting fact again); the
theorem says that beating even \emph{that} baseline, by any
inverse-polynomial margin, is as hard as breaking \emph{every} one-way
function at once. A hair of non-trivial structure-detection, made
reliable, would collapse all of private-key cryptography.

Now read \cref{thm:liupass} in this paper's language. ``$K^{t}(x)$ is
small'' means \emph{$x$ has detectable structure}---a short program
generates it fast. Estimating $K^{t}$ on random inputs is therefore
exactly \emph{hunting for hidden structure in apparently random
strings, on average}---the naturalist's question of \cref{sec:epistemology},
promoted from one function to all strings. The theorem says this hunt is
(mildly) hard \emph{if and only if} objects like \SHA{} can exist at
all. The connection is not a loose analogy; it is the same activity. A
pseudorandom string---say a stretch of the counter stream $\SHA(1),
\SHA(2), \dots$ of \cref{sec:below}---has genuinely \emph{low} $K^{t}$
(the short program ``\,print $\SHA(1..m)$\,'' generates it quickly), yet
is engineered to look maximally incompressible; a distinguisher that
noticed the gap would be estimating $K^{t}$ better than trivial on
exactly these strings, which is to say breaking the generator. So the
empirical program of \cref{sec:gap}---point an instrument at a
pseudorandom object and measure whether any structure pokes through---is
not merely \emph{analogous} to the problem on which the existence of
cryptography rests. Up to the formalization, it \emph{is} that problem,
shrunk to a single function and a laptop. We know no cleaner statement
of why the small experiments of \cref{sec:projects} sit on the main line
of the subject rather than beside it.

\subsection{Two-sided barriers: what is provably out of reach}
\label{sec:barriers}

The deepest results in this corner are \emph{impossibility} theorems,
and they cut in two directions at once.

On the side of \emph{attacks}, the quantum model yields rare
unconditional lower bounds. Grover search inverts any black-box function
on $N$ points in $\Theta(\sqrt{N})$ queries, and this is provably
optimal---no quantum algorithm does better \emph{relative to the
oracle}~\citep{BBBV1997}; for collisions, the Brassard--H{\o}yer--Tapp
birthday-style algorithm costs $\Theta(N^{1/3})$, and Aaronson and Shi
proved a matching $\Omega(N^{1/3})$ quantum lower
bound~\citep{AaronsonShi2004}. These theorems bound what \emph{any}
algorithm can do against a function that behaves as a random oracle, so
they translate into concrete, believed-tight security levels for
\SHA{}---$128$ quantum bits of preimage resistance, and the reason a
$256$-bit digest is considered ``quantum-safe'' for collision resistance.
The catch, familiar by now, is the model: the bounds hold in the
black-box world and become claims about \SHA{} only under the heuristic
that no shortcut exploits its actual circuit.

It is worth pausing to give exponents like $2^{128}$ an empirical
scale, because the planet happens to be running the relevant
experiment. The proof-of-work industry of \cref{sec:uses} evaluates
\SHA{} about $10^{21}$ times per second---roughly $2^{70}$---and its
cumulative total to date is on the order of $2^{96}$ evaluations:
comfortably the largest computation of any kind in human history,
every cycle of it hammering on this one function. Two readings. As a
\emph{statistical} instrument: block arrivals track the uniform model
with textbook fidelity, so this is the most expensive test of the
random-function conjecture ever conducted, passed to date. As a
\emph{yardstick}: the deepest partial inversion all that work has ever
exhibited is a digest beginning with roughly a hundred zero
bits---each mined block is a certified partial preimage, currently
costing about $2^{79}$ evaluations apiece---which still leaves the full
$256$-bit inversion about $2^{156}$ times beyond the species' total
output, and even the $2^{128}$ birthday bound a factor of four billion
past everything computed so far. When \cref{sec:telescope} says the
global geometry of the functional graph is ``operationally
unreachable,'' this is the industrial base against which
``unreachable'' is calibrated.

On the side of \emph{proofs}, the barrier theorems explain the gap of
\cref{sec:gap} as something deeper than a lack of effort. The
\emph{natural proofs} barrier of Razborov and
Rudich~\citep{RazborovRudich1997} shows that any proof of a strong
circuit lower bound which is itself ``constructive and large''---the form
essentially all known lower-bound techniques take---would, ironically,
yield an algorithm breaking the very pseudorandom functions such a lower
bound is meant to certify. In slogan form: \emph{the existence of objects
like \SHA{} is precisely what prevents us from proving that objects like
\SHA{} are hard} by current methods. This is why nobody expects an
unconditional proof of \SHA{}'s security in the foreseeable future, and
why the honest theorems in the subject are either conditional
(\cref{sec:frameworks}), relativized (the quantum bounds above), or about
components rather than the whole (\cref{thm:rivest}).

\subsection{What the proof assistants did prove}\label{sec:verified}

After a section of impossibilities, it is worth recording where
machine-checked mathematics has \emph{succeeded} with \SHA{}---because
it has, completely, one floor down. In 2015 Appel produced a formal,
machine-checked proof, in the Coq proof assistant, that OpenSSL's C
implementation of \SHA{} computes exactly the function specified in
FIPS 180-4: every pointer offset, every rotation, every corner of C
semantics, verified down to a mechanized model of the
compiler's view of the program~\citep{Appel2015}. The HACL\textsuperscript{*}
project carried the method to an entire cryptographic library, written
in the F\textsuperscript{*} proof language and compiled to C; its
verified \SHA{} ships inside Firefox, executing in ordinary users'
browsers~\citep{Zinzindohoue2017}. These are not toy exercises; they
retire a whole species of real historical failure, the
\emph{implementation} bug---and they are proofs about \SHA{} in the
strictest sense of the word.

The epistemological inversion deserves to be savored. \emph{Everything
provable about \SHA{} has now been proved by machines: namely, that the
program is the function.} What no one can prove is any property of the
function itself. Formal verification thus sharpens the boundary of
\cref{sec:gap} into a visible line: below it, mechanized
certainty---code equals specification, checked symbol by symbol; above
it, Rogaway's human ignorance (\cref{sec:definitions}), trust
underwritten by two decades of failed attack. The axiom has quietly
moved. It is no longer ``this function is one-way''; it is ``this
function is the one we wrote down''---and only the second statement is
checkable. A mathematician could not ask for a cleaner exhibit of where,
exactly, in one industrial object, the proved ends and the conjectured
begins.

\subsection{Where this leaves the program}
\label{sec:leaves}

The computability view sharpens, rather than dampens, the case for the
empirical program of this paper. Unconditional hardness is provably out
of reach by known methods; worst-case \textsc{np}-hardness is the wrong
target, and provably unavailable by standard reductions
(\cref{sec:worlds}); the clean theorems that do exist are relativized
bounds assuming random-oracle behavior, or certify the code rather than
the function (\cref{sec:verified}). What is \emph{not} ruled out---by
any barrier theorem---is the discovery of exact, non-statistical
structure in specific reduced systems, of the kind \cref{sec:hope}
catalogued next door and \cref{thm:rivest} exhibited inside the function
itself. The barriers forbid a certain style of \emph{general} proof;
they say nothing against a Josephus-style collapse of a
\emph{particular} finite object. That is exactly the crack the
undergraduate program of \cref{sec:projects} proposes to work in: not to
prove \SHA{} hard---provably hopeless---but to map, instance by instance
and coordinate by coordinate, where the structure ends and the noise
begins.

% ------------------------------------------------------------
% §6  The research program, cut into undergraduate-sized pieces
% ------------------------------------------------------------
\section{Undergraduate short projects}\label{sec:projects}

The preceding sections flagged thirteen concrete projects at the point where
each becomes natural; here we collect them and add the program-level view.
The unifying activity is \emph{measurement}: locating scaled-down members
of the \SHA{} design family in coordinate systems (Flajolet's above all)
where ``looks random'' becomes a number with error bars. None of the
projects requires cryptographic background beyond this paper; all require
programming and a taste for experiment; each produces figures and data
that are a contribution regardless of what they show. That last property
is worth underlining for research with undergraduates in short windows:
these are experiments that \emph{cannot fail}, only surprise.

\begin{center}
\small
\begin{tabular}{@{}p{0.30\linewidth}p{0.22\linewidth}p{0.12\linewidth}p{0.24\linewidth}@{}}
  \toprule
  project (defined in) & prerequisites & window & deliverable \\
  \midrule
  Single ops under the $2$-adic lens (\cref{sec:statespace}) &
    elementary number theory & 3--4 wks &
    verified ergodicity criteria, orbit visualizations \\[2pt]
  The linear algebra of diffusion (\cref{sec:rivest}) &
    linear algebra & 3--4 wks &
    invertibility \& group structure of the $\Sigma/\sigma$ maps \\[2pt]
  Exact inverses of $\Sigma_0, \Sigma_1$ (\cref{sec:rivest}) &
    polynomials over $\Ftwo$ & 2--3 wks &
    closed-form inverse maps, verified two ways \\[2pt]
  Functional-graph atlas of truncated \SHA{} (\cref{sec:telescope}) &
    probability, programming & 6--8 wks &
    atlas of measured vs.\ predicted statistics \\[2pt]
  Collision hunting, van Oorschot--Wiener style (\cref{sec:telescope}) &
    probability, programming & 4--6 wks &
    timing distributions vs.\ theory; parallel implementation \\[2pt]
  The dynamics of iteration (\cref{sec:iterated}) &
    probability, programming & 6--8 wks &
    image-shrinkage law, fixed points, Hellman trade-off \\[2pt]
  Cutoff for toy compression walks (\cref{sec:prngs}) &
    probability, linear algebra & 4--6 wks &
    spectral gaps and cutoff profiles vs.\ rounds \\[2pt]
  Cayley hashes by hand (\cref{sec:provable}) &
    group theory, linear algebra & 3--4 wks &
    spectral gap vs.\ equidistribution; collision gallery \\[2pt]
  Differential trails in miniature (\cref{sec:siege}) &
    probability, programming & 4--6 wks &
    exact toll tables; best-trail probability vs.\ rounds \\[2pt]
  The \textsc{sat} cliff (\cref{sec:npsearch}) &
    programming, logic & 4--6 wks &
    solver-time cliff chart; three frontiers on one axis \\[2pt]
  A neural distinguisher at toy scale (\cref{sec:mlcrypt}) &
    programming, machine learning & 6--8 wks &
    learned-advantage vs.\ rounds; interpretation of the model \\[2pt]
  Where do the tests start failing? (\cref{sec:gap}) &
    statistics, programming & 4--6 wks &
    round-by-round battery failure chart \\[2pt]
  Avalanche cartography (\cref{sec:gap}) &
    linear algebra, programming & 4--6 wks &
    influence-matrix heatmaps by round \\
  \bottomrule
\end{tabular}
\end{center}

\medskip

Three remarks on the program as a whole.

\paragraph{Model organisms.} The full-scale function is beyond experiment
(\cref{sec:telescope}) and beyond proof (\cref{sec:gap}); the program
therefore studies the \emph{family} obtained by shrinking each design
parameter---word size $32 \to 8$ or $16$, state words $8 \to 4$, rounds
$64 \to r$---the way biology studies fruit flies. The two-geometry anatomy
of \cref{sec:statespace} says which shrinkings are honest: one must
preserve the ARX alternation, or the object degenerates into something a
single lens solves. A systematic scaled family, with its statistics
measured, would be a reusable community resource; no standard one exists.

\paragraph{The comparative method.} Every measurement gains meaning from
controls on both sides: genuinely random maps (the Flajolet silhouette)
on one side, and visibly structured maps (bijections, $\Ftwo$-linear
maps, one-round toys) on the other. The interesting science is in the
\emph{trajectory}---how, as rounds accumulate, a \SHA{}-family map walks
from the structured region onto the random silhouette, in which
coordinates it arrives first, and whether the walk is gradual or a phase
transition. SHA-1 and \SHA{} can be given the identical treatment; given
\cref{sec:falling}, any measurable divergence between their trajectories
would be a finding of genuine interest.

\paragraph{Odds of surprise.} Realistically: most measurements will
confirm the random silhouette, and their value is the atlas itself plus
trained students. But the history of \cref{sec:josephus} cuts both
ways---structure, when it exists in innocent-looking discrete processes,
has repeatedly been found by exactly this kind of patient small-case
empiricism, and essentially nobody has pointed the analytic-combinatorics
instrument at this family at model-organism scale. The expected value of
looking is small-positive; the tail is long.

\medskip

\noindent\emph{On tooling.} Everything above runs comfortably in Python
or Julia on a laptop for $n \le 30$; the collision-hunting and atlas
projects scale gracefully to whatever cluster time is available. We
maintain reference implementations and are happy to share scaffolding so
that student time goes to science rather than plumbing.


% ============================================================
% Shared backmatter: appendix + thematically clustered
% references. Each cluster is regrouped inside the single .bbl
% by tools/clusterbbl.py; only cited works print. Green
% \localpdf marks link to archived copies in refs/.
% ============================================================
\appendix
% ------------------------------------------------------------
% §A  Appendix: reference implementation for the Sigma inverses
% ------------------------------------------------------------
\section{Reference implementation: inverting $\Sigma_0$ and $\Sigma_1$}
\label{app:sigmacode}

The listing below (\texttt{tools/invert\_sigma.py} in the distribution
accompanying this document) is the computational companion to
\cref{sec:rivest}. It is included in full for three reasons.

First, it discharges a proof obligation. \Cref{thm:rivest} certifies
that $\Sigma_0$ and $\Sigma_1$ are bijections but exhibits nothing;
\cref{eq:sigmainv} exhibits the inverses but asks to be believed. The
script closes the loop: it constructs each map's $32\times32$ matrix
over $\Ftwo$, inverts it by Gaussian elimination, reads the inverse off
as an xor of rotations (the inverse of a circulant map is circulant, so
the image of the basis word $1$ determines everything), and then
verifies---functionally, against the standard's own
definition~\citep{FIPS180-4} of $\mathrm{ROTR}$---that the seventeen-term
maps of \cref{eq:sigmainv} invert $\Sigma_0$ and $\Sigma_1$ on thousands
of words, along with the order facts $\Sigma_0^{8} = \mathrm{ROTR}^{8}$,
$\Sigma_0^{32} = \Sigma_1^{16} = \mathrm{id}$, and
$\Sigma_i^{-1} = \Sigma_i^{\mathrm{ord}-1}$ from \cref{sec:rivest}.
A skeptical reader can rerun the whole certificate in under a second.

Second, the verification is \emph{convention-proof}, and seeing why is a
small lesson in the algebra of \cref{sec:rivest}. Any claim phrased in
the polynomial ring $R$ silently commits to a bit-numbering convention
(is bit $i$ the coefficient of $x^i$ or of $x^{31-i}$? is rotation
multiplication by $x^r$ or $x^{-r}$?), and such commitments are exactly
where off-by-one and mirror-image errors breed. But the composition law
$\mathrm{ROTR}^{a} \circ \mathrm{ROTR}^{b} = \mathrm{ROTR}^{a+b \bmod 32}$
holds under every convention, so the statement ``the rotation multiset
$S_0 + \{2,13,22\} \bmod 32$ covers $0$ an odd number of times and every
other offset an even number of times'' -- which is what
\cref{eq:sigmainv} asserts -- is convention-free, and the script checks
it against the operational definition rather than any ring
identification. This is the paper's epistemological theme
(\cref{sec:epistemology}) in miniature: replace an argument one must
trust with an artifact anyone can audit.

Third, it is deliberately written as scaffolding. The two project boxes
of \cref{sec:rivest} ask for the same computation done other
ways---extended Euclid in $\Ftwo[x]$, repeated Frobenius squaring---and
for the $\sigma_0, \sigma_1$ question that \cref{eq:sigmainv}
deliberately leaves open; the functions here (\texttt{matrix\_of},
\texttt{invert\_matrix}) apply verbatim to the shift-containing maps,
and changing the single constant \texttt{W} repeats every experiment at
any word length, where Rivest's gcd condition~\citep{Rivest2011} comes
alive (try \texttt{W = 33}).

\medskip
\lstinputlisting[language=Python]{tools/invert_sigma.py}


\clearpage
\phantomsection\addcontentsline{toc}{section}{References}
\section*{References}

\noindent{\small The references are grouped by theme rather than
alphabetically \emph{en masse}, to make the subject's connections
visible. Green marks ({\color{archgreen}$\blacktriangleright$}) link to
archived copies in the \texttt{refs/} folder distributed with this
document (see \texttt{refs/README.md} for provenance); works without a
mark carry a note saying where to find them.\par}

% \refclusteritem is emitted by tools/clusterbbl.py between entry
% groups inside the .bbl's single thebibliography environment.
\newcommand{\refclusteritem}[2]{%
  \item[]\leavevmode
  \hspace*{-\leftmargin}%
  \begin{minipage}{\linewidth}
    \vspace{0.9em}
    {\normalsize\bfseries #1}\\[2pt]
    {\small\itshape #2}
    \vspace{0.1em}
  \end{minipage}}

\bibliographystyle{plainnat}
\bibliography{bib/s1-solvable,bib/s2-prehistory,bib/s3-family,bib/s4-uses,bib/s5-mappings,bib/s6-prng,bib/s7-foundations,bib/s10-provable,bib/s9-differential,bib/s8-computability}


\end{document}
